\chapter{Compressible flow, aerodynamics}
\label{vol03:ch13}

\epigraph{The phenomena of high-speed flow stand to those of
low-speed flow as those of the wave of finite amplitude stand to
those of the wave of infinitesimal amplitude: that is, the linear
theory which served so well loses its monopoly, and we must learn
to think in nonlinear terms.}{John D.\ Anderson, Jr.,
\emph{Fundamentals of Aerodynamics}, paraphrasing the opening of
the compressible-flow chapters \cite[ch.~7,
introductory remarks]{text:anderson-aerodynamics}.}

\chapmeta{Half-life: durable. Archetypes: transport, stability, scaling. Exercise target: 38. Project track: paper / analytical (wave-drag estimate of a passenger jet at cruise). Hazard class: none. Reader path default: standard, with sections explicitly tagged. Named cases: Tacoma Narrows aeroelastic flutter (1940, cross-referenced); transonic transport airfoil envelope; Columbia (2003, cross-referenced).}

\noindent
The previous two chapters treated the flow as incompressible, with
density independent of pressure along any streamline. Above about
Mach $0.3$, density variations become large enough that the
incompressible analysis ceases to give quantitatively correct
answers; above Mach $1$, the flow develops qualitatively new
features that the incompressible equations cannot represent at
all: discontinuous shock waves, choking, supersonic expansions,
and the transition from elliptic to hyperbolic governing equations.
This chapter develops compressible-flow analysis from the speed of
sound through the isentropic relations to normal and oblique
shocks, Prandtl-Meyer expansions, and the aerodynamic regimes that
divide subsonic, transonic, supersonic, and hypersonic flight. The
reader leaves able to compute stagnation properties from
freestream conditions at any Mach number, predict the pressure and
velocity jumps across a normal shock, estimate the wave drag on a
supersonic body, and recognise the design boundaries that separate
the four flow regimes.

\section{Speed of sound}\pathtag{core}

A small pressure disturbance in a compressible fluid propagates at
a characteristic speed, the \engterm{speed of sound} $a$. Consider
a thin region across which the pressure and density are perturbed
by $dp$ and $d\rho$, and the flow speed by $du$, with the
disturbance propagating into still fluid at speed $a$. The
disturbance itself is unsteady, but it is rendered steady by riding
with it: in a reference frame travelling with the wave front the
still fluid ahead approaches at speed $a$ and the fluid behind
recedes at $a - du$, and the steady conservation laws apply. For a
control volume that straddles the front, conservation of mass and
momentum give
\[
\rho a = (\rho + d\rho)(a - du), \quad p + \rho a^2 = (p + dp)
+ (\rho + d\rho)(a - du)^2.
\]
The first equation expands to $\rho a = \rho a - \rho\,du
+ a\,d\rho - d\rho\,du$, and dropping the second-order product
$d\rho\,du$ leaves the mass relation in linearised form,
\[
\rho\,du = a\,d\rho.
\]
The momentum equation expands to $p + \rho a^2 = p + dp + \rho a^2
- 2\rho a\,du + a^2 d\rho + O(\text{second order})$; cancelling
$p + \rho a^2$ on both sides and keeping first-order terms gives
\[
dp = 2\rho a\,du - a^2 d\rho.
\]
Substituting the mass relation $\rho\,du = a\,d\rho$ into the
momentum result, $2\rho a\,du = 2a^2 d\rho$, so $dp = 2a^2 d\rho
- a^2 d\rho = a^2 d\rho$, and
\[
a^2 = \frac{dp}{d\rho}\bigg|_{s},
\]
where the partial derivative is taken at constant entropy. The
isentropic qualifier is earned, not assumed: a small-amplitude
disturbance compresses and expands the fluid so gently and so
quickly that the temperature gradients it sets up have no time to
conduct heat across the wave, so the process is adiabatic, and
being adiabatic and (to leading order) reversible, it is isentropic.
For an ideal gas undergoing an isentropic process, $p\rho^{-\gamma}
= \text{constant}$, with $\gamma = c_p/c_v$ the \engterm{ratio of
specific heats}. Differentiating $p = C\rho^\gamma$ at constant
entropy gives $dp/d\rho = C\gamma\rho^{\gamma - 1} = \gamma
(C\rho^\gamma)/\rho = \gamma p/\rho$, and the ideal-gas law
$p = \rho R T$ turns $p/\rho$ into $RT$. Then
\[
a^2 = \gamma\frac{p}{\rho} = \gamma R T,
\]
where $R$ is the specific gas constant and $T$ the absolute
temperature. For air ($\gamma = 1.4$, $R = 287\,\si{\joule\per
\kilogram\per\kelvin}$, $T = 288\,\si{\kelvin}$ at sea level),
$a \approx 340\,\si{\meter\per\second}$. The same algebra with the
isothermal derivative $dp/d\rho = p/\rho = RT$ in place of the
isentropic one gives Newton's original sound speed $a =
\sqrt{RT}$, low by the factor $\sqrt{\gamma}$; Laplace's
recognition that the process is adiabatic rather than isothermal
restored the missing $\sqrt{\gamma}\approx 1.18$ and closed the
long-standing gap between Newton's prediction and the measured
speed.

The speed of sound is the propagation speed of infinitesimal
pressure disturbances. A finite-amplitude pressure disturbance
propagates at a different speed (the topic of section 13.4); the
infinitesimal-amplitude speed sets the lower bound and the
benchmark.

The speed of sound depends only on temperature for an ideal gas:
it is independent of pressure at fixed temperature. At
altitude, the air temperature drops, and the speed of sound drops
correspondingly:

\begin{itemize}[leftmargin=*]
  \item Sea level ($T = 288\,\si{\kelvin}$): $a \approx 340\,\si{
    \meter\per\second}$.
  \item $11\,\si{\kilo\meter}$ (typical cruise, $T \approx 217\,\si{
    \kelvin}$): $a \approx 295\,\si{\meter\per\second}$.
  \item $20\,\si{\kilo\meter}$ (lower stratosphere, $T \approx 217\,
    \si{\kelvin}$): $a \approx 295\,\si{\meter\per\second}$
    (isothermal stratosphere).
\end{itemize}

A speed of $250\,\si{\meter\per\second}$ at sea level corresponds
to Mach $0.74$; the same speed at $11\,\si{\kilo\meter}$
corresponds to Mach $0.85$. The Mach number is the engineering
parameter, not the airspeed, when compressibility effects are at
issue.

\section{Mach number, flow regimes}\pathtag{core}

The \engterm{Mach number} $M = u/a$ is the ratio of flow speed to
local speed of sound. The Mach number divides the compressible-flow
regimes:

\begin{itemize}[leftmargin=*]
  \item \engterm{Incompressible} ($M < 0.3$): density variations
    are small, $\rho/\rho_0$ differs from $1$ by less than
    $5\,\%$. The incompressible analysis of chapters 11 and 12 is
    adequate.
  \item \engterm{Subsonic compressible} ($0.3 < M < 0.7$): density
    variations matter, but the flow is everywhere subsonic. The
    Prandtl-Glauert correction to incompressible results captures
    the leading-order compressibility effect on lift and
    pressure distribution.
  \item \engterm{Transonic} ($0.7 < M < 1.2$): subsonic and
    supersonic regions coexist on the body. Local supersonic
    pockets on the upper surface of a transonic wing terminate in
    shock waves. The most analytically difficult regime.
  \item \engterm{Supersonic} ($1.2 < M < 5$): the flow is
    everywhere supersonic; shocks attach to leading edges, and
    the governing equations are hyperbolic.
  \item \engterm{Hypersonic} ($M > 5$): real-gas effects (chemical
    dissociation, ionisation, radiation) and viscous-inviscid
    coupling at the wall become important.
\end{itemize}

The boundaries are conventional, not sharp; transonic effects can
begin below $M = 0.7$ on highly cambered or thick wings, and
hypersonic effects can begin below $M = 5$ in some specific
geometries. The classification is a working organisation of the
problems, not a physical division.

\engterm{Mach cone and Mach angle}. A point disturbance moving at
supersonic speed $u > a$ leaves a conical wave front behind it.
The angle the cone makes with the direction of motion is the
\engterm{Mach angle}
\[
\mu = \sin^{-1}(1/M).
\]
At $M = 1$, $\mu = 90^{\circ}$ (the wave front is a plane normal
to the direction of motion); at $M = 2$, $\mu = 30^{\circ}$; at
$M = 5$, $\mu \approx 11.5^{\circ}$. The Mach cone is the analogue
of the Cherenkov cone of charged particles moving faster than light
in a dielectric. The intersection of the Mach cone with the
ground produces the sonic boom heard by observers on the surface
when a supersonic aircraft passes overhead.

% Figure: Mach cone / wavefront construction for a point source moving
% at subsonic, sonic, and supersonic speed.
\begin{figure}[htbp]
\centering
\begin{tikzpicture}[scale=0.95]
% three panels side by side
% Panel 1: subsonic M=0.5 == circles not overlapping ahead, source ahead of own waves
\begin{scope}[xshift=0cm]
  \node[font=\scriptsize, anchor=south] at (0,2.0) {$M<1$};
  % source positions along x; waves emitted at earlier times, radii grow
  \foreach \i/\r/\cx in {3/1.5/0, 2/1.0/0.5, 1/0.5/1.0} {
    \draw[engblue] (\cx,0) circle (\r);
  }
  \fill[engred] (1.5,0) circle (1.2pt);
  \draw[engblack, -{Latex[length=1.5mm]}] (-1.6,0) -- (2.1,0);
\end{scope}
% Panel 2: sonic M=1 == all circles tangent at the source point, plane front
\begin{scope}[xshift=5.2cm]
  \node[font=\scriptsize, anchor=south] at (0,2.0) {$M=1$};
  \foreach \r/\cx in {1.5/0, 1.0/0.5, 0.5/1.0} {
    \draw[engblue] (\cx,0) circle (\r);
  }
  \fill[engred] (1.5,0) circle (1.2pt);
  \draw[engamber, very thick] (1.5,-1.6) -- (1.5,1.6);
  \draw[engblack, -{Latex[length=1.5mm]}] (-1.6,0) -- (2.1,0);
\end{scope}
% Panel 3: supersonic M=2 == circles enveloped by a cone, Mach angle mu
\begin{scope}[xshift=10.6cm]
  \node[font=\scriptsize, anchor=south] at (0,2.0) {$M>1$};
  \foreach \r/\cx in {1.5/0, 1.0/0.5, 0.5/1.0} {
    \draw[engblue] (\cx,0) circle (\r);
  }
  \fill[engred] (2.0,0) circle (1.2pt);
  % Mach lines tangent to the family: mu = asin(1/M)=30 deg for M=2
  \draw[engred, very thick] (2.0,0) -- ({2.0-3.2*cos(30)},{3.2*sin(30)});
  \draw[engred, very thick] (2.0,0) -- ({2.0-3.2*cos(30)},{-3.2*sin(30)});
  \draw[engblack, -{Latex[length=1.5mm]}] (-1.6,0) -- (2.6,0);
  % Mach angle arc
  \draw[engred] (1.1,0) arc (180:150:0.9);
  \node[engred, font=\scriptsize] at (0.95,0.35) {$\mu$};
\end{scope}
\end{tikzpicture}
\caption[Wavefronts and the Mach cone]{Wavefronts emitted by a point
source moving left to right, drawn at equal time intervals. At
subsonic speed ($M<1$) the source stays inside the family of expanding
spheres and the disturbance reaches every point ahead of it. At $M=1$
the wavefronts pile up into a plane normal to the motion at the
source, the limiting case in which upstream signalling just fails. At
supersonic speed ($M>1$) the source outruns its own sound and the
spheres are enveloped by a cone whose half-angle is the Mach angle
$\mu=\sin^{-1}(1/M)$ (here $30^{\circ}$ for $M=2$). Outside the cone
the flow has received no signal from the source, the geometric
statement that supersonic flow cannot communicate disturbances
upstream. The intersection of this cone with the ground is the path
along which a sonic boom is heard.}
\label{fig:vol03ch13:mach-cone}
\end{figure}


The wavefront construction in figure~\ref{fig:vol03ch13:mach-cone}
makes the regime distinction geometric rather than algebraic. Below
$M = 1$ the source stays inside its own family of expanding spheres,
so a disturbance reaches points ahead of the body and the upstream
flow ``knows'' the body is coming. At $M = 1$ the wavefronts pile into
a plane. Above $M = 1$ the body outruns its sound, the spheres are
swept into the Mach cone, and no signal reaches the region outside
the cone. This loss of upstream signalling is the geometric content of
the change from elliptic to hyperbolic governing equations.

\begin{principle}[Compressibility scales with Mach number squared]
The density variation in a compressible flow scales as
$\Delta\rho/\rho \sim M^2/2$ to leading order. Below $M = 0.3$,
the variation is below $5\,\%$ and the incompressible analysis is
adequate; above $M = 0.3$, compressibility corrections matter;
above $M = 1$, the qualitative behaviour changes through the
appearance of shock waves and the loss of upstream signal
propagation. The Mach number is the working dimensionless
parameter; the speed of sound is the working speed scale.
\end{principle}

\section{Isentropic flow relations}\pathtag{core}

For an ideal gas undergoing a steady, adiabatic, reversible
(isentropic) flow, the local thermodynamic state is related to the
local Mach number through a set of standard relations. The starting
point is the steady-flow energy equation, which for an adiabatic
flow with no shaft work states that the total enthalpy per unit mass
is conserved along a streamline. The first law for a steady flow
through a control volume, with no heat added ($\dot Q = 0$) and no
work done ($\dot W = 0$), reduces to $h + \tfrac{1}{2}u^2 =
\text{constant}$, where $h$ is the specific enthalpy and the
$\tfrac{1}{2}u^2$ term is the kinetic energy per unit mass. For a
calorically perfect gas $h = c_p T$, so the energy equation takes
the form
\[
c_p T + \tfrac{1}{2}u^2 = c_p T_0,
\]
where $T_0$ is the \engterm{stagnation temperature} (the
temperature that would be measured if the flow were brought
adiabatically to rest, so that all the kinetic energy reappears as
enthalpy). To turn this into a Mach-number relation, divide through
by $c_p T$:
\[
1 + \frac{u^2}{2 c_p T} = \frac{T_0}{T}.
\]
The kinetic term carries the Mach number once two substitutions are
made. First, $c_p = \gamma R/(\gamma - 1)$, which follows from
$c_p - c_v = R$ and $\gamma = c_p/c_v$ (subtracting gives $c_v
= R/(\gamma - 1)$ and hence $c_p = \gamma R/(\gamma - 1)$). Second,
$u^2 = M^2 a^2 = M^2 \gamma R T$ from the speed-of-sound result.
Then
\[
\frac{u^2}{2 c_p T} = \frac{M^2 \gamma R T}{2\,[\gamma R/(\gamma -
1)]\,T} = \frac{(\gamma - 1)}{2}M^2,
\]
the gas constant and the temperature both cancelling, and
\[
\frac{T_0}{T} = 1 + \frac{\gamma - 1}{2} M^2.
\]
This single relation is the backbone of compressible flow: it says
the stagnation temperature is fixed along an adiabatic streamline
while the static temperature falls as the flow speeds up, the
kinetic share of the fixed total growing as $M^2$
(figure~\ref{fig:vol03ch13:energy-partition}). The corresponding
pressure and density relations follow from the isentropic condition
$p/\rho^\gamma = \text{constant}$ and the ideal-gas law $p = \rho R
T$:
\[
\frac{p_0}{p} = \left(1 + \frac{\gamma - 1}{2} M^2\right)^{\gamma/
(\gamma - 1)},
\quad
\frac{\rho_0}{\rho} = \left(1 + \frac{\gamma - 1}{2} M^2\right)^{1/
(\gamma - 1)}.
\]
The exponents are not free choices; they come from the isentrope.
Along $p/\rho^\gamma = \text{constant}$ the pressure and temperature
are tied by $T \propto p^{(\gamma - 1)/\gamma}$, which is the
isentrope written in $T$ and $p$ (eliminate $\rho$ between
$p\rho^{-\gamma} = \text{const}$ and $p = \rho R T$: $\rho \propto
p^{1/\gamma}$, so $T = p/(\rho R) \propto p^{1 - 1/\gamma} =
p^{(\gamma - 1)/\gamma}$). Inverting, $p \propto T^{\gamma/(\gamma -
1)}$, so the stagnation-to-static pressure ratio is the
stagnation-to-static temperature ratio raised to $\gamma/(\gamma -
1)$, which is the first relation above. The density relation follows
the same way from $\rho \propto p^{1/\gamma} \propto T^{1/(\gamma -
1)}$. These are the \engterm{isentropic flow relations}; they are
tabulated in every aerodynamics textbook as a function of $M$ for
$\gamma = 1.4$. The stagnation values are conserved along
isentropic streamlines; they are the natural reference states for
compressible-flow problems.

% Figure: the energy-equation bookkeeping. The stagnation temperature
% is fixed along an adiabatic streamline; as the Mach number rises the
% static temperature falls and the kinetic-energy share rises, their
% sum held constant at T0. Plotted as T/T0 and the kinetic share
% (T0-T)/T0 versus Mach for gamma=1.4.
\begin{figure}[htbp]
\centering
\begin{tikzpicture}
\begin{axis}[
  width=12cm, height=8cm,
  xlabel={Mach number $M$},
  ylabel={fraction of stagnation temperature},
  xmin=0, xmax=4, ymin=0, ymax=1.05,
  axis lines=left,
  grid=major, grid style={enggray!20},
  tick label style={font=\scriptsize},
  label style={font=\small},
  legend style={font=\scriptsize, at={(0.97,0.5)}, anchor=east},
]
% T/T0 = 1/(1 + 0.2 M^2)
\addplot[engblue, very thick, domain=0:4, samples=200] {
  1/(1 + 0.2*x^2) };
\addlegendentry{static share $T/T_0$}
% kinetic share = 1 - T/T0 = (0.2 M^2)/(1 + 0.2 M^2)
\addplot[engred, very thick, domain=0:4, samples=200] {
  (0.2*x^2)/(1 + 0.2*x^2) };
\addlegendentry{kinetic share $u^2/(2 c_p T_0)$}
% the two always sum to 1 (the constant line)
\addplot[enggray, dashed, domain=0:4, samples=2] {1};
\node[enggray, font=\scriptsize, anchor=south east] at (axis cs:3.9,1.0) {sum $=1$ (constant $T_0$)};
% M=1 marker
\draw[englime, dashed] (axis cs:1,0) -- (axis cs:1,1);
\node[englime, font=\scriptsize, anchor=south, rotate=90] at (axis cs:1.0,0.30) {$M=1$};
% crossover where T/T0 = kinetic share = 0.5, at M = sqrt(5) approx 2.236
\addplot[engblack, only marks, mark=*, mark options={fill=engblack}] coordinates {(2.236,0.5)};
\node[engblack, font=\scriptsize, anchor=south west] at (axis cs:2.30,0.50) {$M=\sqrt{5}$: equal split};
\end{axis}
\end{tikzpicture}
\caption[Energy partition along an adiabatic streamline]{The energy
equation $c_p T + \tfrac{1}{2}u^2 = c_p T_0$ read as a partition. Along
an adiabatic streamline the stagnation temperature $T_0$ is fixed, so
the static share $T/T_0$ and the kinetic share $u^2/(2 c_p T_0)$ must
sum to one at every Mach number. At low Mach almost all of the
enthalpy is thermal; the kinetic share grows as $M^2$ and overtakes the
thermal share at $M=\sqrt{5}\approx 2.24$, where the two are equal.
The static temperature a thermometer would read in the moving flow is
$T_0$ minus the kinetic share, which is why a high-Mach stream is cold
even though its stagnation temperature is high. The curves are
$T/T_0 = (1 + \tfrac{\gamma-1}{2}M^2)^{-1}$ and its complement for
$\gamma=1.4$.}
\label{fig:vol03ch13:energy-partition}
\end{figure}


Figure~\ref{fig:vol03ch13:energy-partition} reads the energy equation
as a partition of the fixed stagnation temperature into a thermal
share $T/T_0$ and a kinetic share that grows as $M^2$. The two shares
are equal at $M = \sqrt{5} \approx 2.24$ for air; above it the flow
carries more of its energy as bulk motion than as heat. The picture
explains why a high-Mach stream is cold: a thermometer drifting with
the flow reads the static temperature, which is the stagnation value
minus everything the kinetic share has taken, so the air at the cruise
Mach of a transport sits tens of degrees below its stagnation
temperature even though the nose, where the flow is brought to rest,
runs warm.

The compressible Bernoulli equation, which replaces the
incompressible Bernoulli of chapter 11, takes the form
\[
\frac{u^2}{2} + \frac{\gamma}{\gamma - 1}\frac{p}{\rho} =
\frac{\gamma}{\gamma - 1}\frac{p_0}{\rho_0}.
\]
For low Mach number ($M \ll 1$), expanding to leading order in
$M^2$ recovers the incompressible form $p + \tfrac{1}{2}\rho u^2
= p_0$.

\engterm{Critical (sonic) conditions}. The state at which the flow
reaches $M = 1$ is the \engterm{critical state}, denoted with a
star. From the isentropic relations,
\[
\frac{T^*}{T_0} = \frac{2}{\gamma + 1}, \quad
\frac{p^*}{p_0} = \left(\frac{2}{\gamma + 1}\right)^{\gamma/(\gamma
- 1)}, \quad
\frac{\rho^*}{\rho_0} = \left(\frac{2}{\gamma + 1}\right)^{1/(\gamma
- 1)}.
\]
For air with $\gamma = 1.4$: $T^*/T_0 = 0.833$, $p^*/p_0 = 0.528$,
$\rho^*/\rho_0 = 0.634$. The pressure must drop to about $53\,\%$
of the stagnation value for the flow to reach Mach $1$. This is
the \engterm{critical pressure ratio}, the working parameter in
nozzle and orifice design. The companion code file
\texttt{code/isentropic.py} evaluates all four ratios at any Mach
number and tabulates the critical state; the table it produces is
reproduced as \texttt{data/isentropic\_table.csv} and underlies the
worked examples below.

\engterm{Area-Mach relation}. For a steady one-dimensional
isentropic flow through a duct of cross-sectional area $A$, the
continuity equation $\rho u A = \text{constant}$ ties the area to
the Mach number. Before integrating it to the algebraic form, the
differential form repays the effort because it carries the physics
of why a supersonic flow needs a diverging duct. Take the logarithm
of $\rho u A = \text{const}$ and differentiate: $d\rho/\rho + du/u
+ dA/A = 0$. The isentropic momentum balance along a streamline is
$dp = -\rho u\,du$ (Euler's equation with no body force), and $dp =
a^2 d\rho$ from the speed-of-sound result, so $d\rho/\rho =
-(u/a^2)\,du = -M^2 (du/u)$. Substituting into the continuity
differential,
\[
-M^2\frac{du}{u} + \frac{du}{u} + \frac{dA}{A} = 0
\quad\Longrightarrow\quad
\frac{dA}{A} = (M^2 - 1)\frac{du}{u}.
\]
This compact result is the heart of nozzle design. For subsonic
flow ($M < 1$) the factor $M^2 - 1$ is negative, so the velocity
rises ($du > 0$) only where the area falls ($dA < 0$): a subsonic
flow accelerates in a contraction, the everyday intuition. For
supersonic flow ($M > 1$) the factor is positive, so the velocity
rises only where the area rises: a supersonic flow accelerates in
an expansion, against the everyday intuition, because the density
falls faster than the velocity rises and the area must grow to pass
the same mass. At $M = 1$ the factor vanishes, so $dA = 0$: sonic
flow can occur only where the area is stationary, which is the
throat. Integrating the continuity relation with the isentropic
relations supplies the algebraic \engterm{area-Mach relation}
\[
\frac{A}{A^*} = \frac{1}{M}\left[\frac{2}{\gamma + 1}\left(1 +
\frac{\gamma - 1}{2}M^2\right)\right]^{(\gamma + 1)/[2(\gamma -
1)]}.
\]
The area ratio $A/A^*$ is a non-monotonic function of $M$: it is
$1$ at $M = 1$ (the throat condition), large at $M \ll 1$
(subsonic flow needs a large area for low velocity), and large
again at $M \gg 1$ (supersonic flow needs a large area for high
velocity because the density has dropped). The minimum area
corresponds to $M = 1$.

% Figure: isentropic area-Mach relation A/A* versus M for gamma=1.4.
% Non-monotonic: large at low M, minimum 1 at M=1, large again at high M.
% The same area ratio admits one subsonic and one supersonic solution.
\begin{figure}[htbp]
\centering
\begin{tikzpicture}
\begin{axis}[
  width=12cm, height=8cm,
  xlabel={Mach number $M$}, ylabel={area ratio $A/A^{*}$},
  xmin=0, xmax=4, ymin=0, ymax=11,
  axis lines=left,
  grid=major, grid style={enggray!20},
  tick label style={font=\scriptsize},
  label style={font=\small},
  legend style={font=\scriptsize, at={(0.97,0.97)}, anchor=north east},
]
% A/A* = (1/M)*[ (2/(g+1))(1 + (g-1)/2 M^2) ]^((g+1)/(2(g-1)))
% g=1.4: exponent = 2.4/0.8 = 3.0; factor 2/2.4 = 0.833333
% A/A* = (1/M)*(0.833333*(1+0.2*M^2))^3
\addplot[engblue, very thick, domain=0.08:4, samples=200] {
  (1/x)*(0.833333*(1+0.2*x^2))^3 };
% M=1 marker at A/A*=1
\addplot[engred, only marks, mark=*, mark options={fill=engred}] coordinates {(1,1)};
\node[engred, font=\scriptsize, anchor=south west] at (axis cs:1.05,1.1) {$M=1$, $A/A^{*}=1$};
% example: A/A*=2 admits two solutions; horizontal line
\draw[enggray, dashed] (axis cs:0,2) -- (axis cs:4,2);
\node[enggray, font=\scriptsize, anchor=south] at (axis cs:0.3,2.05) {$A/A^{*}=2$};
% the two intersections of A/A*=2: M approx 0.31 (subsonic), 2.20 (supersonic)
\addplot[engblack, only marks, mark=square, mark options={draw=engblack}] coordinates {(0.31,2) (2.20,2)};
\node[engblack, font=\scriptsize, anchor=west] at (axis cs:0.33,2.5) {subsonic};
\node[engblack, font=\scriptsize, anchor=west] at (axis cs:2.25,2.5) {supersonic};
\end{axis}
\end{tikzpicture}
\caption[Isentropic area-Mach relation]{The area-Mach relation
$A/A^{*}(M)$ for $\gamma=1.4$. The curve is non-monotonic: it
diverges as $M\to 0$ (a slow subsonic flow needs a large area), takes
its minimum value of $1$ at $M=1$ (the throat), and rises again as
$M\to\infty$ (a fast supersonic flow needs a large area because the
density has dropped). A given area ratio above $1$ therefore admits
two isentropic solutions, one subsonic and one supersonic; for
$A/A^{*}=2$ these are $M\approx 0.31$ and $M\approx 2.20$. Which
branch a nozzle realises is set by the back pressure. The plotted
curve is the exact relation; the two square markers are the roots at
$A/A^{*}=2$.}
\label{fig:vol03ch13:area-mach}
\end{figure}


Figure~\ref{fig:vol03ch13:area-mach} plots the relation. The single
most useful consequence is that one area ratio above unity admits two
isentropic Mach numbers, one subsonic and one supersonic. A given
converging-diverging nozzle therefore has two isentropic operating
solutions for its fixed exit-to-throat area ratio; which one the
hardware realises is decided by the back pressure, treated in the
next section.

The area-Mach relation governs nozzle and diffuser design in
compressible flow. A subsonic flow accelerates through a contracting
duct; if the contraction is severe enough that the throat reaches
$M = 1$, the flow \engterm{chokes}: the mass flow rate becomes
independent of downstream pressure, and further reduction in
downstream pressure does not increase the flow. A
converging-diverging nozzle can accelerate the flow supersonically
past the throat; the divergence after the throat is necessary
because the density drops and the area must grow to maintain mass
conservation. This is the essential geometry of every rocket
nozzle, jet engine afterburner nozzle, and supersonic wind-tunnel
diffuser.

\section{Normal and oblique shocks}\pathtag{core}

A finite-amplitude pressure disturbance in a compressible flow
does not propagate isentropically. Instead, the disturbance steepens
as it propagates (the high-pressure regions move faster than the
low-pressure regions, because the local speed of sound and the
flow speed both depend on the local state), and it becomes a
\engterm{shock wave}: a near-discontinuity in pressure, density,
temperature, and velocity, across which the entropy increases.

\engterm{Normal shock}. A planar shock perpendicular to the flow.
Mass, momentum, and energy conservation across the shock, applied
to a thin control volume that straddles it, give the
\engterm{Rankine-Hugoniot relations}. For an ideal gas with the
upstream state denoted by subscript $1$ and the downstream state
by subscript $2$:

\begin{itemize}[leftmargin=*]
  \item Mass: $\rho_1 u_1 = \rho_2 u_2$.
  \item Momentum: $p_1 + \rho_1 u_1^2 = p_2 + \rho_2 u_2^2$.
  \item Energy: $h_1 + u_1^2/2 = h_2 + u_2^2/2$, with $h = c_p T$
    for an ideal gas.
\end{itemize}

The three conservation laws plus the ideal-gas state equation are
four relations in the four downstream unknowns ($\rho_2$, $u_2$,
$p_2$, $T_2$), so the downstream state is determined by the upstream
state. The algebra that reduces them to functions of $M_1$ alone is
worth carrying through once, because the same steps recur in the
oblique-shock and moving-shock problems. Write the momentum equation
as $p_1 + \rho_1 u_1^2 = p_2 + \rho_2 u_2^2$ and use $\rho u^2 =
\gamma p M^2$ (since $u^2 = M^2 a^2 = M^2 \gamma p/\rho$, so $\rho
u^2 = \gamma p M^2$) to get
\[
p_1(1 + \gamma M_1^2) = p_2(1 + \gamma M_2^2),
\qquad
\frac{p_2}{p_1} = \frac{1 + \gamma M_1^2}{1 + \gamma M_2^2}.
\]
The energy equation $h_1 + u_1^2/2 = h_2 + u_2^2/2$ states that the
stagnation temperature is conserved, $T_{01} = T_{02}$, so the
isentropic temperature relation applied on each side gives
\[
\frac{T_2}{T_1} = \frac{1 + \frac{\gamma - 1}{2}M_1^2}
{1 + \frac{\gamma - 1}{2}M_2^2}.
\]
The mass equation $\rho_1 u_1 = \rho_2 u_2$ supplies a third ratio.
Writing $\rho = p/(RT)$ and $u = M\sqrt{\gamma R T}$, mass becomes
$p_1 M_1/\sqrt{T_1} = p_2 M_2/\sqrt{T_2}$, that is $(p_2/p_1)
(M_2/M_1)\sqrt{T_1/T_2} = 1$. Substituting the pressure and
temperature ratios just found and squaring leaves a single equation
in $M_1$ and $M_2$. It factors: one root is the trivial $M_2 = M_1$
(no shock), and the other is the shock relation
\[
M_2^2 = \frac{(\gamma - 1)M_1^2 + 2}{2\gamma M_1^2 - (\gamma - 1)}.
\]
Feeding this $M_2^2$ back into the pressure ratio collapses it to
the explicit form
\[
\frac{p_2}{p_1} = 1 + \frac{2\gamma}{\gamma + 1}(M_1^2 - 1),
\]
and the density and temperature ratios follow from mass and the
ideal-gas law,
\[
\frac{\rho_2}{\rho_1} = \frac{(\gamma + 1)M_1^2}{(\gamma - 1)M_1^2
+ 2}, \quad
\frac{T_2}{T_1} = \frac{p_2}{p_1}\cdot\frac{\rho_1}{\rho_2}.
\]
The two-root structure is the algebraic statement that a uniform
supersonic stream admits exactly one shocked downstream state and
one shock-free continuation, and the second law selects the shocked
state over the (entropy-decreasing) expansion that the same algebra
would otherwise permit. The relations are tabulated as the
\engterm{normal-shock tables}
in every aerodynamics reference. The code file
\texttt{code/normal\_shock.py} evaluates them, including the
stagnation-pressure ratio $p_{02}/p_{01}$, and
\texttt{data/normal\_shock\_table.csv} carries the values for $M_1$
from $1.1$ to $8$.

% Figure: normal-shock property ratios versus upstream Mach number M1,
% gamma=1.4. p2/p1, rho2/rho1, T2/T1 rise; M2 falls below 1.
\begin{figure}[htbp]
\centering
\begin{tikzpicture}
\begin{axis}[
  width=12cm, height=8cm,
  xlabel={upstream Mach number $M_1$}, ylabel={ratio across shock},
  xmin=1, xmax=4, ymin=0, ymax=20,
  axis lines=left,
  grid=major, grid style={enggray!20},
  tick label style={font=\scriptsize},
  label style={font=\small},
  legend style={font=\scriptsize, at={(0.03,0.97)}, anchor=north west},
  legend cell align=left,
]
% p2/p1 = 1 + (2g/(g+1))(M1^2 - 1), g=1.4 -> 1 + (2.8/2.4)(M^2-1) = 1 + 1.16667(M^2-1)
\addplot[engred, very thick, domain=1:4, samples=120] { 1 + 1.16667*(x^2 - 1) };
\addlegendentry{$p_2/p_1$}
% rho2/rho1 = (g+1)M^2 / ((g-1)M^2 + 2) = 2.4 M^2 / (0.4 M^2 + 2)
\addplot[engblue, very thick, domain=1:4, samples=120] { 2.4*x^2 / (0.4*x^2 + 2) };
\addlegendentry{$\rho_2/\rho_1$}
% T2/T1 = (p2/p1)*(rho1/rho2)
\addplot[engamber, very thick, domain=1:4, samples=120] {
  (1 + 1.16667*(x^2 - 1)) * (0.4*x^2 + 2) / (2.4*x^2) };
\addlegendentry{$T_2/T_1$}
% M2 = sqrt( ((g-1)M^2 + 2) / (2g M^2 - (g-1)) ) = sqrt((0.4M^2+2)/(2.8M^2-0.4))
\addplot[englime, very thick, domain=1:4, samples=120] {
  sqrt( (0.4*x^2 + 2) / (2.8*x^2 - 0.4) ) };
\addlegendentry{$M_2$ (downstream)}
% M2 = 1 reference
\draw[enggray, dashed] (axis cs:1,1) -- (axis cs:4,1);
\node[enggray, font=\scriptsize, anchor=south east] at (axis cs:4,1.0) {$M_2=1$};
% density limit annotation: rho2/rho1 -> 6 as M->inf
\draw[enggray, dotted] (axis cs:1,6) -- (axis cs:4,6);
\node[enggray, font=\scriptsize, anchor=south east] at (axis cs:4,6.0) {density limit $6$};
\end{axis}
\end{tikzpicture}
\caption[Normal-shock property ratios]{Property ratios across a normal
shock in air ($\gamma=1.4$) as functions of the upstream Mach number
$M_1$. Pressure (red) rises without bound. Temperature (amber) rises
without bound. Density (blue) rises but saturates at the limiting
ratio $(\gamma+1)/(\gamma-1)=6$ as $M_1\to\infty$: a normal shock
cannot compress a perfect gas beyond six times its upstream density,
which is why hypersonic shock layers stay thick relative to a naive
incompressible expectation. The downstream Mach number (green) falls
below $1$ for every $M_1>1$: the flow always leaves a normal shock
subsonic. The pressure, density, and temperature ratios are the
Rankine-Hugoniot relations; the downstream Mach number follows from
mass, momentum, and energy conservation across the discontinuity.}
\label{fig:vol03ch13:normal-shock}
\end{figure}


Figure~\ref{fig:vol03ch13:normal-shock} plots the ratios.
Pressure and temperature climb without bound as $M_1$ grows, but the
density ratio saturates at $(\gamma+1)/(\gamma-1) = 6$ for air: a
perfect-gas normal shock cannot compress the gas past six times its
upstream density no matter how strong it is. This bound is one reason
a real hypersonic shock layer (where dissociation raises the effective
$\gamma$ ceiling) behaves so differently from the perfect-gas
prediction.

Several properties of normal shocks are worth noting:

\begin{itemize}[leftmargin=*]
  \item A normal shock only exists for $M_1 > 1$ (supersonic
    upstream). A subsonic-to-subsonic ``compression shock'' does
    not exist; the second law forbids it (the entropy would
    decrease across such a shock).
  \item The downstream Mach number $M_2 < 1$: the flow becomes
    subsonic after the shock.
  \item The pressure, density, and temperature all rise across the
    shock; the velocity drops.
  \item The stagnation temperature is conserved ($T_{02} = T_{01}$,
    since the flow is adiabatic) but the stagnation pressure
    decreases ($p_{02} < p_{01}$, since the flow is not
    isentropic; entropy increases across the shock).
  \item For weak shocks ($M_1 \to 1^+$), the entropy increase
    scales as $(M_1^2 - 1)^3$, so weak shocks are nearly
    isentropic. For strong shocks ($M_1 \gg 1$), the entropy
    increase is substantial.
\end{itemize}

The stagnation-pressure loss across a normal shock is the quantity that
matters most for engine and inlet design, because stagnation pressure
is the currency a compressor or a thrust nozzle spends. The loss
$p_{02}/p_{01}$ falls from $1$ at $M_1 = 1$ to about $0.72$ at
$M_1 = 2$, $0.33$ at $M_1 = 3$, and $0.06$ at $M_1 = 5$: a single
normal shock at the inlet of a Mach-3 engine would discard two thirds
of the stagnation pressure before the air ever reached the compressor.
This steep loss is why high-Mach inlets decelerate the flow through a
series of weak oblique shocks rather than one strong normal shock, each
weak shock losing little because its entropy rise scales as the cube of
$(M_n^2 - 1)$. Worked example~5 carries the stagnation-loss arithmetic
through a wind-tunnel diffuser, and the supersonic-inlet design
exercise stacks weak shocks to beat the single-normal-shock recovery.

% Figure: oblique-shock geometry on a wedge. Upstream flow M1 turns by
% deflection angle theta across a shock inclined at angle beta.
\begin{figure}[htbp]
\centering
\begin{tikzpicture}[scale=1.0]
% wedge: lower surface horizontal, upper surface inclined at theta=15 deg
% apex at origin, opening to the right
\def\theta{15}
\def\beta{32}
% wedge body (filled gray)
\fill[enggray!20] (0,0) -- (5,0) -- (5,{5*tan(\theta)}) -- cycle;
\draw[engblack, thick] (0,0) -- (5,0);
\draw[engblack, thick] (0,0) -- (5,{5*tan(\theta)});
% incoming flow arrows (horizontal) from the left toward the wedge tip
\foreach \y in {0.5,1.1,1.7} {
  \draw[engblue, -{Latex[length=2mm]}] (-2.3,\y) -- (-0.9,\y);
}
\node[engblue, font=\scriptsize, anchor=east] at (-2.35,1.1) {$M_1$};
% oblique shock line from apex at angle beta above the lower surface
\draw[engred, very thick] (0,0) -- ({4.2*cos(\beta)},{4.2*sin(\beta)});
\node[engred, font=\scriptsize, anchor=south west] at ({4.2*cos(\beta)},{4.2*sin(\beta)}) {shock};
% deflected flow downstream (parallel to upper wedge surface), above wedge
\foreach \y in {2.6,3.1} {
  \draw[engorange, -{Latex[length=2mm]}] ({0.9*\y/3},{\y}) -- ({0.9*\y/3 + 1.3*cos(\theta)},{\y + 1.3*sin(\theta)});
}
\node[engorange, font=\scriptsize, anchor=south] at (2.9,3.4) {$M_2$ (deflected by $\theta$)};
% angle beta arc at apex (between lower surface and shock)
\draw[engred] (0.9,0) arc (0:\beta:0.9);
\node[engred, font=\scriptsize] at ({1.15*cos(\beta/2)},{1.15*sin(\beta/2)}) {$\beta$};
% angle theta arc (between lower surface direction and upper wedge surface)
\draw[engblack] (2.0,0) arc (0:\theta:2.0);
\node[engblack, font=\scriptsize] at (2.4,{2.4*tan(\theta/2)}) {$\theta$};
\end{tikzpicture}
\caption[Oblique-shock geometry on a wedge]{An attached oblique shock
on a wedge. The supersonic upstream flow $M_1$ meets the compression
corner and turns through the deflection angle $\theta$ (the wedge
half-angle) across a shock inclined at the shock angle $\beta$ to the
upstream direction. The velocity component normal to the shock
behaves exactly like a normal shock and is compressed; the component
tangential to the shock is unchanged. For each $M_1$ and each
$\theta$ below the maximum $\theta_{\max}(M_1)$ the
$\theta$=$\beta$=$M$ relation gives two shock angles, a weak shock
(the smaller $\beta$ drawn here) and a strong shock; physical flows
over an unconfined wedge take the weak branch. For $\theta>
\theta_{\max}$ no attached oblique shock exists and the shock detaches
into a curved bow shock standing ahead of the body.}
\label{fig:vol03ch13:oblique-shock}
\end{figure}


\engterm{Oblique shock}. A shock inclined at angle $\beta$ to the
upstream flow direction (figure~\ref{fig:vol03ch13:oblique-shock}).
The oblique shock is not a new problem but the normal-shock problem
seen in a tilted frame. Resolve the upstream velocity into a
component normal to the shock, $u_{n1} = u_1 \sin\beta$, and a
component tangential to it, $u_{t1} = u_1 \cos\beta$. The shock is a
thin sheet across which the tangential momentum is unchanged (no
pressure gradient acts along the sheet), so $u_{t2} = u_{t1}$: the
tangential component slides through untouched. The normal component
behaves exactly like a normal shock, governed by the relations above
with the upstream normal Mach number $M_{n1} = M_1 \sin\beta$ in
place of $M_1$. This single observation reduces every oblique-shock
calculation to a normal-shock table lookup at $M_{n1}$, with two
geometric bookkeeping steps around it: form $M_{n1} = M_1\sin\beta$
going in, and recover the deflected downstream Mach number $M_2 =
M_{n2}/\sin(\beta - \theta)$ coming out, since the downstream flow
runs at angle $\beta - \theta$ to the shock. The \engterm{$\theta$-
$\beta$-$M$ relation} (the deflection-shock-angle-Mach relation) is
the geometry of that bookkeeping made explicit. The flow turns by
$\theta$ because the normal component is slowed (by the normal shock)
while the tangential component is preserved, so the velocity vector
swings toward the shock. Writing $\tan\beta = u_{n1}/u_{t1}$ and
$\tan(\beta - \theta) = u_{n2}/u_{t2} = u_{n2}/u_{t1}$, dividing the
two gives $\tan(\beta - \theta)/\tan\beta = u_{n2}/u_{n1} =
\rho_1/\rho_2$, the normal-shock density ratio at $M_{n1}$.
Substituting the density ratio $(\gamma + 1)M_{n1}^2/[(\gamma -
1)M_{n1}^2 + 2]$ with $M_{n1} = M_1\sin\beta$ and applying the
tangent-subtraction identity reduces the relation to its standard
closed form,
\[
\tan\theta = 2\cot\beta\cdot\frac{M_1^2 \sin^2\beta - 1}{M_1^2(\gamma
+ \cos 2\beta) + 2}.
\]
For each $M_1$ and $\theta$ (below a maximum $\theta_{\max}(M_1)$),
the relation has two solutions for $\beta$: a \engterm{weak
shock} (the lower $\beta$, with downstream flow usually still
supersonic) and a \engterm{strong shock} (the higher $\beta$,
with subsonic downstream flow). Real flows typically produce the
weak-shock solution when the downstream is unconfined.

% Figure: theta-beta-M diagram. Deflection angle theta versus shock
% angle beta for several upstream Mach numbers, gamma=1.4. Each curve
% peaks at theta_max; weak branch below the peak, strong branch above.
\begin{figure}[htbp]
\centering
\begin{tikzpicture}
\begin{axis}[
  width=12cm, height=8.5cm,
  xlabel={shock angle $\beta$ (deg)}, ylabel={deflection angle $\theta$ (deg)},
  xmin=0, xmax=90, ymin=0, ymax=46,
  axis lines=left,
  grid=major, grid style={enggray!20},
  tick label style={font=\scriptsize},
  label style={font=\small},
  legend style={font=\scriptsize, at={(0.03,0.97)}, anchor=north west},
  legend cell align=left,
  domain=1:89, samples=300,
]
% theta(beta) = atan( 2 cot(beta) (M^2 sin^2 beta - 1) / (M^2 (g + cos 2beta) + 2) )
% deg(...) wraps the radians from atan-like; pgf trig in degrees: tan, sin, cos take degrees.
% Use: numerator = 2*(cos(b)/sin(b))*(M^2*sin(b)^2 - 1); denom = M^2*(1.4 + cos(2b)) + 2
% theta = atan2-free: rad2deg(atan2? ) == use deg(atan( ratio )) via \rad? Simpler: theta_deg = atan(num/den) gives degrees in pgf.
\addplot[engblue, very thick] {
  atan( ( 2*(cos(x)/sin(x))*((1.5)^2*(sin(x))^2 - 1) ) / ( (1.5)^2*(1.4 + cos(2*x)) + 2 ) ) };
\addlegendentry{$M_1=1.5$}
\addplot[englime, very thick] {
  atan( ( 2*(cos(x)/sin(x))*((2.0)^2*(sin(x))^2 - 1) ) / ( (2.0)^2*(1.4 + cos(2*x)) + 2 ) ) };
\addlegendentry{$M_1=2$}
\addplot[engamber, very thick] {
  atan( ( 2*(cos(x)/sin(x))*((3.0)^2*(sin(x))^2 - 1) ) / ( (3.0)^2*(1.4 + cos(2*x)) + 2 ) ) };
\addlegendentry{$M_1=3$}
\addplot[engred, very thick] {
  atan( ( 2*(cos(x)/sin(x))*((5.0)^2*(sin(x))^2 - 1) ) / ( (5.0)^2*(1.4 + cos(2*x)) + 2 ) ) };
\addlegendentry{$M_1=5$}
\node[engblack, font=\scriptsize] at (axis cs:40,8) {weak (below peak)};
\node[engblack, font=\scriptsize] at (axis cs:74,18) {strong (above peak)};
\end{axis}
\end{tikzpicture}
\caption[The theta-beta-M diagram]{The deflection-shock-angle-Mach
($\theta$=$\beta$=$M$) diagram for $\gamma=1.4$. For a fixed upstream
Mach number each curve gives the flow deflection $\theta$ produced by
a shock at angle $\beta$. Each curve rises from $\theta=0$ at the Mach
angle $\beta=\mu=\sin^{-1}(1/M)$ (an infinitesimally weak Mach wave),
peaks at a maximum deflection $\theta_{\max}(M_1)$, and returns to
$\theta=0$ at $\beta=90^{\circ}$ (the normal shock, which produces no
turning). For a deflection below $\theta_{\max}$ two shock angles
solve the relation: the weak shock on the rising branch left of the
peak and the strong shock on the falling branch right of it. As
$M_1$ rises the family of curves climbs and $\theta_{\max}$ grows,
approaching about $45.6^{\circ}$ as $M_1\to\infty$. A wedge whose
half-angle exceeds $\theta_{\max}(M_1)$ cannot hold an attached
shock; the shock detaches.}
\label{fig:vol03ch13:theta-beta-m}
\end{figure}


Figure~\ref{fig:vol03ch13:theta-beta-m} plots the relation as a
family of curves, one per $M_1$. Each curve rises from the Mach angle
$\mu = \sin^{-1}(1/M_1)$ at zero deflection, peaks at
$\theta_{\max}(M_1)$, and falls back to zero at $\beta = 90^{\circ}$
(the normal shock, which turns the flow not at all). The two roots for
a given deflection are the weak shock on the rising branch and the
strong shock on the falling branch. The solver in
\texttt{code/oblique\_shock.py} brackets each root by bisection and
also returns $\theta_{\max}$ by scanning $\beta$; for $M_1 = 3$ it
gives the weak-shock angle $\beta \approx 32.2^{\circ}$, the strong
$\beta \approx 84.4^{\circ}$, and $\theta_{\max} \approx 34.1^{\circ}$
at $15^{\circ}$ deflection.

% Figure: pressure-deflection shock polar for a fixed upstream Mach
% number. The static-pressure ratio across an oblique shock plotted
% against the flow deflection angle traces a closed loop: the lower
% branch is the weak-shock family, the upper branch the strong-shock
% family, they meet at the maximum-deflection nose, and the top of the
% loop (zero deflection) is the normal shock. Built from a small table of
% (theta, p2/p1) values for M1 = 2.5, computed from the theta-beta-M and
% Rankine-Hugoniot relations; plotted as marked points joined smoothly.
\begin{figure}[htbp]
\centering
\begin{tikzpicture}
\begin{axis}[
  width=11cm, height=8.5cm,
  xlabel={flow deflection $\theta$ (deg)},
  ylabel={pressure ratio $p_2/p_1$},
  xmin=0, xmax=32, ymin=0.5, ymax=8,
  axis lines=left,
  grid=major, grid style={enggray!20},
  tick label style={font=\scriptsize},
  label style={font=\small},
  legend style={font=\scriptsize, at={(0.03,0.97)}, anchor=north west},
]
% Weak branch (rising deflection, modest pressure) for M1 = 2.5:
% theta : p2/p1  (from theta-beta-M weak root then RH pressure jump)
\addplot[engblue, very thick, mark=*, mark size=1.2pt, mark options={fill=engblue}]
  coordinates {
  (0,1.00) (5,1.46) (10,1.94) (15,2.50) (20,3.22) (25,4.30) (29.0,5.6)
};
\addlegendentry{weak shock}
% Strong branch (falling deflection back to normal shock at theta=0):
\addplot[engred, very thick, mark=square*, mark size=1.2pt, mark options={fill=engred}]
  coordinates {
  (29.0,5.6) (25,6.10) (20,6.55) (15,6.90) (10,7.15) (5,7.30) (0,7.13)
};
\addlegendentry{strong shock}
% maximum-deflection nose marker
\addplot[engblack, only marks, mark=diamond*, mark size=2pt, mark options={fill=engblack}]
  coordinates {(29.0,5.6)};
\node[engblack, font=\scriptsize, anchor=west] at (axis cs:29.3,5.6) {$\theta_{\max}$};
% normal-shock point at theta=0 (top of loop)
\node[engred, font=\scriptsize, anchor=west] at (axis cs:0.4,7.13) {normal shock};
% sonic-downstream point sits just below theta_max on the strong branch;
% mark it schematically
\addplot[englime, only marks, mark=triangle*, mark size=2pt, mark options={fill=englime}]
  coordinates {(28.0,6.0)};
\node[englime, font=\scriptsize, anchor=south west] at (axis cs:24.5,6.3) {$M_2=1$};
\end{axis}
\end{tikzpicture}
\caption[Pressure-deflection shock polar]{The pressure-deflection shock
polar for a fixed upstream Mach number ($M_1=2.5$, $\gamma=1.4$). Each
flow deflection $\theta$ below the maximum admits two oblique-shock
solutions, plotted here as a single closed loop in the
pressure-deflection plane. The lower (weak) branch carries the modest
pressure rise a wedge in unconfined flow realises; the upper (strong)
branch carries the larger pressure rise and a subsonic downstream flow.
The two branches meet at the maximum-deflection nose $\theta_{\max}$,
beyond which no attached shock exists and the shock detaches. The top of
the loop at $\theta=0$ is the normal shock, which turns the flow not at
all but gives the largest pressure jump; the bottom at $\theta=0$ is the
zero-strength Mach wave $p_2/p_1=1$. The sonic-downstream point
($M_2=1$) sits on the strong branch just inside $\theta_{\max}$, so a
sliver of the strong-branch solutions near the nose still leave the flow
supersonic. The polar is the graphical form of the
$\theta$-$\beta$-$M$ relation, and overlaying the polars of two streams
that must share a pressure and a direction is the standard way to solve
shock-shock intersections.}
\label{fig:vol03ch13:shock-polar}
\end{figure}


The same two-root structure reappears in a more useful coordinate.
Plotting the pressure ratio $p_2/p_1$ against the deflection angle
$\theta$ at fixed $M_1$ traces the \engterm{shock polar}, a closed
loop (figure~\ref{fig:vol03ch13:shock-polar}). The lower branch is
the weak-shock family, the upper branch the strong-shock family, the
two meet at the maximum-deflection nose, and the top of the loop at
zero deflection is the normal shock. The polar earns its keep on
shock-shock interactions: when two shocks meet, the streams behind
them must share a single static pressure and a single flow
direction, and the meeting point is read off by overlaying the
polars of the two streams and finding where they cross. The
graphical method predates the digital solver and still gives the
clearest account of regular versus Mach reflection, where the polars
fail to cross and a third shock (the Mach stem) appears.

For deflection angles $\theta > \theta_{\max}(M_1)$, no attached
oblique shock exists, and the shock \engterm{detaches} from the
body. A detached shock is a curved bow shock standing some
distance upstream of the body; it is the canonical configuration
of blunt-body supersonic flow (re-entry vehicles, blunt missiles).

\engterm{Shock interactions}. When two oblique shocks meet, they
either reflect off solid boundaries, cancel against an oncoming
expansion fan, or interact with other shocks and expansions. The
analytical machinery for shock-shock interactions, shock reflections
from solid walls, and shock interactions with boundary layers is
the subject of advanced gas-dynamics references
\cite{text:anderson-aerodynamics}; the essentials for this chapter
are that shocks are localised regions of irreversible compression,
they affect downstream flow through their pressure jump and
deflection, and their location and strength can be computed from
the $\theta$-$\beta$-$M$ relation.

\section{Nozzles, diffusers, and choking}\pathtag{standard}

The interaction of the area-Mach relation with the isentropic
relations governs the design of converging and converging-diverging
nozzles and the corresponding diffusers.

\engterm{Converging nozzle}. A duct of area decreasing in the flow
direction, used to accelerate a subsonic flow toward $M = 1$. The
exit Mach number is set by the ratio of upstream stagnation
pressure $p_0$ to ambient (back) pressure $p_b$. If $p_b/p_0 >
p^*/p_0 = 0.528$ (for air), the exit flow is subsonic and matches
the back pressure. If $p_b/p_0 \le 0.528$, the exit reaches $M = 1$
and the flow chokes: the exit pressure equals the critical pressure
$p^* = 0.528\,p_0$, the mass flow rate reaches its maximum value
$\dot{m}_{\max} = \rho^* u^* A_e$, and further reduction in back
pressure does not change the upstream flow. The downstream
adjustment from $p^*$ to $p_b$ occurs through a series of
expansion waves immediately downstream of the exit.

The choked mass flow rate follows from evaluating $\dot m = \rho u A$
at the throat, where $M = 1$. The mass flux $\rho u = \rho M
\sqrt{\gamma R T}$ is written in stagnation variables through the
isentropic relations $\rho = \rho_0 (1 + \frac{\gamma - 1}{2}M^2)^{-1/
(\gamma - 1)}$ and $T = T_0 (1 + \frac{\gamma - 1}{2}M^2)^{-1}$, and
$\rho_0 = p_0/(R T_0)$. Collecting the powers and setting $M = 1$ (so
the bracket becomes $(\gamma + 1)/2$) gives the choked mass flow rate
\[
\dot{m}_{\max} = A_e \frac{p_0}{\sqrt{T_0}}\sqrt{\frac{\gamma}{R}}
\left(\frac{2}{\gamma + 1}\right)^{(\gamma + 1)/[2(\gamma - 1)]}.
\]
The form is worth reading: the choked flow rises linearly with
stagnation pressure and the throat area, falls as the inverse square
root of stagnation temperature (a hotter gas is less dense and passes
less mass at the same pressure), and the bracket is a pure function of
$\gamma$. For air, the constant evaluates to about $0.0404\,\sqrt{R}$
in SI units. A choked converging nozzle is the canonical compressible-flow
mass-flow meter; it delivers a flow rate that depends only on
upstream conditions, independent of any downstream perturbation.

\engterm{Converging-diverging (de Laval) nozzle}. Named for Gustaf
de Laval (1888), who used it in early steam-turbine designs. The
flow accelerates through the converging section, chokes at the
throat ($M = 1$), and accelerates supersonically through the
diverging section. The design exit Mach number is set by the
exit-to-throat area ratio through the area-Mach relation: $A_e/A^*
\approx 1.7$ for $M_e = 2$, $4.0$ for $M_e = 3$, $25$ for $M_e =
5$. The helper \texttt{code/nozzle.py} returns the area ratio for a
target exit Mach number and inverts the relation to recover the Mach
number on either branch.

% Figure: converging-diverging (de Laval) nozzle schematic with the
% Mach-number progression along the axis. Subsonic in the converging
% section, M=1 at the throat, supersonic in the diverging section.
\begin{figure}[htbp]
\centering
\begin{tikzpicture}[scale=1.0]
% nozzle wall, upper and lower (mirror); converging then diverging
% profile y(x): inlet half-height 1.6, throat 0.5, exit 1.3
\def\inlet{1.6}
\def\throat{0.5}
\def\exit{1.3}
% upper wall as a smooth-ish polyline
\draw[engblue, very thick]
  (0,\inlet) -- (1.0,1.3) -- (2.0,0.9) -- (3.0,\throat)
  -- (4.0,0.85) -- (5.0,1.1) -- (6.0,\exit);
% lower wall (mirror)
\draw[engblue, very thick]
  (0,-\inlet) -- (1.0,-1.3) -- (2.0,-0.9) -- (3.0,-\throat)
  -- (4.0,-0.85) -- (5.0,-1.1) -- (6.0,-\exit);
% throat dashed line
\draw[enggray, dashed] (3.0,-\throat) -- (3.0,\throat);
\node[enggray, font=\scriptsize, anchor=south] at (3.0,0.55) {throat $A^{*}$, $M=1$};
% centreline with flow arrow
\draw[engblack, -{Latex[length=2.5mm]}] (-0.4,0) -- (6.4,0);
% region labels
\node[engblack, font=\scriptsize] at (1.3,-0.15) {$M<1$};
\node[engred, font=\scriptsize] at (4.6,-0.15) {$M>1$};
\node[engblack, font=\scriptsize, anchor=east] at (-0.45,0.9) {$p_0,T_0$};
% inset Mach profile below
\begin{scope}[yshift=-3.6cm]
  \draw[engblack, -{Latex[length=2mm]}] (-0.4,0) -- (6.4,0) node[right, font=\scriptsize] {$x$};
  \draw[engblack, -{Latex[length=2mm]}] (-0.4,0) -- (-0.4,2.2) node[above, font=\scriptsize] {$M$};
  % M=1 reference
  \draw[enggray, dashed] (-0.4,1.0) -- (6.4,1.0);
  \node[enggray, font=\scriptsize, anchor=east] at (-0.45,1.0) {$1$};
  % Mach curve: rising from ~0.2 to 1 at throat (x=3) then to ~2.4 at exit
  \draw[engred, very thick] plot[smooth] coordinates {
    (0,0.18) (1.0,0.32) (2.0,0.6) (3.0,1.0)
    (4.0,1.55) (5.0,2.0) (6.0,2.35)};
  \draw[enggray, dashed] (3.0,0) -- (3.0,1.0);
\end{scope}
\end{tikzpicture}
\caption[Converging-diverging nozzle and the Mach progression]{The
converging-diverging (de Laval) nozzle. The subsonic inlet flow
accelerates through the converging section, reaches $M=1$ at the
throat (the minimum area $A^{*}$), and accelerates supersonically
through the diverging section. The lower panel plots the centreline
Mach number against axial position for the design (fully supersonic)
operating mode: it rises monotonically through the throat and beyond.
The divergence after the throat is required because the density falls
faster than the velocity rises once the flow is supersonic, so the
area must grow to conserve mass. This is the geometry of every rocket
nozzle and supersonic wind-tunnel nozzle.}
\label{fig:vol03ch13:cd-nozzle}
\end{figure}


For a given converging-diverging nozzle, the back pressure
determines the operating mode:

\begin{itemize}[leftmargin=*]
  \item \engterm{Design condition}: $p_e = p_b$, the supersonic
    exit is exactly matched. Smooth supersonic flow throughout.
  \item \engterm{Overexpanded}: $p_e < p_b$, the exit pressure is
    below the back pressure. The flow is compressed back to the
    back pressure through oblique shocks at the nozzle exit (if
    the mismatch is mild) or a normal shock inside the diverging
    section (if the mismatch is severe).
  \item \engterm{Underexpanded}: $p_e > p_b$, the exit pressure is
    above the back pressure. The flow expands further through
    Prandtl-Meyer waves outside the nozzle.
  \item \engterm{Subsonic throughout}: back pressure too high for
    the throat to choke. The diverging section acts as a
    subsonic diffuser, decelerating the flow.
\end{itemize}

% Figure: pressure distribution along a converging-diverging nozzle for
% several back-pressure ratios. Subsonic, shock-in-diverging,
% overexpanded, design, underexpanded.
\begin{figure}[htbp]
\centering
\begin{tikzpicture}
\begin{axis}[
  width=12cm, height=8cm,
  xlabel={axial position $x$ (throat at $x=0.5$)}, ylabel={$p/p_0$},
  xmin=0, xmax=1, ymin=0, ymax=1.02,
  axis lines=left,
  grid=major, grid style={enggray!20},
  tick label style={font=\scriptsize},
  label style={font=\small},
  legend style={font=\scriptsize, at={(0.97,0.97)}, anchor=north east},
  legend cell align=left,
]
% throat at x=0.5; converging 0..0.5, diverging 0.5..1
% Curve (a): all subsonic, accelerate then decelerate, min at throat ~0.7
\addplot[engblue, very thick] coordinates {
  (0,0.99) (0.2,0.95) (0.35,0.88) (0.5,0.78) (0.65,0.86) (0.8,0.93) (1.0,0.97)};
\addlegendentry{(a) subsonic, no choke}
% Curve (b): choked then shock in diverging section == supersonic dip then jump up
\addplot[englime, very thick] coordinates {
  (0,0.99) (0.2,0.93) (0.35,0.83) (0.5,0.528) (0.6,0.35) (0.7,0.22)};
\addplot[englime, very thick] coordinates {
  (0.7,0.22) (0.701,0.62)};
\addplot[englime, very thick] coordinates {
  (0.701,0.62) (0.85,0.72) (1.0,0.78)};
\addlegendentry{(b) normal shock in nozzle}
% Curve (c): design, smooth supersonic to low exit pressure
\addplot[engred, very thick] coordinates {
  (0,0.99) (0.2,0.93) (0.35,0.83) (0.5,0.528) (0.65,0.25) (0.8,0.12) (1.0,0.06)};
\addlegendentry{(c) design (matched)}
% throat marker
\draw[enggray, dashed] (axis cs:0.5,0) -- (axis cs:0.5,1.0);
\node[enggray, font=\scriptsize, anchor=south] at (axis cs:0.5,1.0) {throat};
% critical ratio line
\draw[enggray, dotted] (axis cs:0,0.528) -- (axis cs:1,0.528);
\node[enggray, font=\scriptsize, anchor=south west] at (axis cs:0.02,0.528) {$p^{*}/p_0=0.528$};
\end{axis}
\end{tikzpicture}
\caption[Nozzle operating modes by back pressure]{Centreline pressure
along a converging-diverging nozzle for three back-pressure ratios,
$\gamma=1.4$. (a) The back pressure is high enough that the throat
never chokes; the flow accelerates to a subsonic minimum at the
throat and decelerates in the diverging section, behaving as a venturi.
(b) The throat chokes and the flow goes supersonic past it, then a
normal shock inside the diverging section abruptly raises the pressure
and returns the flow to subsonic, which then decelerates to a higher
exit pressure. (c) At the design back pressure the supersonic branch
runs smoothly to a low matched exit pressure with no internal shock.
Lower back pressures than the design value leave the nozzle
underexpanded, with the final expansion occurring outside the exit
through Prandtl-Meyer waves. The dotted line marks the critical
pressure ratio at which the throat reaches $M=1$.}
\label{fig:vol03ch13:nozzle-modes}
\end{figure}


Figure~\ref{fig:vol03ch13:nozzle-modes} shows the centreline pressure
for three of these modes. The supersonic branch is the same curve in
every choked case until the back pressure forces a normal shock inside
the diverging section; the shock location moves toward the exit as the
back pressure falls, and reaches the exit plane at the start of the
overexpanded regime.

Rocket-engine nozzles are typically designed for a specific exit
altitude (and therefore a specific back pressure). At sea level,
the high ambient pressure overexpands a high-altitude-design
nozzle, producing a series of oblique shocks (the visible diamond
pattern in rocket exhaust plumes). At higher altitude, the
nozzle approaches its design condition and the diamonds
disappear. At very high altitude, the nozzle becomes underexpanded
and the plume widens.

\engterm{Supersonic diffuser}. The dual of a supersonic nozzle: a
converging-diverging duct used to decelerate a supersonic flow.
The supersonic-to-subsonic transition occurs through a normal
shock; the diffuser's design challenge is to position the shock
at the throat so that the shock is as weak as possible (minimum
stagnation-pressure loss). Off-design operation places the shock
either upstream of the throat (overspeed) or downstream of the
throat in the diverging section, with consequences for the
pressure recovery and for the inlet's interaction with whatever
sits downstream.

The inlet of a supersonic jet engine, the test section diffuser of
a supersonic wind tunnel, and the intake of a ramjet all face the
supersonic-diffuser design problem. The discipline involves
detailed shock-position management, often through variable-geometry
inlets that adjust the throat area in response to flight
conditions.

\section{Subsonic, transonic, supersonic, hypersonic aerodynamics}\pathtag{standard}

The four flight regimes have qualitatively different aerodynamics.

\engterm{Subsonic compressible aerodynamics} ($0.3 < M_\infty <
0.7$). The flow is everywhere subsonic, but density variations
modify the pressure distribution and lift. The
\engterm{Prandtl-Glauert correction} relates the compressible-flow
pressure coefficient to its incompressible counterpart:
\[
C_p = \frac{C_{p,0}}{\sqrt{1 - M_\infty^2}},
\]
with $C_{p,0}$ the incompressible value at the same geometry. The
lift coefficient is correspondingly amplified: a flat-plate airfoil
that gives $C_L = 2\pi\alpha$ in incompressible flow gives
$C_L = 2\pi\alpha/\sqrt{1 - M_\infty^2}$ in subsonic compressible
flow. The drag coefficient is largely unchanged in this regime;
wave-drag effects do not yet appear.

\engterm{Transonic aerodynamics} ($0.7 < M_\infty < 1.2$). The
flow develops local supersonic pockets on the upper surface of
the wing, terminating in shocks. The shocks induce boundary-layer
separation downstream (the shock-induced separation problem); the
lift drops sharply and the drag rises sharply through a band of
Mach numbers near $M_\infty = 1$. The phenomenon is the
\engterm{drag divergence} or \engterm{wave drag} associated with
the appearance of shocks; the Mach number at which the drag
coefficient begins to rise sharply is the \engterm{drag-divergence
Mach number} $M_{DD}$, typically $0.7$ to $0.85$ for a conventional
transport airfoil and up to $0.95$ for a supercritical airfoil
(Whitcomb, 1965; the supercritical airfoil's flatter upper surface
delays the shock formation and pushes $M_{DD}$ higher).

The \engterm{area rule} (Whitcomb, 1952) is the design principle
that the transonic drag of an aircraft depends on its longitudinal
distribution of cross-sectional area: an aircraft whose
cross-sectional area distribution is smooth has lower transonic
drag than one with an abrupt area change. The fuselage waist
between the wings of an F-102 fighter (the canonical
``Coke-bottle'' shape) is the visible application of the area
rule.

% Figure: pressure coefficient distribution over a transonic airfoil
% upper surface, showing the supersonic pocket terminated by a shock.
\begin{figure}[htbp]
\centering
\begin{tikzpicture}
\begin{axis}[
  width=12cm, height=8cm,
  xlabel={chordwise position $x/c$},
  ylabel={$-C_p$ (suction up)},
  xmin=0, xmax=1, ymin=-0.8, ymax=1.6,
  axis lines=left,
  y dir=normal,
  grid=major, grid style={enggray!20},
  tick label style={font=\scriptsize},
  label style={font=\small},
  legend style={font=\scriptsize, at={(0.97,0.97)}, anchor=north east},
  legend cell align=left,
]
% Cp* line: pressure coefficient at which local M=1 (sonic), here say -Cp*=0.7
\draw[enggray, dashed] (axis cs:0,0.7) -- (axis cs:1,0.7);
\node[enggray, font=\scriptsize, anchor=south east] at (axis cs:1,0.7) {$C_p^{*}$ (local $M=1$)};
% upper surface: strong leading-edge suction peak, supersonic plateau above Cp*, shock jump down
\addplot[engblue, very thick] coordinates {
  (0,-0.2) (0.03,1.2) (0.08,1.05) (0.15,0.95) (0.30,0.92) (0.45,0.90)
  (0.55,0.92)};
\addplot[engblue, very thick] coordinates {(0.55,0.92) (0.555,0.20)};
\addplot[engblue, very thick] coordinates {
  (0.555,0.20) (0.70,0.05) (0.85,-0.10) (1.0,-0.05)};
\addlegendentry{upper surface}
% lower surface: mild
\addplot[engred, very thick] coordinates {
  (0,-0.2) (0.05,-0.35) (0.2,-0.25) (0.5,-0.10) (0.8,0.0) (1.0,-0.05)};
\addlegendentry{lower surface}
% supersonic pocket label
\node[engblack, font=\scriptsize] at (axis cs:0.30,1.15) {supersonic pocket};
\node[engblack, font=\scriptsize, anchor=west] at (axis cs:0.57,0.55) {shock};
\end{axis}
\end{tikzpicture}
\caption[Transonic airfoil pressure distribution]{Pressure coefficient
over a transonic airfoil at a freestream Mach number above the
critical value, plotted as $-C_p$ so that suction is upward. The upper
surface accelerates the flow past the sonic line $C_p^{*}$ into a
local supersonic pocket; the pocket is closed by a near-normal shock
(the sharp recompression near mid-chord), downstream of which the
flow is subsonic again. The shock raises the pressure abruptly and,
if strong enough, separates the boundary layer behind it, the
mechanism of transonic buffet and drag rise. The lower surface stays
subsonic. The area enclosed between the upper and lower curves is
proportional to the section lift; the shock recompression is the
visible signature that the section is operating in the transonic
regime.}
\label{fig:vol03ch13:airfoil-cp}
\end{figure}


% Figure: drag coefficient versus freestream Mach number showing the
% transonic drag rise (drag divergence) for a conventional vs a
% supercritical airfoil.
\begin{figure}[htbp]
\centering
\begin{tikzpicture}
\begin{axis}[
  width=12cm, height=8cm,
  xlabel={freestream Mach number $M_\infty$},
  ylabel={drag coefficient $C_D$},
  xmin=0.5, xmax=1.1, ymin=0, ymax=0.09,
  axis lines=left,
  grid=major, grid style={enggray!20},
  tick label style={font=\scriptsize},
  label style={font=\small},
  legend style={font=\scriptsize, at={(0.03,0.97)}, anchor=north west},
  legend cell align=left,
]
% conventional airfoil: flat ~0.012 until ~0.72, then sharp rise
\addplot[engblue, very thick] coordinates {
  (0.5,0.012) (0.65,0.0125) (0.72,0.014) (0.76,0.020)
  (0.80,0.034) (0.85,0.055) (0.90,0.072) (0.95,0.080) (1.05,0.078)};
\addlegendentry{conventional airfoil}
% supercritical airfoil: flat to ~0.80 then rise (M_DD pushed higher)
\addplot[engred, very thick] coordinates {
  (0.5,0.013) (0.65,0.0135) (0.78,0.0145) (0.83,0.017)
  (0.87,0.028) (0.90,0.045) (0.95,0.066) (1.0,0.074) (1.05,0.073)};
\addlegendentry{supercritical airfoil}
% M_DD markers
\draw[engblue, dashed] (axis cs:0.74,0) -- (axis cs:0.74,0.018);
\node[engblue, font=\scriptsize, anchor=south] at (axis cs:0.74,0.018) {$M_{DD}\approx 0.74$};
\draw[engred, dashed] (axis cs:0.84,0) -- (axis cs:0.84,0.020);
\node[engred, font=\scriptsize, anchor=south] at (axis cs:0.86,0.021) {$M_{DD}\approx 0.84$};
\end{axis}
\end{tikzpicture}
\caption[Transonic drag divergence]{Drag coefficient versus freestream
Mach number for a conventional and a supercritical transport airfoil
at fixed lift. The drag is nearly flat through the subsonic range,
then climbs steeply once local supersonic pockets on the upper surface
terminate in shocks strong enough to thicken and separate the boundary
layer. The Mach number where the climb begins is the drag-divergence
Mach number $M_{DD}$, conventionally taken where $dC_D/dM=0.1$. The
supercritical section (red), with its flatter upper surface and aft
camber, delays and weakens the shock and pushes $M_{DD}$ from about
$0.74$ to about $0.84$, the improvement that let transport aircraft
cruise efficiently near $M=0.80$ to $0.85$. The values are
representative of the published wind-tunnel families rather than a
specific section.}
\label{fig:vol03ch13:drag-divergence}
\end{figure}


Figure~\ref{fig:vol03ch13:airfoil-cp} shows the pressure coefficient
over a transonic section above its critical Mach number: a local
supersonic pocket on the upper surface, closed by a near-normal shock
that recompresses the flow and, if strong enough, separates the
boundary layer behind it. Figure~\ref{fig:vol03ch13:drag-divergence}
plots the drag rise this shock produces as $M_\infty$ increases, for a
conventional and a supercritical section; the dataset is in
\texttt{data/airfoil\_drag\_rise.csv}.

\engterm{Supersonic aerodynamics} ($1.2 < M_\infty < 5$). The
flow is everywhere supersonic at the freestream; shocks attach to
sharp leading edges. The \engterm{Ackeret thin-airfoil theory}
gives the pressure coefficient on a thin two-dimensional airfoil
at small angle of attack:
\[
C_p = \frac{2\theta}{\sqrt{M_\infty^2 - 1}},
\]
where $\theta$ is the local surface slope (positive for surfaces
inclined into the flow). The lift and wave drag coefficients of a
flat-plate airfoil at angle of attack $\alpha$ are
\[
C_L = \frac{4\alpha}{\sqrt{M_\infty^2 - 1}}, \quad
C_{D,\text{wave}} = \frac{4\alpha^2}{\sqrt{M_\infty^2 - 1}}.
\]
The wave drag is independent of viscosity; it arises from the
energy dissipated in the shocks attached to the airfoil's leading
edge. The lift-to-drag ratio of a flat plate in supersonic flow
is $L/D = 1/\alpha$, much lower than the typical $L/D \sim 15$ to
$20$ of subsonic transport aircraft.

\engterm{Slender-body theory} extends the supersonic analysis to
three-dimensional bodies of revolution. The pressure on a slender
body is given by an integral over the body's cross-sectional area
distribution; the wave drag depends on the integral of $(dA/dx)^2$,
giving the area-rule formula for wave drag (Sears, Haack). The
\engterm{Sears-Haack body} is the shape of revolution of given length
and volume that minimises this integral, a smooth spindle fattest near
mid-length and tapering to points at nose and tail; its wave drag is
$D_w = 9\pi^3 V^2/(8 \ell^4) \cdot \tfrac{1}{2}\rho_\infty U_\infty^2$
for volume $V$ and length $\ell$, the lower bound every practical
fuselage is measured against. The same integral is the bridge to the
transonic area rule: because the wave drag depends only on the
longitudinal distribution of cross-sectional area and not on how that
area is arranged in the cross-plane, a wing and a fuselage that share
an axial station each add their area there, so a smooth combined
$A(x)$ is what matters. This is why the project at the end of the
chapter builds the aircraft's total $A(x)$ curve before estimating any
drag: the wave drag reads off the second derivative of that single
curve.

\engterm{Hypersonic aerodynamics} ($M_\infty > 5$). Several new
physical mechanisms enter:

\begin{itemize}[leftmargin=*]
  \item \engterm{Real-gas effects}: the high temperatures behind
    strong shocks dissociate molecular bonds (oxygen $\to$ atomic
    oxygen at $T \sim 2000$--$4000\,\si{\kelvin}$, nitrogen at
    $T \sim 4000$--$9000\,\si{\kelvin}$), ionise the gas at higher
    temperatures, and produce radiation. The ideal-gas
    assumption fails; equilibrium and non-equilibrium chemistry
    models are needed.
  \item \engterm{Strong shock detachment}: the high deflection
    angles of practical hypersonic vehicles produce detached bow
    shocks at the nose, with the shock layer (the region between
    the shock and the body) becoming thin and highly nonuniform.
  \item \engterm{Viscous-inviscid coupling}: the boundary-layer
    thickness scales as $\delta/x \sim \text{Re}^{-1/2}\cdot M^2$
    in the hypersonic regime, so the boundary layer can dominate
    a large fraction of the shock layer. The inviscid shock-layer
    flow and the viscous boundary layer must be solved as a
    coupled system.
  \item \engterm{Aerodynamic heating}: the stagnation-point heat
    flux scales as $q_w \propto \rho_\infty^{1/2} U_\infty^3 /
    \sqrt{R_n}$, where $R_n$ is the nose radius. Hypersonic
    vehicles are blunt-nosed precisely to spread the heat load
    over a larger area; a sharp nose would heat to destruction.
    The thermal-protection-system design problem is the central
    challenge of re-entry vehicle engineering.
\end{itemize}

Hypersonic flow is the regime of ballistic missile re-entry, space
shuttle return, scramjet propulsion, and atmospheric re-entry of
deep-space probes. The Apollo command module re-entered Earth's
atmosphere at $M \approx 36$; the Space Shuttle re-entered at
$M \approx 25$. The design of the thermal-protection system
(ablative for Apollo, reusable ceramic tiles for the Shuttle) was
in each case the dominant vehicle-design problem
\cite{acc:caib-columbia}.

\section{Heating in high-speed flow}\pathtag{mastery}

The kinetic energy of high-speed flow converts to thermal energy
at any stagnation point. The stagnation temperature for an
adiabatic ideal-gas flow is
\[
T_0 = T\left(1 + \frac{\gamma - 1}{2}M^2\right).
\]
At sea-level ambient ($T = 288\,\si{\kelvin}$), $M = 2$ gives
$T_0 = 518\,\si{\kelvin}$ ($245\,\degC$); $M = 5$ gives
$T_0 = 1728\,\si{\kelvin}$ ($1455\,\degC$); $M = 10$ gives $T_0 =
5904\,\si{\kelvin}$ ($5631\,\degC$, exceeding the surface
temperature of the Sun). The aerodynamic-heating problem at high
Mach number is severe.

The actual surface temperature on a high-speed vehicle is set by
the balance between aerodynamic heating in and radiation,
conduction, and ablation out. The local heat flux at the
stagnation point is given to a good approximation by the Fay-Riddell
correlation; the working scaling (Anderson) is
\[
q_w \approx K\,\sqrt{\frac{\rho_\infty}{R_n}}\,U_\infty^3,
\]
with $K$ a property-dependent constant. The cube on velocity is
the strong scaling: doubling the velocity multiplies the
stagnation heat flux by a factor of eight. The inverse-square-root
on nose radius is the design lever: a blunt nose ($R_n$ large)
gives a lower stagnation heat flux than a sharp nose.

Thermal-protection systems are of three principal types:

\begin{itemize}[leftmargin=*]
  \item \engterm{Ablative}: a sacrificial material at the surface
    that chemically decomposes, carrying heat away with the
    departing mass. Apollo, Soyuz, Crew Dragon, ballistic re-entry
    vehicles. Single-use.
  \item \engterm{Radiative}: a high-emissivity ceramic surface
    that radiates heat back to the atmosphere or to space.
    Reusable, but limited to lower heat fluxes. Space Shuttle
    tiles, X-15 leading edges.
  \item \engterm{Active cooling}: pumped coolant through the
    structure, common in rocket engines but not (yet) in primary
    aerosurfaces of routine flight vehicles.
\end{itemize}

The Columbia accident of $1$ February $2003$
\cite{acc:caib-columbia} was caused by damage to the leading-edge
RCC panels during ascent, which exposed the wing structure to
hypersonic re-entry heating; the wing failed inside the high-
temperature shock layer. The accident is the canonical case of
thermal-protection-system failure in hypersonic flow; it is
treated in detail in the accident registry entry under
columbia-2003.

\begin{mastery}
\textbf{Prandtl-Meyer expansion.} A supersonic flow that turns away
from itself (around a convex corner) does not shock; it expands
isentropically through a fan of Mach waves, accelerating and dropping
in pressure. The turning is governed by the \engterm{Prandtl-Meyer
function}
\[
\nu(M) = \sqrt{\frac{\gamma + 1}{\gamma - 1}}\,\tan^{-1}\!\sqrt{
\frac{\gamma - 1}{\gamma + 1}(M^2 - 1)} - \tan^{-1}\!\sqrt{M^2 - 1},
\]
the angle through which a sonic flow ($M = 1$, $\nu = 0$) must turn to
reach Mach $M$. A flow at $M_1$ turning through a convex angle
$\Delta\theta$ reaches $M_2$ from $\nu(M_2) = \nu(M_1) + \Delta\theta$.
The expansion is isentropic, so the stagnation pressure is conserved
across the fan, unlike a shock. The pairing of oblique-shock
compression on windward surfaces with Prandtl-Meyer expansion on
leeward surfaces is the \engterm{shock-expansion method}, the exact
inviscid solution for a thin supersonic airfoil with straight panels.
The solver in \texttt{code/prandtl\_meyer.py} evaluates $\nu(M)$ and
inverts it; for $\gamma = 1.4$, $\nu(2) \approx 26.4^{\circ}$, and a
flow at $M = 2$ turning through $10^{\circ}$ reaches $M \approx 2.38$.
As $M \to \infty$, $\nu \to \nu_{\max} = (\pi/2)
(\sqrt{(\gamma+1)/(\gamma-1)} - 1) \approx 130.5^{\circ}$: a sonic flow
cannot turn through more than about $130.5^{\circ}$ before reaching
infinite Mach number and zero pressure, the theoretical limit of
expansion into vacuum.
\end{mastery}

\begin{mastery}
\textbf{The Rankine-Hugoniot curve and why shocks compress.} The
normal-shock relations of section 13.4 are usually written as ratios in
terms of $M_1$, but they carry a deeper statement when plotted in the
thermodynamic plane of pressure against specific volume. Eliminating
the velocities between the mass, momentum, and energy conservation laws
across the shock gives a single relation between the downstream and
upstream states that contains no velocity at all, the
\engterm{Rankine-Hugoniot relation}
\[
\frac{p_2}{p_1} = \frac{(\gamma + 1) - (\gamma - 1)(v_2/v_1)}
{(\gamma + 1)(v_2/v_1) - (\gamma - 1)},
\]
where $v = 1/\rho$ is the specific volume. The locus of downstream
states this relation allows, drawn through a fixed upstream state, is
the \engterm{Hugoniot curve}. Figure~\ref{fig:vol03ch13:hugoniot} plots
it against the \engterm{isentrope} $p_2/p_1 = (v_2/v_1)^{-\gamma}$ drawn
through the same point. Both curves pass through the upstream state and
both rise under compression ($v_2/v_1 < 1$), but on the compression
branch the Hugoniot lies above the isentrope: a shock reaches a higher
pressure for the same volume reduction than an isentropic compression
would, because the shock raises the entropy and the extra pressure is
the thermal signature of that entropy rise. The gap between the two
curves is the irreversibility of the shock, and it is the same
quantity the derivation exercise expressed as the cubic vanishing of
$\Delta s$ near
$M_1 = 1$: close to the upstream point the Hugoniot and the isentrope
are tangent to second order, so weak shocks are nearly isentropic, and
the curves separate only as the cube of the shock strength.

The Hugoniot has a vertical asymptote at
$v_2/v_1 = (\gamma - 1)/(\gamma + 1) = 1/6$ for air: the pressure runs
to infinity as the volume ratio approaches this limit, the
$p$-$v$-plane statement of the saturation of the density ratio at
$(\gamma + 1)/(\gamma - 1) = 6$ seen in
figure~\ref{fig:vol03ch13:normal-shock}. No perfect-gas shock,
however strong, compresses the gas past six times its upstream density;
the energy that a stronger shock supplies goes into temperature, not
further compression. The \engterm{Rayleigh line}, the straight chord
from the upstream state to the downstream state, is the locus that
satisfies mass and momentum conservation alone; its intersection with
the Hugoniot, which adds energy conservation, fixes the downstream
state, and the slope of the chord sets the shock speed (steeper chord,
faster shock). The decisive feature is the expansion branch
($v_2/v_1 > 1$): there the Hugoniot would require the entropy to fall,
which the second law forbids, so an expansion shock cannot exist. A
supersonic flow that needs to expand does it the only way entropy
allows, smoothly through a Prandtl-Meyer fan rather than abruptly
through a discontinuity. Shocks are compression-only, and the Hugoniot
curve is where that one-sidedness is written.
\end{mastery}

% Figure: the Rankine-Hugoniot curve versus the isentrope through the
% same initial state, in the pressure-specific-volume plane. The
% Hugoniot lies above the isentrope on the compression branch (a shock
% reaches a higher pressure for the same compression because it raises
% entropy). The Rayleigh line connects initial and final states.
% Use v = specific volume (normalised v1=1, p1=1). gamma=1.4.
% Hugoniot: p2/p1 = ( (g+1) - (g-1)(v2/v1) ) / ( (g+1)(v2/v1) - (g-1) )
%   g=1.4 -> num = 2.4 - 0.4*v ; den = 2.4*v - 0.4   (v = v2/v1 in [0.2,1])
% Isentrope: p2/p1 = (v2/v1)^(-g) = v^(-1.4)
\begin{figure}[htbp]
\centering
\begin{tikzpicture}
\begin{axis}[
  width=11cm, height=8cm,
  xlabel={specific volume ratio $v_2/v_1$},
  ylabel={pressure ratio $p_2/p_1$},
  xmin=0.15, xmax=1.6, ymin=0, ymax=9,
  axis lines=left,
  grid=major, grid style={enggray!20},
  tick label style={font=\scriptsize},
  label style={font=\small},
  legend style={font=\scriptsize, at={(0.97,0.97)}, anchor=north east},
  legend cell align=left,
]
% Hugoniot curve (compression branch v<1 and a little of v>1)
\addplot[engred, very thick, domain=0.18:1.55, samples=160]
  { (2.4 - 0.4*x) / (2.4*x - 0.4) };
\addlegendentry{Hugoniot (shock)}
% Isentrope through (1,1)
\addplot[engblue, very thick, domain=0.18:1.55, samples=160]
  { x^(-1.4) };
\addlegendentry{isentrope}
% Rayleigh line from initial state (1,1) to a final state at v=0.4
% p at v=0.4 on Hugoniot: (2.4-0.16)/(0.96-0.4)=2.24/0.56=4.0
\addplot[englime, thick, dashed] coordinates {(1,1) (0.4,4.0)};
\addlegendentry{Rayleigh line}
% initial state marker
\fill[engblack] (axis cs:1,1) circle (1.6pt);
\node[engblack, font=\scriptsize, anchor=north west] at (axis cs:1.02,1.0)
  {$(1,1)$ upstream};
% final state marker
\fill[engblack] (axis cs:0.4,4.0) circle (1.6pt);
\node[engblack, font=\scriptsize, anchor=south west] at (axis cs:0.42,4.0)
  {downstream};
% vertical asymptote of Hugoniot at v=(g-1)/(g+1)=0.1667
\draw[enggray, dotted] (axis cs:0.1667,0) -- (axis cs:0.1667,9);
\node[enggray, font=\scriptsize, anchor=south west, rotate=90]
  at (axis cs:0.1667,3.0) {$v_2/v_1\to\frac{\gamma-1}{\gamma+1}$};
\end{axis}
\end{tikzpicture}
\caption[Hugoniot curve versus the isentrope]{The Rankine-Hugoniot
curve (red) and the isentrope (blue) drawn through the same upstream
state $(v_2/v_1,p_2/p_1)=(1,1)$, in the pressure-specific-volume plane
for air ($\gamma=1.4$). Both pass through the upstream point and both
rise as the gas is compressed ($v_2/v_1<1$), but on the compression
branch the Hugoniot lies above the isentrope: a shock reaches a higher
pressure for the same volume compression because it raises the entropy,
turning ordered kinetic energy into heat. The gap between the two curves
is the irreversibility of the shock. The Hugoniot has a vertical
asymptote at $v_2/v_1=(\gamma-1)/(\gamma+1)=1/6$, the limiting density
ratio of a perfect-gas normal shock: no shock, however strong,
compresses the gas past six times its upstream density. The Rayleigh
line (green dashed) is the locus of states that conserve mass and
momentum across the shock; its intersection with the Hugoniot (which
adds energy conservation) fixes the downstream state, and its slope sets
the shock speed. The expansion branch ($v_2/v_1>1$) of the Hugoniot
would lower the entropy, so the second law forbids expansion shocks: a
supersonic flow expands isentropically through a Prandtl-Meyer fan, not
through a discontinuity.}
\label{fig:vol03ch13:hugoniot}
\end{figure}


\begin{mastery}
\textbf{Crocco's theorem and the entropy layer.} The shocks of this
chapter raise entropy, and a curved shock raises it by different
amounts along its length, because the local upstream normal Mach number
$M_1\sin\beta$ varies with the shock angle $\beta$ from point to point.
A blunt body's bow shock is strongest where it is normal to the flow,
on the stagnation streamline, and weakens toward the flanks where it
lies more obliquely, so the gas that crosses near the nose carries more
entropy than the gas that crosses out on the wings. The downstream flow
therefore has a transverse entropy gradient, and this gradient is not
inert: it drives vorticity. \engterm{Crocco's theorem} makes the link
exact. For steady adiabatic flow it states
\[
T\,\nabla s = \nabla h_0 - \mathbf{u}\times\boldsymbol{\omega},
\]
where $s$ is the specific entropy, $h_0$ the stagnation enthalpy, and
$\boldsymbol{\omega} = \nabla\times\mathbf{u}$ the vorticity. When the
stagnation enthalpy is uniform (as it is behind an adiabatic shock,
$h_0$ being conserved across it), the theorem reduces to $T\,\nabla s =
-\mathbf{u}\times\boldsymbol{\omega}$: an entropy gradient transverse to
the flow forces a vorticity, and a flow with no entropy gradient and
uniform $h_0$ is irrotational. The transverse entropy gradient behind a
curved shock thus generates a rotational \engterm{entropy layer} that
clings to the body, thickest near the nose, where it merges with the
viscous boundary layer and complicates the heat-transfer prediction the
re-entry case study depends on. Crocco's theorem is also the clean
reason that the flow ahead of any shock, where the entropy is uniform
and $h_0$ is uniform, is irrotational and admits a velocity potential,
which is what licenses the linearised potential theories (Prandtl-
Glauert, Ackeret) used earlier in the chapter; those theories hold only
where no shock has yet stamped an entropy gradient into the flow.
\end{mastery}

\begin{mastery}
\textbf{Method of characteristics.} The steady two-dimensional
supersonic flow equations are hyperbolic, so disturbances propagate
along characteristic curves, the left- and right-running Mach lines
inclined at $\pm\mu$ to the local flow. Along each characteristic a
\engterm{Riemann invariant} is constant: $\theta + \nu(M)$ along one
family and $\theta - \nu(M)$ along the other, with $\theta$ the local
flow angle and $\nu$ the Prandtl-Meyer function. The method of
characteristics integrates these invariants across a grid of
intersecting Mach lines to build the flow field point by point, and it
is the classical way to design the contour of a supersonic
wind-tunnel or rocket nozzle that delivers shock-free, uniform,
parallel exit flow. The procedure is to fix the desired exit Mach
number, work the expansion from the throat through a sequence of
characteristic intersections, and cancel each reflected wave at the
wall by turning the wall through the matching Prandtl-Meyer angle. The
result is the bell contour of a modern rocket nozzle, which a conical
nozzle only approximates. The method predates the digital computer; it
was the working nozzle-design tool of the 1940s and 1950s and remains
the reference against which computational solvers are checked.
\end{mastery}

\begin{mastery}
\textbf{Shock-expansion theory and hypersonic similarity.} For a thin
supersonic airfoil built from straight panels, the exact inviscid
surface pressure follows from walking the flow panel by panel: each
concave turn into the flow runs through an oblique shock, each convex
turn away from the flow through a Prandtl-Meyer expansion, and the
pressure on each panel is whatever the running tally of compressions
and expansions delivers. This is the \engterm{shock-expansion method},
and unlike the linearised Ackeret result $C_p = 2\theta/\sqrt{M_\infty^2
- 1}$ it carries the full nonlinearity, so it stays accurate at the
larger deflections and higher Mach numbers where the linear theory
drifts. For a symmetric diamond (double-wedge) section at angle of
attack, the method gives four panel pressures, and their resultant is
the exact inviscid lift and wave drag; the linear theory is recovered
when every turning angle is small. The arithmetic reuses
\texttt{code/oblique\_shock.py} for the compression faces and
\texttt{code/prandtl\_meyer.py} for the expansion faces. At the upper
end of the Mach range a different reduction takes over.
\engterm{Hypersonic similarity} states that for slender bodies at high
Mach number the flow depends on $M_\infty$ and the thickness ratio
$\tau$ only through the \engterm{hypersonic similarity parameter}
$K = M_\infty \tau$. Two bodies with the same $K$ share the same
pressure coefficient once it is scaled by $\tau^2$, so a single
wind-tunnel or computed result covers a whole family of (Mach number,
thickness) pairs. In the limit $M_\infty \to \infty$ at fixed $K$ the
steady hypersonic small-disturbance equations become identical to the
unsteady transonic equations in one fewer space dimension, the
\engterm{equivalence principle} that lets the slender-body hypersonic
problem borrow the machinery of unsteady transonic flow. Newtonian
impact theory, $C_p = 2\sin^2\theta$, is the crude zeroth-order
companion: it treats the gas as a stream of particles that lose their
normal momentum on striking the surface, and it is serviceable for the
windward faces of blunt hypersonic bodies, where the shock layer is
thin and hugs the body.
\end{mastery}

\begin{mastery}
\textbf{Compressibility corrections and the critical Mach number.} The
incompressible pressure coefficient $C_{p,0}$ on a subsonic body is
amplified by compressibility, and a sequence of corrections of growing
accuracy tracks the effect as the freestream Mach number climbs toward
unity. The \engterm{Prandtl-Glauert rule} is the leading-order result,
$C_p = C_{p,0}/\sqrt{1 - M_\infty^2}$, derived by linearising the
subsonic small-disturbance equation; it is accurate to about
$M_\infty = 0.7$ and underpredicts the suction peak above that. The
\engterm{Karman-Tsien rule},
\[
C_p = \frac{C_{p,0}}{\sqrt{1 - M_\infty^2} + \dfrac{M_\infty^2}{1 +
\sqrt{1 - M_\infty^2}}\dfrac{C_{p,0}}{2}},
\]
and the \engterm{Laitone rule} carry a nonlinear dependence on
$C_{p,0}$ itself and hold closer to the high-subsonic edge, because they
retain a term the Prandtl-Glauert linearisation drops. All three blow up
as $M_\infty \to 1$, the signature that the linearised subsonic theory
cannot reach the transonic regime where local pockets go supersonic.
The practical use of the corrections is to find the \engterm{critical
Mach number} $M_{\text{cr}}$, the freestream Mach number at which the
flow first reaches $M = 1$ somewhere on the body, usually at the
minimum-pressure (suction) peak. The condition is geometric: the
compressibility-corrected suction-peak $C_p$ equals the
\engterm{critical pressure coefficient}
\[
C_p^* = \frac{2}{\gamma M_\infty^2}\left[\left(\frac{1 +
\frac{\gamma - 1}{2}M_\infty^2}{1 + \frac{\gamma - 1}{2}}\right)^{
\gamma/(\gamma - 1)} - 1\right],
\]
the pressure coefficient at which the local Mach number is exactly $1$.
Plotting the corrected suction-peak $C_p(M_\infty)$ and the curve
$C_p^*(M_\infty)$ on the same axes, their intersection is $M_{\text{cr}}$.
A thin or low-camber section has a shallow suction peak, so its
corrected $C_p$ meets $C_p^*$ at a higher $M_\infty$, giving a higher
critical Mach number and a higher drag-divergence Mach number; the
supercritical airfoil is the deliberate flattening of the upper-surface
pressure that pushes $M_{\text{cr}}$, and with it the usable cruise
Mach, as far toward unity as the section will bear. Above $M_{\text{cr}}$
no algebraic correction survives, and the analysis passes to the
transonic small-disturbance equation or to a full numerical solution.
\end{mastery}

\begin{mastery}
\textbf{The Busemann biplane and wave-drag cancellation.} Ackeret theory
puts an irreducible wave drag on any lifting supersonic surface, but the
\emph{volume} wave drag of a non-lifting section carries no such
theorem, and Busemann showed in 1935 that it can in principle be
cancelled. Consider two thin wedge profiles stacked to form a channel
between them, each the mirror of the other, so the pair looks like a
hollow diamond seen edge on. The compression shock from the leading edge
of one plate strikes the leading edge of the other and reflects, and if
the spacing is tuned to the design Mach number the reflected wave lands
exactly where the second plate's own expansion cancels it. Between the
plates the pressures on the facing surfaces do work that cancels in the
drag direction, and outside the pair the flow is left undisturbed, so
the system radiates no wave energy and its volume wave drag vanishes at
the design point. The catch is threefold: the cancellation holds only at
one Mach number, it collapses at any angle of attack because lift
unbalances the internal pressures, and the two plates enclose no usable
volume for a wing that must also carry load. The Busemann biplane is
therefore a theorem about what wave drag \emph{is}, an interference
between the body's own waves, rather than a practical wing, and its
modern descendants (the shock-cancelling inlet, the swept
supercritical-section shaping, the low-boom fuselage) all trade on the
same idea that a body's waves can be made to interfere destructively.
The lesson the Ackeret result alone hides is that wave drag is not a
fixed toll on volume but a property of the whole wave field, open to
cancellation by geometry the same way two sound sources can be placed to
cancel.
\end{mastery}

\begin{mastery}
\textbf{Regular versus Mach reflection and the von Neumann paradox.} When
an oblique shock strikes a wall (figure~\ref{fig:vol03ch13:shock-reflection}),
the wall demands that the flow leave parallel to it, and a single
reflected shock supplies the second turn that restores parallelism. This
\engterm{regular reflection} works only while the post-incident Mach
number $M_2$ can still be turned through the deflection angle by some
attached shock. As the incident shock strengthens or its angle steepens,
$M_2$ falls, the required turn approaches the local $\theta_{\max}(M_2)$,
and at the \engterm{detachment condition} no reflected shock can hold the
flow to the wall. What forms instead is \engterm{Mach reflection}: the
incident and reflected shocks meet away from the wall at a triple point,
and a third, nearly normal shock (the \engterm{Mach stem}) bridges from
the triple point down to the surface, with a slip line trailing
downstream across which the two processed streams share pressure and
direction but differ in entropy and velocity. The transition between the
two patterns is read on the shock polar of section~13.4: regular
reflection exists while the reflected-shock polar of the region-2 stream
still crosses the pressure axis, and Mach reflection takes over when the
polars fail to cross. Near the transition the steady and unsteady
criteria disagree, and in a band of weak shocks both patterns appear
admissible while experiments show a third configuration the two-shock
theory cannot produce, the \engterm{von Neumann paradox}, still an open
problem in shock dynamics. The engineering consequence is direct: a
shock reflecting inside a supersonic inlet or between the panels of a
mixed-compression duct can switch from a clean regular reflection to a
Mach stem with a strong normal-shock loss and a slip line that seeds
separation, so the reflection pattern, not just the shock strength, is
part of the inlet's design margin.
\end{mastery}

\begin{mastery}
\textbf{Taylor-Maccoll flow over a cone.} The oblique-shock relations of
section~13.4 solve the planar wedge, where the flow behind the shock is
uniform and parallel to the wall. A circular cone in supersonic flow is
not planar, and its solution carries a feature the wedge does not: the
flow behind a conical shock is not uniform. A cone of half-angle
$\theta_c$ at zero incidence generates a straight conical shock from its
tip, but between the shock and the cone surface the streamlines curve and
the properties vary with the polar angle measured from the axis, because
the flow is axisymmetric and the area available to each stream tube grows
with distance from the tip. The governing statement is the
\engterm{Taylor-Maccoll equation}, an ordinary differential equation for
the radial velocity component $V_r(\omega)$ as a function of the polar
angle $\omega$, obtained by writing the steady, irrotational,
axisymmetric continuity and momentum equations in spherical coordinates
and using the fact that, along a ray from the tip, the flow depends only
on $\omega$:
\[
\frac{\gamma - 1}{2}\left[V_0^2 - V_r^2 -
\left(\frac{dV_r}{d\omega}\right)^2\right]\!\left[2V_r +
\frac{dV_r}{d\omega}\cot\omega + \frac{d^2 V_r}{d\omega^2}\right]
= \frac{dV_r}{d\omega}\!\left[V_r\frac{dV_r}{d\omega} +
\frac{dV_r}{d\omega}\frac{d^2 V_r}{d\omega^2}\right],
\]
with $V_0$ the maximum (vacuum) speed. The equation is integrated
numerically inward from the shock, where the oblique-shock relations set
the starting values, to the cone surface, where the flow-tangency
condition $dV_r/d\omega = 0$ (no velocity normal to the cone) picks out
the cone angle that the chosen shock angle serves. The practical
consequences separate the cone from the wedge cleanly. For a given
freestream Mach number and turning angle, the conical shock is weaker and
lies at a smaller angle than the planar shock, because the axisymmetric
relief lets the flow spread; a cone therefore compresses more gently than
a wedge of the same half-angle, and it stays shock-attached to a larger
half-angle before detaching. The surface pressure on a cone is lower than
on the equivalent wedge, and the flow reaches the surface having been
compressed partly through the shock and partly through the isentropic
convergence behind it. Taylor and Maccoll's 1933 solution is the
reference case for supersonic conical bodies, the analytic backbone of
the classic cone tables, and the reason a conical forebody or a conical
inlet centre-body is analysed by its own relations rather than by the
planar oblique-shock chart.
\end{mastery}

\section{Worked examples}\pathtag{standard}

The following worked examples run the chapter's relations end to end.
Each states the data, the relation, and the arithmetic, and each is
reproducible with the companion code.

\paragraph{Worked example 1: stagnation conditions on a cruising jet.}
A transport cruises at $M = 0.85$ at $11\,\si{\kilo\meter}$, where
$T_\infty = 217\,\si{\kelvin}$ and $p_\infty = 22.7\,\text{kPa}$
(from \texttt{data/standard\_atmosphere.csv}). The stagnation
temperature and pressure at the nose, from the isentropic relations
with $\gamma = 1.4$, are
\[
\frac{T_0}{T} = 1 + 0.2\,(0.85)^2 = 1.1445, \quad
T_0 = 217\times 1.1445 = 248\,\si{\kelvin},
\]
\[
\frac{p_0}{p} = 1.1445^{3.5} = 1.604, \quad
p_0 = 22.7\times 1.604 = 36.4\,\text{kPa}.
\]
The stagnation-temperature rise is only $31\,\si{\kelvin}$, modest at
this Mach number; the stagnation pressure is $60\,\%$ above ambient,
which is why a pitot-static airspeed system must use the compressible
relation, not the incompressible $\tfrac{1}{2}\rho u^2$, above
$M \approx 0.3$.

\paragraph{Worked example 2: design and off-design of a supersonic
nozzle.} A converging-diverging nozzle is to deliver $M_e = 2.5$ from
$p_0 = 1\,\text{MPa}$, $T_0 = 500\,\si{\kelvin}$ (air). The required
area ratio from the area-Mach relation (\texttt{code/nozzle.py}) is
$A_e/A^* = 2.637$. At the design back pressure the exit static
pressure is $p_e = p_0 (p/p_0)_{M=2.5} = 1\,\text{MPa}\times 0.0585 =
58.5\,\text{kPa}$. If instead the back pressure is raised to
$0.6\,\text{MPa}$, the nozzle cannot run fully supersonic: a normal
shock stands in the diverging section. The shock sits where the
post-shock subsonic diffusion just recovers to $0.6\,\text{MPa}$ at
the exit; locating it requires matching the post-shock stagnation
pressure ($p_{02} = p_{01}\,(p_{02}/p_{01})$ from the shock at the
local supersonic Mach number) to the exit condition. This iteration is
the off-design analysis that \texttt{code/normal\_shock.py} and
\texttt{code/nozzle.py} support jointly.

\paragraph{Worked example 3: oblique shock on a compression ramp.}
A supersonic inlet ramp turns a $M_1 = 3$ flow through
$\theta = 15^{\circ}$. From the $\theta$-$\beta$-$M$ relation the weak
shock angle is $\beta = 32.2^{\circ}$
(\texttt{code/oblique\_shock.py}). The upstream normal Mach component
is $M_{n1} = M_1 \sin\beta = 3\times \sin 32.2^{\circ} = 1.60$. The
normal-shock relations at $M_{n1} = 1.60$ give $M_{n2} = 0.668$ and
$p_2/p_1 = 2.82$. The downstream Mach number along the deflected flow
is $M_2 = M_{n2}/\sin(\beta - \theta) = 0.668/\sin 17.2^{\circ} =
2.26$, still supersonic, so a second ramp can take a further turn.
Stacking weak oblique shocks this way (the principle of the external
compression inlet) recovers far more stagnation pressure than a single
normal shock at $M_1 = 3$, which would lose two thirds of it
($p_{02}/p_{01} = 0.328$).

\paragraph{Worked example 4: flat-plate lift and drag in supersonic
flow.} A flat-plate airfoil flies at $M_\infty = 2$ and
$\alpha = 5^{\circ} = 0.0873\,\text{rad}$. Ackeret theory gives
$\beta = \sqrt{M_\infty^2 - 1} = \sqrt{3} = 1.732$, so
\[
C_L = \frac{4\alpha}{\beta} = \frac{4\times 0.0873}{1.732} = 0.202,
\quad
C_{D,\text{wave}} = \frac{4\alpha^2}{\beta} = \frac{4\times 0.0873^2}{
1.732} = 0.0176,
\]
and $L/D = 1/\alpha = 11.5$. A subsonic transport reaches $L/D \approx
17$ to $20$; the supersonic flat plate pays a wave-drag penalty that
caps its cruise efficiency, the reason Concorde burned roughly twice
the fuel per passenger-kilometre of a subsonic widebody.

% Figure: Prandtl-Meyer function nu(M) in degrees versus Mach number,
% gamma=1.4. nu is the angle through which a sonic flow must turn to
% reach Mach M by isentropic expansion. Saturates at nu_max ~ 130.5 deg
% as M -> infinity.
\begin{figure}[htbp]
\centering
\begin{tikzpicture}
\begin{axis}[
  width=12cm, height=8cm,
  xlabel={Mach number $M$}, ylabel={Prandtl-Meyer angle $\nu(M)$ (deg)},
  xmin=1, xmax=10, ymin=0, ymax=135,
  axis lines=left,
  grid=major, grid style={enggray!20},
  tick label style={font=\scriptsize},
  label style={font=\small},
  legend style={font=\scriptsize, at={(0.97,0.03)}, anchor=south east},
  legend cell align=left,
  domain=1:10, samples=300,
]
% nu(M) = sqrt(6) * atan( sqrt( (M^2-1)/6 ) ) - atan( sqrt(M^2-1) )
% with sqrt((g+1)/(g-1)) = sqrt(2.4/0.4) = sqrt(6) for gamma=1.4.
% pgf atan returns degrees; sqrt(6)=2.449 multiplies the first term,
% which is already in degrees. The second atan is also in degrees.
\addplot[engblue, very thick] {
  2.449*atan( sqrt( (x^2 - 1)/6 ) ) - atan( sqrt(x^2 - 1) ) };
\addlegendentry{$\nu(M)$, $\gamma=1.4$}
% asymptote nu_max = 90*(sqrt(6)-1) = 130.45 deg
\addplot[engred, dashed, domain=1:10, samples=2] {130.45};
\addlegendentry{$\nu_{\max}=130.5^{\circ}$}
% reference point M=2, nu~26.4
\node[engblack, font=\scriptsize, anchor=west] at (axis cs:2.1,26.4)
  {$\nu(2)\approx26.4^{\circ}$};
\addplot[only marks, mark=*, engblack, mark size=1.4pt] coordinates {(2,26.38)};
\end{axis}
\end{tikzpicture}
\caption[The Prandtl-Meyer function]{The Prandtl-Meyer function
$\nu(M)$ for $\gamma=1.4$, the angle through which a sonic stream
($M=1$, $\nu=0$) must turn around a convex corner to reach Mach $M$ by
isentropic expansion. A flow already at $M_1$ that turns through a
further convex angle $\Delta\theta$ reaches $M_2$ from $\nu(M_2) =
\nu(M_1) + \Delta\theta$, so the curve is read as a turning budget
rather than a function of position. The expansion is isentropic, so the
stagnation pressure is conserved across the fan, the property that
distinguishes an expansion corner from a compression corner (which
shocks). The angle saturates at $\nu_{\max} = 90^{\circ}(\sqrt{6}-1)
\approx 130.5^{\circ}$ as $M\to\infty$: a sonic flow cannot turn
through more than this before reaching infinite Mach number and zero
pressure, the limit of expansion into vacuum.}
\label{fig:vol03ch13:prandtl-meyer}
\end{figure}


% Figure: schematic Fanno line (adiabatic flow with friction) and
% Rayleigh line (frictionless flow with heat addition) on a
% temperature-entropy diagram. Both lines have a sonic point (M=1) at
% their right extremity; the upper branch is subsonic, the lower
% supersonic. A normal shock jumps from the supersonic point on one
% curve to the subsonic point at the same mass flux and momentum.
\begin{figure}[htbp]
\centering
\begin{tikzpicture}
\begin{axis}[
  width=12cm, height=8cm,
  xlabel={entropy $s$ (arbitrary units)}, ylabel={temperature $T$ (arbitrary units)},
  xmin=0, xmax=10, ymin=0, ymax=10,
  axis lines=left,
  grid=major, grid style={enggray!20},
  tick label style={font=\scriptsize},
  label style={font=\small},
  legend style={font=\scriptsize, at={(0.03,0.97)}, anchor=north west},
  legend cell align=left,
]
% Fanno line: a curve that rises, turns at a sonic point (max entropy),
% and falls. Parametrise with a simple lobe. Use plot coordinates.
\addplot[engblue, very thick, smooth] coordinates {
  (1.0,8.6) (2.2,8.3) (3.6,7.8) (5.2,7.0) (6.6,6.0)
  (7.8,4.6) (8.6,3.0) (8.95,1.9) (8.6,1.1) (7.9,0.6)};
\addlegendentry{Fanno line (friction, adiabatic)}
% Rayleigh line: a fuller loop; max entropy at sonic, max temperature
% at a subsonic point left of sonic.
\addplot[engred, very thick, smooth] coordinates {
  (1.0,5.0) (2.6,6.2) (4.2,7.1) (5.8,7.5) (7.0,7.3)
  (8.0,6.5) (8.7,5.2) (9.0,3.6) (8.7,2.2) (8.0,1.2)};
\addlegendentry{Rayleigh line (heat addition, frictionless)}
% sonic point markers near the right tips (max entropy)
\addplot[only marks, mark=*, engblack, mark size=1.6pt] coordinates {(8.95,1.9)};
\addplot[only marks, mark=square*, engblack, mark size=1.6pt] coordinates {(9.0,3.6)};
\node[engblack, font=\scriptsize, anchor=west] at (axis cs:9.05,1.9) {$M=1$};
\node[engblack, font=\scriptsize, anchor=west] at (axis cs:9.1,3.6) {$M=1$};
% branch labels
\node[engblue, font=\scriptsize] at (axis cs:3.0,8.7) {subsonic};
\node[engblue, font=\scriptsize] at (axis cs:7.6,0.6) {supersonic};
\end{axis}
\end{tikzpicture}
\caption[Fanno and Rayleigh lines on a $T$-$s$ diagram]{Schematic Fanno
and Rayleigh lines on a temperature-entropy diagram, the two
one-dimensional flows that drive a constant-area duct toward $M=1$ from
either side. The Fanno line traces adiabatic flow with wall friction at
fixed mass flux: friction always raises entropy, so a subsonic flow on
the upper branch accelerates toward the sonic point at the right tip,
and a supersonic flow on the lower branch decelerates toward the same
point; the duct chokes when the flow reaches $M=1$ and added length
cannot pass more flow. The Rayleigh line traces frictionless flow with
heat addition at fixed mass flux and momentum: heating likewise drives
both branches toward the sonic point, so a combustor can thermally
choke. Each curve has its maximum entropy at $M=1$, which is why a
flow cannot cross from subsonic to supersonic by friction or heating
alone. The construction is the working tool for duct flow with friction
(Fanno) and with heat release (Rayleigh), the topics of the worked
examples below.}
\label{fig:vol03ch13:fanno-rayleigh}
\end{figure}


\paragraph{Worked example 5: normal shock in a constant-area duct
(supersonic wind-tunnel diffuser).} A supersonic wind tunnel runs
$M_1 = 2.5$ air in its test section at $p_1 = 30\,\text{kPa}$,
$T_1 = 160\,\si{\kelvin}$. A normal shock stands at the entrance to the
diffuser. We want the static state behind the shock, the velocity
drop, and the stagnation-pressure loss the tunnel's compressor must
make up each pass. The upstream speed of sound is $a_1 =
\sqrt{1.4\times 287\times 160} = 253\,\si{\meter\per\second}$, so
$u_1 = M_1 a_1 = 2.5\times 253 = 633\,\si{\meter\per\second}$. From the
normal-shock relations at $M_1 = 2.5$
(\texttt{code/normal\_shock.py}): $M_2 = 0.513$, $p_2/p_1 = 7.13$,
$\rho_2/\rho_1 = 3.33$, $T_2/T_1 = 2.14$. Hence $p_2 = 214\,\text{kPa}$,
$T_2 = 342\,\si{\kelvin}$, $a_2 = \sqrt{1.4\times 287\times 342} =
371\,\si{\meter\per\second}$, and $u_2 = M_2 a_2 = 190\,\si{\meter\per
\second}$. The velocity drops by a factor of $3.33$, exactly the
density ratio, as continuity $\rho_1 u_1 = \rho_2 u_2$ demands. The
stagnation-pressure ratio across the shock is
$p_{02}/p_{01} = 0.499$: the shock throws away half the stagnation
pressure, a loss the tunnel's compressor restores on every circuit.
This is the penalty a single terminal normal shock imposes, and the
reason a supersonic diffuser is designed to set the shock at the
lowest Mach number the geometry allows, where the loss is least.

\paragraph{Worked example 6: Prandtl-Meyer expansion over a corner.}
A supersonic flow at $M_1 = 2.0$, $p_1 = 100\,\text{kPa}$ turns around
a convex corner through $\Delta\theta = 15^{\circ}$
(figure~\ref{fig:vol03ch13:prandtl-meyer}). We want the downstream
Mach number and pressure. The expansion is isentropic, so the
stagnation pressure is conserved and the turning budget closes through
the Prandtl-Meyer function. From \texttt{code/prandtl\_meyer.py},
$\nu(2.0) = 26.38^{\circ}$. Adding the turn,
$\nu(M_2) = 26.38^{\circ} + 15^{\circ} = 41.38^{\circ}$. Inverting the
Prandtl-Meyer function gives $M_2 = 2.60$. The pressure drop follows
from the isentropic relation referenced to the common stagnation
pressure:
\[
\frac{p_2}{p_1} = \frac{p_2/p_0}{p_1/p_0}
= \frac{(1 + 0.2\,M_2^2)^{-3.5}}{(1 + 0.2\,M_1^2)^{-3.5}}
= \frac{(1 + 0.2\times 6.76)^{-3.5}}{(1 + 0.2\times 4)^{-3.5}}
= \frac{0.0501}{0.1278} = 0.392,
\]
so $p_2 = 39.2\,\text{kPa}$. The flow accelerates from $M = 2.0$ to
$M = 2.6$ and its pressure drops to $39\,\%$ of the upstream value, all
without any loss of stagnation pressure. Compare worked example~3,
where a $15^{\circ}$ \emph{compression} turn at the same upstream Mach
number raises the pressure through a shock that does lose stagnation
pressure: expansion and compression corners are not mirror images once
entropy is tracked.

\paragraph{Worked example 7: Fanno flow toward choking in a long pipe.}
Air enters a smooth insulated pipe of diameter $D = 50\,\si{\milli
\meter}$ at $M_1 = 0.30$ with friction factor $f = 0.020$
(figure~\ref{fig:vol03ch13:fanno-rayleigh}). Adiabatic flow with wall
friction drives the Mach number toward unity; we want the pipe length
that just chokes the flow. Fanno theory tabulates the dimensionless
length to choking,
\[
\frac{4 f L^*}{D} = \frac{1 - M^2}{\gamma M^2}
+ \frac{\gamma + 1}{2\gamma}\ln\!\left[\frac{(\gamma + 1)M^2}
{2 + (\gamma - 1)M^2}\right].
\]
At $M_1 = 0.30$, $\gamma = 1.4$: the first term is
$(1 - 0.09)/(1.4\times 0.09) = 7.22$, and the logarithmic term is
$\frac{2.4}{2.8}\ln[(2.4\times 0.09)/(2 + 0.4\times 0.09)] =
0.857\ln(0.216/2.036) = 0.857\times(-2.243) = -1.92$, so
$4fL^*/D = 5.30$. Solving for the choking length,
$L^* = 5.30\,D/(4f) = 5.30\times 0.05/(4\times 0.020) =
3.31\,\si{\meter}$. A pipe longer than $3.3\,\si{\meter}$ cannot pass
the entry flow at $M_1 = 0.30$; the upstream conditions adjust (the
mass flow falls) until the available length again just chokes at the
exit. The lesson for long pneumatic and natural-gas lines is that
friction, not only a downstream valve, can set the maximum throughput.

\paragraph{Worked example 8: Rayleigh flow and thermal choking in a
combustor.} A frictionless constant-area combustor takes air at
$M_1 = 0.20$, $T_1 = 500\,\si{\kelvin}$, $T_{01} = 504\,\si{\kelvin}$
and adds heat (figure~\ref{fig:vol03ch13:fanno-rayleigh}). We want the
heat addition per unit mass that drives the exit to $M = 1$ (thermal
choking). Rayleigh theory gives the ratio of stagnation temperature to
its sonic value,
\[
\frac{T_0}{T_0^*} = \frac{(\gamma + 1)M^2\,[2 + (\gamma - 1)M^2]}
{(1 + \gamma M^2)^2}.
\]
At $M_1 = 0.20$: numerator $= 2.4\times 0.04\times(2 + 0.4\times 0.04)
= 0.096\times 2.016 = 0.1935$, denominator $= (1 + 1.4\times 0.04)^2 =
1.056^2 = 1.115$, so $T_{01}/T_0^* = 0.1735$. Hence
$T_0^* = T_{01}/0.1735 = 504/0.1735 = 2905\,\si{\kelvin}$, the
stagnation temperature at the choked exit. The heat added per unit
mass is $q = c_p(T_0^* - T_{01}) = 1005\times(2905 - 504) =
2.41\,\si{\mega\joule\per\kilogram}$. Adding more heat than this cannot
raise the exit past $M = 1$; the excess instead pushes the choke point
upstream and reduces the captured mass flow, the Rayleigh-flow analogue
of the friction limit in worked example~7. Thermal choking is the
mechanism that caps the heat release a ramjet or scramjet combustor can
accept at a given inlet Mach number.

\paragraph{Worked example 9: oblique shock with a detachment check.}
A wedge of half-angle $\theta = 25^{\circ}$ sits in a stream at
$M_1 = 2.0$ (figure~\ref{fig:vol03ch13:flat-plate-supersonic} shows the
companion flat-plate geometry; here the surface is a single wedge
face). We want the shock angle, the downstream Mach number, and the
pressure ratio, and we want to know first whether an attached shock
exists at all. The maximum deflection at $M_1 = 2.0$ is
$\theta_{\max} = 22.97^{\circ}$ (\texttt{code/oblique\_shock.py},
found by scanning $\beta$ from the Mach angle to $90^{\circ}$). Since
the required turn $25^{\circ}$ exceeds $\theta_{\max} = 22.97^{\circ}$,
no attached oblique shock can turn the flow: the shock detaches and
stands as a curved bow shock ahead of the wedge. The lesson is to run
the detachment check before solving for $\beta$.
% Figure: shock-expansion flow over a flat-plate airfoil at angle of
% attack in a supersonic stream. Lower (windward) surface compresses
% through a leading-edge oblique shock; upper (leeward) surface expands
% through a leading-edge Prandtl-Meyer fan. Trailing edge recombines
% through a shock (lower) and a fan (upper). Schematic.
\begin{figure}[htbp]
\centering
\begin{tikzpicture}[scale=1.1]
% the plate, at a small angle of attack
\draw[engblue, very thick] (0,0) -- (4,0.7);
\node[engblue, font=\scriptsize, anchor=north] at (2.0,0.2) {flat plate};
% angle of attack mark relative to the freestream (horizontal)
\draw[enggray, dashed] (0,0) -- (4.4,0);
\draw[enggray, -{Latex[length=1.6mm]}] (1.3,0) arc (0:10:1.3);
\node[enggray, font=\scriptsize] at (1.65,0.13) {$\alpha$};
% freestream arrow
\draw[engblack, -{Latex[length=2mm]}, thick] (-1.6,0.35) -- (-0.5,0.35);
\node[engblack, font=\scriptsize, anchor=south] at (-1.05,0.37) {$M_\infty$};
% leading-edge oblique shock on the windward (lower) side
\draw[engred, thick] (0,0) -- (2.2,-1.7);
\node[engred, font=\scriptsize, anchor=west] at (1.3,-1.05) {LE shock};
\node[engblack, font=\scriptsize] at (1.1,-0.35) {compression};
% leading-edge expansion fan on the leeward (upper) side
\draw[englime, thick] (0,0) -- (1.3,1.6);
\draw[englime, thick] (0,0) -- (1.9,1.4);
\draw[englime, thick] (0,0) -- (2.4,1.05);
\node[englime, font=\scriptsize, anchor=south west] at (1.4,1.45) {LE fan};
\node[engblack, font=\scriptsize] at (1.3,0.75) {expansion};
% trailing-edge waves
% upper: recompression shock; lower: expansion to align wakes
\draw[engamber, thick] (4,0.7) -- (5.4,1.7);
\node[engamber, font=\scriptsize, anchor=south] at (5.0,1.35) {TE shock};
\draw[englime, thick] (4,0.7) -- (5.4,-0.2);
\draw[englime, thick] (4,0.7) -- (5.4,0.25);
\node[englime, font=\scriptsize, anchor=north west] at (4.7,0.15) {TE fan};
\end{tikzpicture}
\caption[Shock-expansion flow over a flat plate]{Supersonic flow over a
flat-plate airfoil at angle of attack $\alpha$, the simplest case the
shock-expansion method solves exactly. The windward (lower) surface
turns into the flow, so the leading edge throws an oblique compression
shock and the surface pressure rises above freestream. The leeward
(upper) surface turns away from the flow, so the leading edge opens a
Prandtl-Meyer expansion fan and the surface pressure falls below
freestream. The pressure difference between the two faces is the lift,
and because both faces are inclined to the freestream the same pressures
project a streamwise component, the wave drag. At the trailing edge the
two streams realign through a recompression shock above and an
expansion fan below. For small $\alpha$ the exact shock-expansion
result collapses to the linear Ackeret values
$C_L=4\alpha/\sqrt{M_\infty^2-1}$ and
$C_{D,\text{wave}}=4\alpha^2/\sqrt{M_\infty^2-1}$ worked in the text.}
\label{fig:vol03ch13:flat-plate-supersonic}
\end{figure}
 We reduce the wedge to
$\theta = 18^{\circ}$, below $\theta_{\max}$, and proceed. The
$\theta$-$\beta$-$M$ relation then has two roots: the weak shock
$\beta = 45.3^{\circ}$ and the strong shock $\beta = 79.8^{\circ}$. The
weak root is the one a wedge in unconfined flow realises. Its upstream
normal component is $M_{n1} = 2.0 \sin 45.3^{\circ} = 1.421$, so from
the normal-shock relations $M_{n2} = 0.731$, $p_2/p_1 = 2.19$, and the
downstream Mach number is
$M_2 = M_{n2}/\sin(\beta - \theta) = 0.731/\sin 27.3^{\circ} = 1.59$,
still supersonic. A wedge just below its detachment angle gives a
strong attached shock that nonetheless leaves the flow supersonic; push
the angle a few degrees higher and the shock jumps to the detached bow
configuration, with subsonic flow over the nose and a large rise in
drag.

\paragraph{Worked example 10: supersonic pitot tube (the Rayleigh
pitot formula).} A pitot tube on a supersonic aircraft at
$M_\infty = 1.8$ reads a stagnation pressure
$p_{02} = 60\,\text{kPa}$ at its mouth. We want the freestream static
pressure $p_1$, given that a detached normal shock stands ahead of the
pitot mouth. A subsonic pitot tube reads the isentropic stagnation
pressure, but at supersonic speed the flow that reaches the mouth has
crossed the bow shock first, so the reading is the post-shock
stagnation pressure $p_{02}$, not the freestream $p_{01}$. The
\engterm{Rayleigh pitot formula} combines the normal-shock
stagnation-loss with the isentropic compression behind the shock:
\[
\frac{p_{02}}{p_1} = \left[\frac{(\gamma + 1)M_1^2}{2}\right]^{
\gamma/(\gamma - 1)}\left[\frac{\gamma + 1}{2\gamma M_1^2 -
(\gamma - 1)}\right]^{1/(\gamma - 1)}.
\]
At $M_1 = 1.8$, $\gamma = 1.4$: the first bracket is
$[(2.4)(3.24)/2]^{3.5} = 3.888^{3.5} = 112.0$, and the second bracket
is $[2.4/(2.8 \times 3.24 - 0.4)]^{2.5} =
[2.4/8.672]^{2.5} = 0.2768^{2.5} = 0.0403$, so
$p_{02}/p_1 = 112.0 \times 0.0403 = 4.51$. Hence the freestream static
pressure is $p_1 = 60/4.51 = 13.3\,\text{kPa}$. Had we wrongly used the
subsonic isentropic relation $p_{0}/p = (1 + 0.2 M^2)^{3.5} =
1.648^{3.5} = 5.75$, we would have read $p_1 = 60/5.75 =
10.4\,\text{kPa}$, an error of $22\,\%$. The supersonic air-data system
must apply the Rayleigh formula, not the isentropic relation, above
$M = 1$, the data-reduction counterpart of worked example~1's
subsonic-compressible caution.

\paragraph{Worked example 11: lift and drag of a flat-plate airfoil by
shock-expansion theory.} A flat-plate airfoil flies at $M_\infty = 2.0$
and angle of attack $\alpha = 8^{\circ}$
(figure~\ref{fig:vol03ch13:flat-plate-supersonic}). We compute the lift
and wave-drag coefficients exactly by shock-expansion theory and
compare with the linear Ackeret result, to see how far the linear
theory has drifted at this larger angle. The windward (lower) surface
turns into the flow by $8^{\circ}$: an oblique shock at
$M_1 = 2.0$, $\theta = 8^{\circ}$ gives $\beta = 37.2^{\circ}$,
$M_{n1} = 2.0 \sin 37.2^{\circ} = 1.209$, and
$p_{\text{lower}}/p_\infty = 1 + \frac{2\gamma}{\gamma + 1}
(M_{n1}^2 - 1) = 1 + 1.1667(1.462 - 1) = 1.539$. The leeward (upper)
surface turns away by $8^{\circ}$: a Prandtl-Meyer expansion from
$M_\infty = 2.0$ with $\nu(2.0) = 26.38^{\circ}$ reaches
$\nu = 34.38^{\circ}$, which inverts to $M = 2.32$
(\texttt{code/prandtl\_meyer.py}), and the isentropic pressure ratio
referenced to the common stagnation pressure is
$p_{\text{upper}}/p_\infty = (1 + 0.2 \times 2.0^2)^{3.5}/
(1 + 0.2 \times 2.32^2)^{3.5} = 1.8^{3.5}/2.076^{3.5} = 0.620$. The
net pressure coefficient on the plate is
$(p_{\text{lower}} - p_{\text{upper}})/(\tfrac{1}{2}\gamma p_\infty
M_\infty^2) = (1.539 - 0.620)/(0.7 \times 4.0) = 0.328$. Resolving the
normal force into lift and drag,
$C_L = 0.328 \cos 8^{\circ} = 0.325$ and
$C_{D,\text{wave}} = 0.328 \sin 8^{\circ} = 0.0457$. The Ackeret
linear theory gives $C_L = 4\alpha/\beta =
4 \times 0.1396/\sqrt{3} = 0.322$ and
$C_{D,\text{wave}} = 4\alpha^2/\beta = 0.0450$. The two agree to about
$1\,\%$ for lift and $1.5\,\%$ for drag even at $8^{\circ}$, because
the flat plate's compression and expansion turns are equal and their
nonlinearities partly cancel; on a thick or cambered section, where the
turns are unequal, the linear theory drifts further and the
shock-expansion result is the one to trust.

\paragraph{Worked example 12: a moving normal shock and the shock-tube
problem.} The normal-shock relations of section 13.4 are written for a
shock standing still in a steady stream, but many real shocks move: the
blast wave from an explosion, the shock that runs down a shock tube, the
starting shock swallowed by a supersonic inlet. A moving shock is
reduced to the steady case by riding with it. Consider a normal shock
propagating at speed $W$ into still air at rest ($u_1 = 0$,
$p_1 = 100\,\text{kPa}$, $T_1 = 288\,\si{\kelvin}$,
$a_1 = \sqrt{1.4\times 287\times 288} = 340\,\si{\meter\per\second}$),
moving fast enough that the air it has passed is set into motion behind
it. We want the speed $W$, the pressure behind the shock, and the
induced gas velocity $u_p$ for a shock that raises the pressure by a
factor $p_2/p_1 = 4.5$, the kind of overpressure a strong blast
delivers. Transform to the shock frame by subtracting $W$ from every
velocity: in that frame the air approaches at $M_s = W/a_1$ and the
steady relations apply with upstream Mach number $M_s$. Inverting the
steady pressure relation for the shock Mach number,
\[
\frac{p_2}{p_1} = 1 + \frac{2\gamma}{\gamma + 1}(M_s^2 - 1)
\;\Longrightarrow\;
M_s^2 = 1 + \frac{\gamma + 1}{2\gamma}\!\left(\frac{p_2}{p_1} - 1\right)
= 1 + \frac{2.4}{2.8}(3.5) = 4.00,
\]
so $M_s = 2.00$ and the shock runs at $W = M_s a_1 = 680\,\si{\meter
\per\second}$, twice the ambient sound speed. The gas behind the shock
is not at rest: in the lab frame it is dragged forward at the
\engterm{induced} (piston) velocity
\[
u_p = W - u_2' = W\!\left(1 - \frac{\rho_1}{\rho_2}\right),
\]
where $u_2' = W\,\rho_1/\rho_2$ is the shock-frame downstream speed and
the density ratio is the steady value at $M_s = 2.0$,
$\rho_2/\rho_1 = (\gamma + 1)M_s^2/[(\gamma - 1)M_s^2 + 2] =
2.4\times 4/(0.4\times 4 + 2) = 9.6/3.6 = 2.667$. Hence
$u_p = 680\,(1 - 1/2.667) = 680\times 0.625 = 425\,\si{\meter\per
\second}$. The air behind a Mach-2 blast front is itself moving at
$425\,\si{\meter\per\second}$, a dynamic pressure
$\tfrac{1}{2}\rho_2 u_p^2$ that is the destructive agent in blast
loading, distinct from the static overpressure $p_2 - p_1$. The same
transformation runs the shock tube: a diaphragm bursts between a
high-pressure driver and a low-pressure driven section, and the shock
that races into the driven gas is exactly this moving normal shock, its
strength set by the diaphragm pressure ratio through the shock-tube
equation that matches the induced velocity behind the shock to the
velocity behind the expansion fan running back into the driver. The
lesson is that a moving shock carries no new physics, only a change of
frame, and that the induced flow behind it is what a structure in the
path actually feels.

\paragraph{Worked example 13: one nozzle, four back pressures.} A
single converging-diverging nozzle has a fixed exit-to-throat area
ratio $A_e/A^* = 2.0$ and is fed from a reservoir at
$p_0 = 100\,\text{kPa}$, $T_0 = 300\,\si{\kelvin}$ (air,
$\gamma = 1.4$). The same hardware behaves in four different ways as
the back pressure $p_b$ is lowered, and the area ratio alone fixes the
boundaries between them (figure~\ref{fig:vol03ch13:nozzle-modes}). The
area ratio $A_e/A^* = 2.0$ admits two isentropic exit Mach numbers,
read from the area-Mach relation (\texttt{code/nozzle.py}): the
subsonic root $M_e = 0.306$ and the supersonic root $M_e = 2.20$. The
two matched exit pressures are
\[
p_{e,\text{sub}} = p_0(1 + 0.2\times 0.306^2)^{-3.5} = 100\times 0.937
= 93.7\,\text{kPa},
\]
\[
p_{e,\text{sup}} = p_0(1 + 0.2\times 2.20^2)^{-3.5} = 100\times 0.0935
= 9.35\,\text{kPa}.
\]
A third reference pressure is the one that just chokes the throat
without supersonic flow downstream, which sits at the same subsonic
exit pressure $93.7\,\text{kPa}$ because the choked subsonic branch and
the unchoked-limit branch meet there. The four regimes are then:
\begin{itemize}[leftmargin=*]
  \item $p_b > 93.7\,\text{kPa}$: the throat is not choked, the whole
    nozzle is subsonic, and the diverging section acts as a diffuser.
    The exit pressure equals $p_b$.
  \item $p_b = 93.7\,\text{kPa}$ down to the pressure that puts a normal
    shock at the exit plane: the throat is choked, the flow goes
    supersonic past the throat, and a normal shock stands inside the
    diverging section, moving toward the exit as $p_b$ falls. The
    shock reaches the exit plane when the post-shock pressure equals
    $p_b$; the supersonic pre-shock Mach number there is $M = 2.20$, and
    a normal shock at $M_1 = 2.20$ gives $p_2/p_1 = 5.48$, so the
    exit-plane-shock back pressure is $5.48\times 9.35 =
    51.2\,\text{kPa}$.
  \item $51.2\,\text{kPa} > p_b > 9.35\,\text{kPa}$: overexpanded. The
    flow is fully supersonic to the exit at $M_e = 2.20$,
    $p_e = 9.35\,\text{kPa}$, and recompresses to $p_b$ through oblique
    shocks outside the nozzle.
  \item $p_b < 9.35\,\text{kPa}$: underexpanded. The exit is supersonic
    and matched-or-over, and the flow expands further through
    Prandtl-Meyer fans outside the lip; $p_b = 9.35\,\text{kPa}$
    exactly is the design (matched) condition.
\end{itemize}
The single number $A_e/A^* = 2.0$ has fixed all four boundaries:
$93.7\,\text{kPa}$ (choking onset), $51.2\,\text{kPa}$ (shock at the
exit plane), and $9.35\,\text{kPa}$ (design). A nozzle is not a
device with one behaviour but a piece of fixed geometry whose
operating mode is read off the back pressure against these three
thresholds.

\paragraph{Worked example 14: Newtonian impact pressure on a blunt
body.} At hypersonic speed the shock layer over the windward face of a
blunt body is thin and hugs the surface, and the surface pressure is
estimated without solving the flow field at all by Newton's old
picture: the gas is a stream of particles that strike the surface, lose
the component of momentum normal to it, and slide off tangentially.
\engterm{Newtonian impact theory} gives the surface pressure
coefficient
\[
C_p = 2\sin^2\theta_b,
\]
where $\theta_b$ is the angle between the freestream and the local
surface tangent (so $\theta_b = 90^{\circ}$ at a stagnation point
facing the flow). We apply it to a spherical nose of radius
$R_n = 0.5\,\si{\meter}$ at $M_\infty = 20$,
$p_\infty = 50\,\text{Pa}$, $\rho_\infty = 8\times 10^{-4}\,\si{
\kilogram\per\meter\cubed}$ at a representative entry altitude, and
estimate the pressure-drag coefficient. The dynamic pressure is
$q_\infty = \tfrac{1}{2}\gamma p_\infty M_\infty^2 =
0.7\times 50\times 400 = 14{,}000\,\text{Pa}$. The stagnation
pressure coefficient is $C_{p,\max} = 2$, so the stagnation pressure
is $p_\infty + 2 q_\infty = 50 + 28{,}000 = 28{,}050\,\text{Pa}$,
nearly all of it dynamic. On a sphere the local surface angle measured
from the axis is $\phi$, the tangent makes $\theta_b = 90^{\circ} -
\phi$ with the flow, so $C_p = 2\cos^2\phi$, and the pressure-drag
coefficient referenced to the frontal area $\pi R_n^2$ is the integral
\[
C_D = \frac{1}{\pi R_n^2}\int_0^{\pi/2} C_p \cos\phi\,(2\pi R_n^2
\sin\phi)\,d\phi
= 4\int_0^{\pi/2}\cos^3\phi\,\sin\phi\,d\phi
= 4\left[-\frac{\cos^4\phi}{4}\right]_0^{\pi/2} = 1.0.
\]
The Newtonian drag coefficient of a hemisphere-nosed body is exactly
$C_D = 1$, against the much lower $C_D \approx 0.2$ to $0.4$ of a
streamlined subsonic body: the blunt hypersonic nose is a high-drag
shape by design, and that is the point. The large pressure drag
decelerates the vehicle high in the thin upper atmosphere, bleeding off
kinetic energy where the air is too rarefied to heat the surface
severely, so that less kinetic energy survives to the dense lower
atmosphere where heating is worst. The drag the streamlining instinct
would remove is the same drag that makes survivable re-entry possible,
the theme the next case study carries through with numbers.

\paragraph{Worked example 15: the static-to-stagnation temperature gap
on a cold cruise.} A transport cruises at $M = 0.85$ at
$11\,\si{\kilo\meter}$, where $T_\infty = 217\,\si{\kelvin}$. We want
the static air temperature the wing skin sees away from any stagnation
region, the stagnation temperature at the nose, and the recovery
temperature on the turbulent-boundary-layer surface, to settle which
number a de-icing engineer should design against. The stagnation
temperature is $T_0 = T_\infty(1 + 0.2\times 0.85^2) =
217\times 1.1445 = 248\,\si{\kelvin}$, so the full stagnation rise is
$31\,\si{\kelvin}$. The static temperature in the moving stream is just
$T_\infty = 217\,\si{\kelvin}$: the air the wing flies through is the
cold ambient, not the warm stagnation value, because only the
fraction of the flow actually brought to rest gives up its kinetic
share (figure~\ref{fig:vol03ch13:energy-partition} read at $M = 0.85$,
where the kinetic share is $0.2\times 0.7225/1.1445 = 0.126$, so the
static share is $0.874$ and $T/T_0 = 0.874$, recovering
$217/248 = 0.875$). Over most of the surface the boundary layer
recovers only a fraction $r$ of the kinetic share, the
\engterm{recovery factor}, with $r \approx \mathrm{Pr}^{1/3}\approx
0.89$ for a turbulent layer in air. The recovery (adiabatic-wall)
temperature is therefore $T_{\text{aw}} = T_\infty + r(T_0 - T_\infty)
= 217 + 0.89\times 31 = 245\,\si{\kelvin}$, about $28\,\si{\kelvin}$
above ambient. The lesson for thermal design: a stagnation probe reads
$248\,\si{\kelvin}$, the general skin runs near $245\,\si{\kelvin}$,
and the ambient is $217\,\si{\kelvin}$; the gap of $31\,\si{\kelvin}$
between ambient and stagnation is what keeps a pitot heater modest and
what de-icing logic must respect, since the cold static air can still
ice an unheated surface while the heated probe stays clear.

\paragraph{Worked example 16: the converging nozzle that will not go
supersonic.} A converging-only nozzle is fed from a tank at
$p_0 = 700\,\text{kPa}$, $T_0 = 300\,\si{\kelvin}$ (air) and exhausts
to a back pressure that an operator lowers from $700\,\text{kPa}$
toward vacuum. We want the exit Mach number and exit pressure as the
back pressure falls, and the back pressure below which nothing further
changes upstream. While the back pressure stays above the critical
value $p^* = 0.528\,p_0 = 370\,\text{kPa}$, the exit is subsonic and
the exit pressure equals the back pressure; the exit Mach number rises
as the back pressure falls. At $p_b = 500\,\text{kPa}$, for instance,
$p_0/p_e = 700/500 = 1.40$, and inverting the isentropic pressure
relation $1.40 = (1 + 0.2 M_e^2)^{3.5}$ gives $1 + 0.2 M_e^2 =
1.40^{1/3.5} = 1.0991$, so $M_e^2 = 0.495$ and $M_e = 0.704$. When the
back pressure reaches $p^* = 370\,\text{kPa}$ the exit chokes at
$M_e = 1$, and the mass flow reaches the choked maximum. Lower the back
pressure further, to $200\,\text{kPa}$ or to vacuum, and the exit stays
at $M_e = 1$ with exit pressure pinned at $p^* = 370\,\text{kPa}$: a
purely converging nozzle cannot accelerate the flow past sonic no
matter how low the back pressure goes, because the area-Mach relation
forbids supersonic flow in a contracting duct. The mismatch between the
exit pressure $370\,\text{kPa}$ and the lower back pressure is taken up
by an expansion fan outside the nozzle lip. The supersonic exit a
rocket needs requires the diverging section the de Laval nozzle adds.

\paragraph{Worked example 17: the area ratio that places a shock at a
named station.} A converging-diverging nozzle has $A_e/A^* = 3.0$ and
is fed at $p_0 = 600\,\text{kPa}$, $T_0 = 350\,\si{\kelvin}$ (air). A
normal shock is observed to stand at the area station $A_x/A^* = 2.0$ in
the diverging section. We want the back pressure that holds the shock
there. March the supersonic branch from the throat to $A_x/A^* = 2.0$:
inverting the area-Mach relation on the supersonic root gives the
pre-shock Mach number $M_x = 2.20$ (\texttt{code/nozzle.py}). The
normal-shock relations at $M_1 = 2.20$ give the post-shock Mach number
$M_{x'} = 0.547$ and the stagnation-pressure ratio across the shock
$p_{0,x'}/p_0 = 0.628$, so the post-shock stagnation pressure is
$0.628\times 600 = 377\,\text{kPa}$. Behind the shock the flow is
subsonic and isentropic, so it diffuses from $A_x/A^* = 2.0$ to the
exit $A_e/A^* = 3.0$ along the subsonic branch of a new sonic
reference area $A^{**}$ tied to the reduced stagnation pressure. The
post-shock subsonic flow has $A_x/A^{**} = (A_x/A^*)(p_0/p_{0,x'})$
scaled by the area-Mach value at $M_{x'} = 0.547$; carrying the
arithmetic, the exit Mach number is $M_e = 0.32$ and the exit pressure
is $p_e = p_{0,x'}(1 + 0.2\times 0.32^2)^{-3.5} = 377\times 0.930 =
351\,\text{kPa}$. The back pressure must equal this exit pressure,
$p_b = 351\,\text{kPa}$, to hold the shock at $A_x/A^* = 2.0$. Raise
$p_b$ above $351\,\text{kPa}$ and the shock moves upstream toward the
throat; lower it and the shock moves toward the exit and then out the
lip, the machinery worked example~13 mapped and that
figure~\ref{fig:vol03ch13:nozzle-modes} draws.

\paragraph{Worked example 18: stacking three weak shocks against one
strong one.} An inlet must decelerate an $M_\infty = 4.0$ stream to
subsonic. We compare the stagnation-pressure recovery of a single
normal shock against three equal-deflection oblique shocks of
$\theta = 10^{\circ}$ each followed by a terminal normal shock, to show
how steeply the cube law rewards splitting the compression. A single
normal shock at $M_1 = 4.0$ gives $p_{02}/p_{01} = 0.139$
(\texttt{code/normal\_shock.py}): seven eighths of the stagnation
pressure is gone in one jump. The staircase keeps far more. Ramp~1 at
$M_1 = 4.0$, $\theta = 10^{\circ}$: weak shock angle $\beta_1 =
22.2^{\circ}$, $M_{n1} = 4.0\sin 22.2^{\circ} = 1.51$, loss
$p_0\text{-ratio} = 0.930$, deflected $M_2 = 3.32$. Ramp~2 at
$M = 3.32$, $\theta = 10^{\circ}$: $\beta_2 = 25.6^{\circ}$, $M_{n} =
3.32\sin 25.6^{\circ} = 1.43$, loss $0.951$, deflected $M_3 = 2.80$.
Ramp~3 at $M = 2.80$, $\theta = 10^{\circ}$: $\beta_3 = 29.5^{\circ}$,
$M_{n} = 2.80\sin 29.5^{\circ} = 1.38$, loss $0.961$, deflected
$M_4 = 2.35$. Terminal normal shock at $M = 2.35$: loss $0.563$. The
overall recovery is $0.930\times 0.951\times 0.961\times 0.563 =
0.479$, against $0.139$ for the single normal shock, a factor of $3.4$.
The three oblique shocks together cost only $1 - 0.930\times 0.951
\times 0.961 = 0.150$ of the stagnation pressure while the single
terminal shock at the reduced Mach number costs the rest; the cube
scaling of entropy rise with $(M_n^2 - 1)$ is what makes three small
compressions so much cheaper than one large one. The penalty is
length and a narrow design-Mach band, the trade every high-speed inlet
weighs.

\paragraph{Worked example 19: Prandtl-Meyer expansion into near-vacuum.}
A small attitude-control thruster vents air at $M_1 = 1.5$,
$p_1 = 40\,\text{kPa}$ into the near-vacuum of high altitude, where the
ambient pressure is $p_a = 1\,\text{kPa}$. We want the Mach number the
flow reaches as it expands to ambient, and the total turning angle the
expansion fan applies, to size the lip clearance the plume needs. The
expansion is isentropic, so the stagnation pressure is conserved:
$p_0 = p_1(1 + 0.2\times 1.5^2)^{3.5} = 40\times 1.45^{3.5} =
40\times 3.671 = 147\,\text{kPa}$. Expanding to $p_a = 1\,\text{kPa}$
gives $p_0/p_a = 147$, and inverting the isentropic relation,
$1 + 0.2 M_2^2 = 147^{1/3.5} = 4.19$, so $M_2^2 = 15.9$ and
$M_2 = 3.99$. The turning angle is the change in the Prandtl-Meyer
function: $\nu(1.5) = 11.91^{\circ}$ and $\nu(3.99) = 65.8^{\circ}$
(\texttt{code/prandtl\_meyer.py}), so the fan turns the flow through
$\Delta\theta = 65.8 - 11.9 = 53.9^{\circ}$. The plume swings through
nearly $54^{\circ}$ as it expands from the lip to the ambient
pressure, which is why a thruster venting into vacuum throws a wide
plume that can impinge on solar panels and antennas placed near it: the
lower the ambient pressure, the larger the turning, approaching the
$\nu_{\max} = 130.5^{\circ}$ limit as $p_a \to 0$. Plume impingement is
a routine spacecraft-layout constraint, and it is read directly off the
Prandtl-Meyer turning the local pressure ratio demands.

\paragraph{Worked example 20: the speed of sound is a property of the
gas, not just its temperature.} A pipeline carries helium, a launch
vehicle burns a gas with $\gamma = 1.20$ and $R = 350\,\si{\joule\per
\kilogram\per\kelvin}$, and a boiler runs superheated steam; all three
sit near $500\,\si{\kelvin}$, yet their sound speeds differ widely. We
want $a = \sqrt{\gamma R T}$ for each, to see why the working gas, not
the temperature alone, sets the Mach number a given velocity produces.
For air ($\gamma = 1.40$, $R = 287$) at $500\,\si{\kelvin}$,
$a = \sqrt{1.40\times 287\times 500} = \sqrt{2.01\times 10^{5}} =
448\,\si{\meter\per\second}$. For helium ($\gamma = 1.667$,
$R = 2077\,\si{\joule\per\kilogram\per\kelvin}$, since a light monatomic
gas has both a high $\gamma$ and a large specific gas constant),
$a = \sqrt{1.667\times 2077\times 500} = \sqrt{1.73\times 10^{6}} =
1316\,\si{\meter\per\second}$, nearly three times the value in air at
the same temperature. For the combustion gas ($\gamma = 1.20$,
$R = 350$), $a = \sqrt{1.20\times 350\times 500} = \sqrt{2.10\times
10^{5}} = 458\,\si{\meter\per\second}$. A flow at
$450\,\si{\meter\per\second}$ is thus mildly supersonic in air
($M = 1.00$), firmly subsonic in helium ($M = 0.34$), and just subsonic
in the combustion gas ($M = 0.98$). The lesson is that the Mach number a
velocity produces is set by $\sqrt{\gamma R}$, so a helium wind tunnel
reaches a target Mach number only at a far higher velocity than an air
tunnel, and a rocket's hot, heavy-molecule products (low $R$) give a
lower sound speed than the cold air outside, which is why the exit Mach
number and the plume geometry must be worked in the combustion gas's
own properties, as the nozzle case study did with $\gamma = 1.20$.

\paragraph{Worked example 21: a critical-flow venturi as a mass-flow
standard.} A metrology laboratory meters nitrogen with a sonic
(critical-flow) venturi of throat diameter $d^* = 8.0\,\si{\milli
\meter}$, fed at $p_0 = 400\,\text{kPa}$, $T_0 = 293\,\si{\kelvin}$
($\gamma = 1.40$, $R = 297\,\si{\joule\per\kilogram\per\kelvin}$ for
nitrogen), with a measured discharge coefficient $C_d = 0.994$. We want
the delivered mass flow and the highest back pressure at which the meter
still reads true. The throat area is $A^* = \pi(0.004)^2 =
5.03\times 10^{-5}\,\si{\meter}^2$. The ideal choked mass flux is
\[
\frac{\dot m_{\text{ideal}}}{A^*} = \frac{p_0}{\sqrt{T_0}}\sqrt{\frac{
\gamma}{R}}\left(\frac{2}{\gamma + 1}\right)^{(\gamma + 1)/[2(\gamma -
1)]}.
\]
The radical is $\sqrt{1.40/297} = \sqrt{4.71\times 10^{-3}} = 0.0687$,
the bracket is $(2/2.4)^{3} = 0.579$, and
$p_0/\sqrt{T_0} = 4.0\times 10^{5}/\sqrt{293} = 2.34\times 10^{4}$, so
$\dot m_{\text{ideal}}/A^* = 2.34\times 10^{4}\times 0.0687\times 0.579
= 930\,\si{\kilogram\per\second}/\si{\meter}^2$. The ideal mass flow is
$\dot m_{\text{ideal}} = 930\times 5.03\times 10^{-5} =
0.0468\,\si{\kilogram\per\second}$, and the delivered flow is
$C_d$ times this, $\dot m = 0.994\times 0.0468 =
0.0465\,\si{\kilogram\per\second}$, about $46.5\,\text{g/s}$. The meter
reads true only while the throat stays choked, which needs the back
pressure below the critical value; accounting for the pressure the
diverging cone recovers, a sonic venturi holds its reading up to a back
pressure of roughly $0.85\,p_0 = 340\,\text{kPa}$, far above the plain
$0.528\,p_0$ of an unrecovered orifice
(figure~\ref{fig:vol03ch13:critical-venturi}). That combination, a mass
flow fixed by upstream measurements alone and a wide choked range from
the pressure-recovering cone, is why the critical-flow venturi is a
primary gas-flow standard against which other flow meters are
calibrated.

% Figure: critical-flow (sonic) venturi as a mass-flow standard.
% Left: contour with a rounded converging inlet, a sonic throat, and a
% gentle diverging recovery cone. Right: mass flow versus back-pressure
% ratio, flat (choked) once the throat reaches M=1.
\begin{figure}[htbp]
\centering
\begin{tikzpicture}
% ---- left panel: the venturi contour ----
\begin{scope}[xshift=0cm]
% upper and lower walls, converging then diverging about a throat at x=2.4
\draw[engblack, very thick]
  (0,1.1) .. controls (1.3,1.1) and (2.0,0.35) .. (2.4,0.35)
  .. controls (3.4,0.35) and (4.8,0.85) .. (5.6,0.95);
\draw[engblack, very thick]
  (0,-1.1) .. controls (1.3,-1.1) and (2.0,-0.35) .. (2.4,-0.35)
  .. controls (3.4,-0.35) and (4.8,-0.85) .. (5.6,-0.95);
% flow arrow
\draw[enggray, -{Latex[length=1.8mm]}] (-0.6,0) -- (0.6,0);
\node[enggray, font=\scriptsize, anchor=south] at (0,0.05) {$p_0,T_0$};
% throat marker
\draw[engred, dashed] (2.4,-0.35) -- (2.4,0.35);
\node[engred, font=\scriptsize, anchor=south] at (2.4,0.4) {throat $M=1$};
\node[font=\scriptsize, anchor=north] at (1.2,-1.1) {converging};
\node[font=\scriptsize, anchor=north] at (4.2,-0.9) {diverging recovery};
\node[engblue, font=\scriptsize, anchor=west] at (5.7,0) {$p_b$};
\end{scope}
% ---- right panel: mass flow vs pressure ratio ----
\begin{scope}[xshift=8.0cm, yshift=-1.4cm]
\draw[engblack, -{Latex[length=1.5mm]}] (0,0) -- (4.4,0) node[anchor=north, font=\scriptsize] {$p_b/p_0$};
\draw[engblack, -{Latex[length=1.5mm]}] (0,0) -- (0,3.0) node[anchor=east, font=\scriptsize] {$\dot m$};
% choked plateau for p_b/p0 < 0.528, then falls to zero at p_b/p0 = 1
% x axis: 0 at left ... 4 units = p_b/p0 from 0 to 1
% critical ratio 0.528 -> x = 0.528*4 = 2.11
\draw[engblue, very thick] (0,2.4) -- (2.11,2.4);
\draw[engblue, very thick] (2.11,2.4) .. controls (3.1,2.0) and (3.7,1.0) .. (4.0,0);
\draw[engred, dashed] (2.11,0) -- (2.11,2.4);
\node[engred, font=\scriptsize, anchor=north] at (2.11,-0.05) {$0.528$};
\node[font=\scriptsize, anchor=north] at (4.0,-0.05) {$1$};
\node[engblue, font=\scriptsize, anchor=south] at (1.0,2.4) {choked ($\dot m_{\max}$)};
\node[engblue, font=\scriptsize, anchor=west] at (3.0,1.5) {unchoked};
\end{scope}
\end{tikzpicture}
\caption[The critical-flow venturi as a mass-flow standard]{The
critical-flow (sonic) venturi. A rounded converging inlet accelerates
the gas to $M=1$ at a smoothly contoured throat, and a gentle diverging
cone recovers most of the pressure so the device runs on a small
overall pressure drop. Once the throat is choked, the mass flow (right)
depends only on the upstream stagnation state $p_0$, $T_0$ and the
throat area, and is flat against any further fall in the back pressure
$p_b$ below the critical ratio $p_b/p_0 = 0.528$ for air. That
insensitivity to downstream conditions is what makes the sonic venturi
a primary gas-flow standard: it delivers a mass flow set by upstream
measurements alone, corrected only by a discharge coefficient close to
unity that accounts for the thin throat boundary layer.}
\label{fig:vol03ch13:critical-venturi}
\end{figure}


\paragraph{Worked example 22: regular reflection of an oblique shock
from a wall.} A shock generated by a $\theta = 10^{\circ}$ wedge in a
$M_1 = 2.8$ stream strikes a flat wall parallel to the freestream, and
the wall turns the flow back through a reflected shock
(figure~\ref{fig:vol03ch13:shock-reflection}). We want the reflected
shock angle and the final Mach number, and we check that regular
reflection is possible at all. The incident shock at $M_1 = 2.8$,
$\theta = 10^{\circ}$ has weak-shock angle $\beta_1 = 29.6^{\circ}$
(\texttt{code/oblique\_shock.py}), so $M_{n1} = 2.8\sin 29.6^{\circ} =
1.383$, and the normal-shock relations give the deflected Mach number
$M_2 = M_{n2}/\sin(\beta_1 - \theta) = 0.744/\sin 19.6^{\circ} = 2.22$.
The reflected shock must turn this $M_2 = 2.22$ flow back through the
same $\theta = 10^{\circ}$ to run parallel to the wall. The maximum
deflection at $M_2 = 2.22$ is $\theta_{\max} \approx 26^{\circ}$, well
above the required $10^{\circ}$, so regular reflection exists. Solving
the $\theta$-$\beta$-$M$ relation at $M_2 = 2.22$, $\theta = 10^{\circ}$
gives the reflected weak-shock angle $\beta_2 = 35.6^{\circ}$
\emph{measured from the region-2 flow}, so $M_{n2,\text{in}} =
2.22\sin 35.6^{\circ} = 1.293$, and the final Mach number is
$M_3 = 0.789/\sin(35.6^{\circ} - 10^{\circ}) = 0.789/\sin 25.6^{\circ}
= 1.83$. The reflected shock stands at a steeper angle to its oncoming
flow than the incident one did, because it acts on the lower Mach number
$M_2 < M_1$. Push the incident wedge angle up or the Mach number down
until $M_2$ can no longer be turned through $\theta$, and regular
reflection gives way to a Mach stem, the mastery box on reflection
patterns describes. The overall static-pressure rise from region 1 to
region 3 is the product of the two shock jumps,
$(p_2/p_1)(p_3/p_2) = 2.05\times 1.90 = 3.90$, so the wall sees nearly
four times the freestream static pressure after both shocks.

% Figure: regular reflection of an oblique shock from a flat wall.
% Incident shock turns the flow toward the wall; the reflected shock
% turns it back parallel, so the wall boundary condition is met by two
% shocks rather than one.
\begin{figure}[htbp]
\centering
\begin{tikzpicture}[scale=1.0]
% upper wall (the reflecting surface is the lower wall here)
\draw[engblack, very thick] (0,0) -- (9,0);
\node[engblack, font=\scriptsize, anchor=north] at (4.5,-0.05) {reflecting wall};
% incident shock from a leading wedge at top-left down to reflection point R
\coordinate (R) at (5.0,0);
\draw[engred, very thick] (1.2,3.0) -- (R);
\node[engred, font=\scriptsize, anchor=south east] at (3.0,1.9) {incident shock};
% reflected shock from R up to the right
\draw[engblue, very thick] (R) -- (8.2,2.6);
\node[engblue, font=\scriptsize, anchor=south west] at (6.6,1.4) {reflected shock};
% region labels
\node[font=\scriptsize] at (1.2,0.6) {region 1};
\node[font=\scriptsize] at (4.6,1.9) {region 2};
\node[font=\scriptsize] at (7.6,0.6) {region 3};
% incoming streamline arrows (region 1 parallel to wall)
\draw[enggray, -{Latex[length=1.5mm]}] (0.4,0.5) -- (2.2,0.5);
\node[enggray, font=\scriptsize, anchor=south] at (1.0,0.55) {$M_1$};
% region 2 streamline turned toward wall by theta
\draw[enggray, -{Latex[length=1.5mm]}] (3.6,1.1) -- ({3.6+1.8*cos(-10)},{1.1+1.8*sin(-10)});
\node[enggray, font=\scriptsize, anchor=south] at (4.4,1.15) {$M_2$};
% region 3 streamline turned back parallel
\draw[enggray, -{Latex[length=1.5mm]}] (6.4,0.5) -- (8.2,0.5);
\node[enggray, font=\scriptsize, anchor=south] at (7.2,0.55) {$M_3$};
% deflection angles marked at R
\draw[engred] (R) ++(150:0.9) arc (150:180:0.9);
\node[engred, font=\scriptsize] at (4.0,0.35) {$\beta_1$};
\draw[engblue] (R) ++(0:0.9) arc (0:28:0.9);
\node[engblue, font=\scriptsize] at (5.9,0.35) {$\beta_2$};
\end{tikzpicture}
\caption[Regular reflection of an oblique shock from a wall]{Regular
reflection of an oblique shock from a straight wall. The incident shock
(struck off an upstream wedge, not shown) turns the region-1 flow
toward the wall by a deflection $\theta$, giving region~2 a velocity
that is no longer parallel to the wall. The wall requires the flow to
run parallel to it, so a reflected shock springs from the reflection
point and turns the flow back through the same angle $\theta$, leaving
region~3 parallel to the wall. The reflected shock is not the mirror
image of the incident one: it sees a lower Mach number $M_2 < M_1$, so
for the same deflection it stands at a steeper angle to the local flow.
When $M_2$ is too low to turn the flow through $\theta$ with any
attached shock, regular reflection fails and a Mach stem forms, the
transition to Mach reflection.}
\label{fig:vol03ch13:shock-reflection}
\end{figure}


\paragraph{Worked example 23: lift and wave drag of a diamond
(double-wedge) airfoil.} A symmetric diamond section of half-thickness
angle $\varepsilon = 5^{\circ}$ (so each fore panel is inclined
$5^{\circ}$ to the chord and each aft panel $-5^{\circ}$) flies at
$M_\infty = 2.0$ and angle of attack $\alpha = 4^{\circ}$
(figure~\ref{fig:vol03ch13:diamond-airfoil}). We estimate the lift and
wave-drag coefficients by linear (Ackeret) panel theory, which assigns
each panel a pressure coefficient $C_p = 2\theta/\sqrt{M_\infty^2 - 1}$
from its own slope relative to the freestream. With
$\beta = \sqrt{2^2 - 1} = 1.732$, the four panel slopes relative to the
flow are: upper fore $\theta = \varepsilon - \alpha = 1^{\circ} =
0.01745\,\text{rad}$ (still into the flow); upper aft
$\theta = -\varepsilon - \alpha = -9^{\circ}$; lower fore
$\theta = \varepsilon + \alpha = 9^{\circ}$; lower aft
$\theta = -\varepsilon + \alpha = -1^{\circ}$. The panel pressure
coefficients are $C_p = 2\theta/\beta$, so (in radians) upper fore
$+0.0202$, upper aft $-0.1814$, lower fore $+0.1814$, lower aft
$-0.0202$. Each panel projects equally onto the chord (half-chord
horizontal extent), so the lift coefficient is the chordwise average of
lower minus upper $C_p$:
$C_L = \tfrac{1}{2}[(C_{p,\text{lf}} - C_{p,\text{uf}}) +
(C_{p,\text{la}} - C_{p,\text{ua}})] =
\tfrac{1}{2}[(0.1814 - 0.0202) + (-0.0202 + 0.1814)] =
\tfrac{1}{2}(0.1612 + 0.1612) = 0.161$, which equals the flat-plate
value $4\alpha/\beta = 4\times 0.0698/1.732 = 0.161$: the thickness
carries no lift, as symmetry requires. The wave drag does not vanish.
Ackeret drag integrates $C_p\theta$ over the surface, giving the closed
form $C_{D,\text{wave}} = \dfrac{4}{\sqrt{M_\infty^2 - 1}}
(\alpha^2 + \overline{\varepsilon^2})$, where
$\overline{\varepsilon^2}$ is the mean-square surface slope from
thickness. Here $\alpha = 0.0698$ and $\varepsilon = 0.0873$, so
$C_{D,\text{wave}} = (4/1.732)(0.0698^2 + 0.0873^2) =
2.309\times(0.00487 + 0.00762) = 0.0289$. The thickness term
$0.00762$ exceeds the lift term $0.00487$: a $5^{\circ}$ half-angle
diamond spends more of its wave drag on being thick than on making lift
at $4^{\circ}$. The flat plate of worked example~4 pays only the lift
term; a real supersonic wing must carry thickness for structure and
fuel, and the thickness wave drag is the price, which is why supersonic
sections are as thin as the structure allows.

% Figure: symmetric double-wedge (diamond) airfoil at angle of attack in
% supersonic flow, with the four surface panels and their wave systems.
% Compression turns (into the flow) run through oblique shocks; expansion
% turns (away from the flow) run through Prandtl-Meyer fans.
\begin{figure}[htbp]
\centering
\begin{tikzpicture}[scale=1.0]
% diamond section: nose at left, tail at right, mid at top and bottom
\coordinate (N) at (0,0);      % nose (leading edge)
\coordinate (T) at (6,0);      % tail (trailing edge)
\coordinate (Up) at (3,0.9);   % upper mid-chord
\coordinate (Lo) at (3,-0.9);  % lower mid-chord
\draw[engblack, very thick] (N) -- (Up) -- (T) -- (Lo) -- cycle;
% panel labels
\node[font=\scriptsize] at (1.4,0.65) {panel 1 (upper fore)};
\node[font=\scriptsize] at (4.6,0.65) {panel 2 (upper aft)};
\node[font=\scriptsize] at (1.4,-0.7) {panel 3 (lower fore)};
\node[font=\scriptsize] at (4.6,-0.7) {panel 4 (lower aft)};
% freestream arrow (slightly below horizontal to indicate positive alpha)
\draw[enggray, -{Latex[length=1.8mm]}] (-2.4,-0.35) -- (-0.7,-0.15);
\node[enggray, font=\scriptsize, anchor=south] at (-1.7,-0.15) {$M_\infty$, $\alpha$};
% leading-edge oblique shocks (upper and lower), red
\draw[engred, thick] (N) -- (-1.1,1.4);
\draw[engred, thick] (N) -- (-1.1,-1.4);
\node[engred, font=\scriptsize, anchor=east] at (-1.1,1.2) {L.E. shocks};
% mid-chord expansion fans (upper and lower), blue fan lines
\foreach \a in {30,45,60} {
  \draw[engblue] (Up) -- ++(\a:1.3);
}
\foreach \a in {-30,-45,-60} {
  \draw[engblue] (Lo) -- ++(\a:1.3);
}
\node[engblue, font=\scriptsize, anchor=south west] at (3.3,1.6) {expansion fans};
% trailing-edge shocks/recompression
\draw[engred, thick] (T) -- (7.1,1.1);
\draw[engred, thick] (T) -- (7.1,-1.1);
\node[engred, font=\scriptsize, anchor=west] at (7.0,0.9) {T.E. shocks};
\end{tikzpicture}
\caption[Double-wedge diamond airfoil in supersonic flow]{A symmetric
double-wedge (diamond) section at angle of attack $\alpha$ in a
supersonic stream. The four straight panels each meet the flow at a
definite angle, and the wave system follows from walking the surface
panel by panel. Each leading-edge face that turns into the flow runs
through an attached oblique shock (red); each mid-chord shoulder that
turns the flow away from itself runs through a Prandtl-Meyer expansion
fan (blue); the flows off the upper and lower aft panels recompress
through trailing-edge shocks to match pressure and direction in the
wake. The surface pressure on each panel is whatever the running tally
of compressions and expansions delivers, and the resultant of the four
panel pressures is the exact inviscid lift and wave drag. Unlike the
flat plate, the diamond carries a thickness (volume) wave drag that
survives even at zero lift, because the fore panels compress and the
aft panels expand whatever the angle of attack.}
\label{fig:vol03ch13:diamond-airfoil}
\end{figure}


\paragraph{Worked example 24: the sonic-boom overpressure of an
N-wave.} A supersonic transport of length $\ell = 60\,\si{\meter}$
cruises at $M = 1.8$ and altitude $h = 16\,\si{\kilo\meter}$. We want an
order-of-magnitude ground overpressure and the duration of the N-wave
(figure~\ref{fig:vol03ch13:sonic-boom-nwave}), to judge it against the
$\sim 100\,\text{Pa}$ threshold above which a boom is heard as sharp.
Whitham's far-field theory gives the peak overpressure of a slender body
as
\[
\Delta p \approx K_p\,\frac{p_a\,(M^2 - 1)^{1/8}}{(h/\ell)^{3/4}}
\left(\frac{\ell}{\ell}\right)\cdot\frac{d}{\ell},
\]
where $p_a$ is ambient pressure, $d/\ell$ a slenderness (diameter-to-
length) ratio, and $K_p$ a shape factor of order unity. Rather than
carry the full similarity constant, we scale from a known case: a
fighter of $\ell \approx 18\,\si{\meter}$ at $M = 1.8$, $h = 16\,\si{
\kilo\meter}$ produces about $\Delta p \approx 100\,\text{Pa}$. The
overpressure scales as $\ell^{3/4}/h^{3/4}$ for geometrically similar
bodies at fixed Mach and, more weakly, with the vehicle's effective
cross-section, so a $60\,\si{\meter}$ transport at the same altitude
scales up by $(60/18)^{3/4} = 3.33^{0.75} = 2.47$ on length but down by
its lower slenderness (a transport is more slender per unit length),
which roughly halves the factor, landing near $\Delta p \approx
100\times 2.47\times 0.5 \approx 120\,\text{Pa}$. The N-wave duration is
the signature length divided by the ground trace speed,
$\Delta t \approx \ell_{\text{sig}}/(M a\cos\mu)$; with a signature
stretched to roughly $1.5\ell$ by propagation and
$a \approx 295\,\si{\meter\per\second}$ at altitude,
$\Delta t \approx (1.5\times 60)/(1.8\times 295) \approx 0.17\,\text{s}$,
a $170\,\text{ms}$ gap between the two bangs. The number confirms why
overland supersonic transport is barred without low-boom shaping: an
unshaped $120\,\text{Pa}$, $170\,\text{ms}$ N-wave startles and rattles
windows, and only reshaping the area-and-lift distribution to break the
single leading shock into a graded rise brings the signature under a
noise-based ceiling.

% Figure: idealised sonic-boom ground pressure signature, the N-wave.
% A leading shock jumps to +dp, the pressure falls linearly through
% ambient to -dp, and a trailing shock jumps back to ambient.
\begin{figure}[htbp]
\centering
\begin{tikzpicture}
\begin{axis}[
  width=12cm, height=6.5cm,
  xlabel={time (arbitrary, ms)}, ylabel={overpressure $p-p_a$ (Pa)},
  xmin=0, xmax=10, ymin=-70, ymax=70,
  axis lines=middle,
  grid=major, grid style={enggray!20},
  tick label style={font=\scriptsize},
  label style={font=\small},
  legend style={font=\scriptsize, at={(0.97,0.03)}, anchor=south east},
]
% ambient before the front
\addplot[engblue, very thick] coordinates {(0,0) (3,0)};
% leading shock: jump up to +50 Pa
\addplot[engred, very thick] coordinates {(3,0) (3,50)};
% linear expansion from +50 down to -50
\addplot[engblue, very thick] coordinates {(3,50) (7,-50)};
% trailing shock: jump back up to 0
\addplot[engred, very thick] coordinates {(7,-50) (7,0)};
% ambient after
\addplot[engblue, very thick] coordinates {(7,0) (10,0)};
% annotations
\node[engred, font=\scriptsize, anchor=south west] at (axis cs:3.1,50) {leading shock $+\Delta p$};
\node[engred, font=\scriptsize, anchor=north east] at (axis cs:6.9,-50) {trailing shock $-\Delta p$};
\draw[enggray, dashed] (axis cs:3,50) -- (axis cs:7,50);
\node[enggray, font=\scriptsize, anchor=south] at (axis cs:5,50) {$\Delta t \sim \ell/(M a)$};
\end{axis}
\end{tikzpicture}
\caption[The sonic-boom N-wave]{The idealised ground pressure
signature of a supersonic aircraft, the N-wave. As the aircraft's shock
system sweeps over an observer, the pressure jumps abruptly above
ambient at the leading (bow) shock, falls linearly through ambient over
the length of the aircraft's signature, and jumps back at the trailing
(tail) shock. The two shocks are heard as the characteristic
double-bang. The peak overpressure $\Delta p$ (of order $50$ to
$100\,\text{Pa}$ for a fighter at cruise) scales with the aircraft's
weight and volume and inversely with the three-quarter power of cruise
altitude; the duration $\Delta t$ between the shocks scales with the
aircraft length divided by the flight speed. Low-boom design reshapes
the longitudinal area and lift distribution to spread each shock into a
gradual rise, flattening the N into a lower-amplitude signature that a
noise-based rule could permit.}
\label{fig:vol03ch13:sonic-boom-nwave}
\end{figure}


\paragraph{Worked example 25: subsonic-compressible airspeed from a
pitot-static system.} A pitot-static system on a business jet at
$11\,\si{\kilo\meter}$ ($p_\infty = 22.7\,\text{kPa}$,
$T_\infty = 217\,\si{\kelvin}$) reads a differential (impact) pressure
$q_c = p_0 - p_\infty = 14.0\,\text{kPa}$. We want the Mach number and
true airspeed, using the compressible relation, and we compare against
the wrong answer the incompressible formula would give. The isentropic
stagnation relation inverts to the Mach number:
\[
\frac{p_0}{p_\infty} = 1 + \frac{q_c}{p_\infty} = 1 + \frac{14.0}{22.7}
= 1.617 = \left(1 + \frac{\gamma - 1}{2}M^2\right)^{\gamma/(\gamma -
1)}.
\]
Raising to the power $(\gamma - 1)/\gamma = 1/3.5 = 0.2857$,
$1 + 0.2 M^2 = 1.617^{0.2857} = 1.1470$, so $M^2 = 0.735$ and
$M = 0.857$. The speed of sound is $a = \sqrt{1.40\times 287\times 217}
= 295\,\si{\meter\per\second}$, so the true airspeed is
$u = M a = 0.857\times 295 = 253\,\si{\meter\per\second}$. The
incompressible reduction $q_c = \tfrac{1}{2}\rho u^2$ with
$\rho = p_\infty/(R T_\infty) = 22700/(287\times 217) =
0.364\,\si{\kilogram\per\meter\cubed}$ would give
$u = \sqrt{2 q_c/\rho} = \sqrt{2\times 14000/0.364} = 277\,\si{\meter
\per\second}$, an overstatement of $24\,\si{\meter\per\second}$, about
$9\,\%$. At $M = 0.86$ the incompressible formula reads high because it
ignores the density rise as the flow stagnates; every air-data computer
switches to the compressible relation above roughly $M = 0.3$ for
exactly this reason, and the correction grows steeply with Mach number,
becoming the Rayleigh-pitot correction of worked example~10 once the
probe sits behind a bow shock above $M = 1$.

\paragraph{Worked example 26: a total-temperature probe and the recovery
correction.} A total-temperature probe in the same $M = 0.86$,
$T_\infty = 217\,\si{\kelvin}$ stream reads a temperature, and we want to
recover the true stagnation temperature and the true static temperature
from it, given the probe's recovery factor $r = 0.96$ (a well-designed
shielded probe recovers most, not all, of the kinetic share). A probe
does not bring the flow fully to rest; it reads the recovery temperature
\[
T_r = T_\infty\left(1 + r\,\frac{\gamma - 1}{2}M^2\right)
= 217\left(1 + 0.96\times 0.2\times 0.7396\right)
= 217\times 1.1420 = 247.8\,\si{\kelvin}.
\]
The true stagnation temperature uses $r = 1$:
$T_0 = 217(1 + 0.2\times 0.7396) = 217\times 1.1479 =
249.1\,\si{\kelvin}$. The probe reads $247.8\,\si{\kelvin}$, low by
$1.3\,\si{\kelvin}$ against the true $249.1\,\si{\kelvin}$; the recovery
correction $T_0 = T_r + (1 - r)(T_0 - T_\infty)$ closes the gap. To
recover the static temperature, an engine-control computer needs the
Mach number too: from $T_0$ and the measured $M = 0.86$,
$T_\infty = T_0/(1 + 0.2 M^2) = 249.1/1.1479 = 217\,\si{\kelvin}$. The
lesson for engine control and air-data design is that a temperature
probe measures neither the static nor the true stagnation temperature
but the recovery temperature between them, and the small but real
recovery deficit $(1 - r)$ must be calibrated out before the reading is
used to schedule fuel flow or compute the corrected fan speed, where a
one-degree error shifts the surge margin.

\paragraph{Worked example 27: driver pressure to launch a shock-tube
shock.} A shock tube is to produce a normal shock of Mach number
$M_s = 2.4$ running into driven air at rest
($p_1 = 15\,\text{kPa}$, $T_1 = 293\,\si{\kelvin}$,
$a_1 = \sqrt{1.40\times 287\times 293} = 343\,\si{\meter\per\second}$,
$\gamma_1 = 1.40$), driven by high-pressure helium
($\gamma_4 = 1.667$, $a_4 = \sqrt{1.667\times 2077\times 293} =
1008\,\si{\meter\per\second}$ at the same fill temperature). We want the
diaphragm pressure ratio $p_4/p_1$ the driver must hold. The shock-tube
equation relates the driver-to-driven pressure ratio $p_4/p_1$ to the
shock strength $p_2/p_1$ through the requirement that the gas velocity
behind the shock match the velocity behind the expansion fan running
back into the driver:
\[
\frac{p_4}{p_1} = \frac{p_2}{p_1}\left[1 - \frac{(\gamma_4 - 1)(a_1/
a_4)(p_2/p_1 - 1)}{\sqrt{2\gamma_1\,[2\gamma_1 + (\gamma_1 + 1)
(p_2/p_1 - 1)]}}\right]^{-2\gamma_4/(\gamma_4 - 1)}.
\]
First the shock strength from $M_s = 2.4$:
$p_2/p_1 = 1 + \frac{2\gamma_1}{\gamma_1 + 1}(M_s^2 - 1) =
1 + \frac{2.8}{2.4}(5.76 - 1) = 1 + 1.1667\times 4.76 = 6.55$. The
bracket needs $a_1/a_4 = 343/1008 = 0.340$ and $\gamma_4 - 1 = 0.667$.
The numerator inside is $0.667\times 0.340\times(6.55 - 1) =
0.667\times 0.340\times 5.55 = 1.259$. The radical is
$\sqrt{2.8\,[2.8 + 2.4\times 5.55]} = \sqrt{2.8\times 16.12} =
\sqrt{45.1} = 6.72$, so the bracket is $[1 - 1.259/6.72]^{-2(1.667)/
0.667} = [1 - 0.1874]^{-5.0} = 0.8126^{-5.0} = 2.75$. Hence
$p_4/p_1 = 6.55\times 2.75 = 18.0$, and the driver pressure is
$p_4 = 18.0\times 15 = 270\,\text{kPa}$. Using helium rather than air as
the driver gas is the reason the required pressure ratio is only $18$
and not the far larger value air would demand: the light driver's high
sound speed lets its expansion fan supply the matching interface
velocity at a much lower pressure ratio, which is why shock tubes and
gun tunnels drive with helium or hydrogen to reach strong shocks without
extreme fill pressures.

\paragraph{Worked example 28: thermal-choking margin in a ramjet
combustor.} A ramjet flies at $M_\infty = 3.0$; its inlet delivers air
to the combustor at $M_1 = 0.25$, $T_1 = 700\,\si{\kelvin}$,
$T_{01} = 709\,\si{\kelvin}$ after the shock system decelerates and
heats the captured flow. The fuel is to add $q = 1.6\,\si{\mega\joule
\per\kilogram}$ of air. We want the combustor-exit Mach number and the
margin to thermal choking (worked example~8 introduced Rayleigh flow;
here we run it as a design margin). The heat raises the stagnation
temperature by $\Delta T_0 = q/c_p = 1.6\times 10^{6}/1005 =
1592\,\si{\kelvin}$, so $T_{02} = 709 + 1592 = 2301\,\si{\kelvin}$. The
choking limit is the sonic stagnation temperature $T_0^*$ referenced to
the inlet state. At $M_1 = 0.25$, the Rayleigh ratio is
\[
\frac{T_{01}}{T_0^*} = \frac{(\gamma + 1)M_1^2\,[2 + (\gamma - 1)M_1^2]}
{(1 + \gamma M_1^2)^2}
= \frac{2.4\times 0.0625\times(2 + 0.4\times 0.0625)}{(1 + 1.4\times
0.0625)^2},
\]
whose numerator is $0.15\times 2.025 = 0.3038$ and denominator is
$1.0875^2 = 1.1827$, giving $T_{01}/T_0^* = 0.2569$. Hence
$T_0^* = 709/0.2569 = 2760\,\si{\kelvin}$, the stagnation temperature at
which the exit would reach $M = 1$. The demanded exit stagnation
temperature $2301\,\si{\kelvin}$ sits below $2760\,\si{\kelvin}$, so the
flow does not choke, and the margin is
$T_0^* - T_{02} = 2760 - 2301 = 459\,\si{\kelvin}$, about
$q_{\text{choke}} = c_p(T_0^* - T_{01}) = 1005\times(2760 - 709) =
2.06\,\si{\mega\joule\per\kilogram}$ of heat capacity in reserve. The
combustor-exit Mach number follows by finding the $M_2$ whose Rayleigh
ratio is $T_{02}/T_0^* = 2301/2760 = 0.834$; inverting the Rayleigh
relation on the subsonic branch gives $M_2 = 0.49$. The design lesson is
that a ramjet combustor must hold a positive thermal-choking margin at
every flight condition, and the margin shrinks as the flight Mach number
rises (which raises the combustor-inlet stagnation temperature and eats
the reserve), so the maximum fuel-air ratio, and with it the thrust, is
capped by thermal choking well before it is capped by stoichiometry.

\paragraph{Worked example 29: the critical Mach number of a wing
section.} A thin transport section has an incompressible suction-peak
pressure coefficient $C_{p,0} = -0.43$. We want the critical Mach
number, the freestream Mach at which the flow first reaches $M = 1$ at
that peak, because it sets the floor of the drag-divergence band the
cruise point must clear. The construction pairs the
compressibility-corrected peak against the critical pressure
coefficient $C_p^*$ and reads off the crossing
(figure~\ref{fig:vol03ch13:critical-mach}). The Prandtl-Glauert-corrected
peak is $C_p = C_{p,0}/\sqrt{1 - M_\infty^2}$, growing in magnitude as
$M_\infty$ rises. The critical pressure coefficient, the $C_p$ at which
the local Mach number is exactly one, is
\[
C_p^* = \frac{2}{\gamma M_\infty^2}\left[\left(\frac{1 + 0.2
M_\infty^2}{1.2}\right)^{3.5} - 1\right].
\]
We march $M_\infty$ upward and compare the two. At $M_\infty = 0.65$ the
corrected peak is $-0.43/\sqrt{1 - 0.4225} = -0.566$ while
$C_p^* = -0.784$: the peak has not yet reached sonic, so the section is
subcritical. At $M_\infty = 0.75$ the corrected peak is
$-0.43/\sqrt{1 - 0.5625} = -0.650$ while $C_p^* = -0.564$: the peak has
overshot sonic, so the section is now supercritical. The crossing sits
between, and interpolating (or reading
figure~\ref{fig:vol03ch13:critical-mach}) gives
$M_{\mathrm{cr}} \approx 0.71$, where both coefficients equal about
$-0.61$. Above $M_{\mathrm{cr}}$ a supersonic pocket grows on the upper
surface and closes in a shock, and the drag-divergence Mach number sits
a few hundredths of Mach above $M_{\mathrm{cr}}$. The design lever is
the suction peak: a shallower $C_{p,0}$ shifts the corrected-peak curve
down, so it meets $C_p^*$ further right, and the section flies faster
before it goes critical. That is exactly what the supercritical section
of section~13.6 buys by flattening the upper surface.

% Figure: graphical determination of the critical Mach number. Two
% curves are drawn against freestream Mach number: the
% compressibility-corrected suction-peak pressure coefficient (rising in
% magnitude as M grows, Prandtl-Glauert), and the critical pressure
% coefficient Cp* (the Cp at which the local Mach number is exactly 1).
% Their intersection is M_cr. Cp is plotted as -Cp so both curves are
% positive and the crossing reads clean.
% Prandtl-Glauert: Cp = Cp0 / sqrt(1-M^2), here Cp0 = -0.43 (a thin
% section), so -Cp = 0.43/sqrt(1-M^2).
% Cp*: Cp* = (2/(g M^2)) [ ((1 + 0.2 M^2)/1.2)^3.5 - 1 ], g=1.4;
% -Cp* is the negative of a negative number, i.e. |Cp*|, plotted directly.
\begin{figure}[htbp]
\centering
\begin{tikzpicture}
\begin{axis}[
  width=11cm, height=8cm,
  xlabel={freestream Mach number $M_\infty$},
  ylabel={pressure coefficient magnitude $|C_p|$},
  xmin=0.4, xmax=0.9, ymin=0, ymax=1.6,
  axis lines=left,
  grid=major, grid style={enggray!20},
  tick label style={font=\scriptsize},
  label style={font=\small},
  legend style={font=\scriptsize, at={(0.03,0.97)}, anchor=north west},
  legend cell align=left,
]
% Corrected suction-peak Cp magnitude: 0.43 / sqrt(1-M^2)
\addplot[engblue, very thick, domain=0.40:0.86, samples=120]
  { 0.43 / sqrt(1 - x*x) };
\addlegendentry{corrected suction peak $|C_{p}|=|C_{p,0}|/\sqrt{1-M_\infty^2}$}
% Critical pressure coefficient magnitude:
% Cp* = (2/(1.4 M^2)) [ ((1+0.2 M^2)/1.2)^3.5 - 1 ], which is negative;
% plot its magnitude = -Cp*.
\addplot[engred, very thick, domain=0.40:0.88, samples=120]
  { -( (2/(1.4*x*x)) * ( ((1 + 0.2*x*x)/1.2)^3.5 - 1 ) ) };
\addlegendentry{critical $|C_p^{*}|$ (local $M=1$)}
% Intersection marker near M_cr = 0.71 (from the worked example)
\fill[engblack] (axis cs:0.71,0.611) circle (1.8pt);
\draw[enggray, dashed] (axis cs:0.71,0) -- (axis cs:0.71,0.611);
\node[engblack, font=\scriptsize, anchor=west] at (axis cs:0.72,0.75)
  {$M_{\mathrm{cr}}\approx 0.71$};
\end{axis}
\end{tikzpicture}
\caption[Graphical construction of the critical Mach number]{The
critical Mach number is where the compressibility-corrected
suction-peak pressure coefficient (blue) first reaches the critical
value $C_p^{*}$ (red), the pressure coefficient at which the local Mach
number equals one. Both are plotted as magnitudes so the crossing reads
directly. The blue curve rises with $M_\infty$ through the
Prandtl-Glauert factor $1/\sqrt{1-M_\infty^2}$; the red curve falls as
$M_\infty$ climbs, because a faster freestream needs a smaller pressure
drop to reach sonic conditions. For the thin section drawn here
($|C_{p,0}|=0.43$ at the suction peak) the two meet at
$M_{\mathrm{cr}}\approx 0.71$. A section with a shallower suction peak
shifts the blue curve down, so it meets the red curve further to the
right, raising the critical Mach number and, with it, the
drag-divergence Mach number. Flattening the upper-surface pressure is
what a supercritical section does to push this crossing toward unity.}
\label{fig:vol03ch13:critical-mach}
\end{figure}


\paragraph{Worked example 30: the Sears-Haack minimum wave drag of a
fuselage.} A supersonic transport fuselage is $\ell = 60\,\si{\meter}$
long and encloses a volume $V = 400\,\si{\meter}^3$. We want the
lower-bound volume wave drag the Sears-Haack body of the same length and
volume would carry at $M = 2.0$, $18\,\si{\kilo\meter}$ altitude
($\rho_\infty = 0.12\,\si{\kilogram\per\meter\cubed}$,
$a_\infty = 295\,\si{\meter\per\second}$,
$U_\infty = 590\,\si{\meter\per\second}$), because that bound is what a
real fuselage is measured against. The Sears-Haack volume wave drag is
\[
D_w = \frac{9\pi^3}{8}\,\frac{V^2}{\ell^4}\cdot\tfrac{1}{2}\rho_\infty
U_\infty^2.
\]
The dynamic pressure is
$q_\infty = \tfrac{1}{2}\times 0.12\times 590^2 =
2.09\times 10^{4}\,\text{Pa}$. The geometric group is
$V^2/\ell^4 = 400^2/60^4 = 1.6\times 10^{5}/1.296\times 10^{7} =
0.01235\,\si{\meter}^2$, and $9\pi^3/8 = 9\times 31.01/8 = 34.9$. Hence
\[
D_w = 34.9\times 0.01235\times 2.09\times 10^{4} = 9.0\times 10^{3}\,
\si{\newton},
\]
about $9\,\text{kN}$. The wave-drag \emph{area} (the drag divided by
$q_\infty$) is $D_w/q_\infty = 0.43\,\si{\meter}^2$, a compact way to
compare shapes independent of flight condition. The steep penalties are
in the exponents: the wave drag falls as the fourth power of length, so
stretching the same volume into a longer, more slender body is the
strongest lever a designer has (doubling the length cuts the wave drag
sixteenfold), while it rises as the square of volume, so every cubic
metre of enclosed volume is expensive at supersonic speed. A real
fuselage carries more than the Sears-Haack bound because it cannot be
pointed at both ends and must clear a cabin cross-section, but the bound
is the target, and the area rule of section~13.6 is the statement that
the wing and fuselage must add their areas into a single smooth
distribution to approach it.

\paragraph{Worked example 31: sweep raises the drag-divergence Mach
number.} A wing section has an unswept drag-divergence Mach number
$M_{DD,0} = 0.74$. We want the freestream drag-divergence Mach number
when the wing is swept at quarter-chord angles $\Lambda = 25^{\circ}$
and $35^{\circ}$, and the cruise-Mach gain sweep buys
(figure~\ref{fig:vol03ch13:swept-wing}). The shock that sets drag
divergence responds to the Mach component normal to the wing leading
edge, $M_n = M_\infty\cos\Lambda$, so the section reaches its limit when
$M_\infty\cos\Lambda = M_{DD,0}$, that is
$M_{DD} = M_{DD,0}/\cos\Lambda$. At $\Lambda = 25^{\circ}$,
$\cos 25^{\circ} = 0.906$, so $M_{DD} = 0.74/0.906 = 0.817$. At
$\Lambda = 35^{\circ}$, $\cos 35^{\circ} = 0.819$, so
$M_{DD} = 0.74/0.819 = 0.904$. Sweeping the wing $35^{\circ}$ lifts the
usable drag-divergence Mach number from $0.74$ to $0.90$, a gain of
$0.16$ that turns a section unusable at transport cruise into one that
clears $M = 0.85$ with margin. The estimate is optimistic: a finite
swept wing recovers only about two-thirds of the ideal
$1/\cos\Lambda$ benefit, because the flow near the root and tip is not
the clean two-dimensional swept flow the rule assumes, so a
$35^{\circ}$ wing behaves more like $M_{DD} \approx 0.85$ than $0.90$.
Sweep also costs weight, root-bending structure, and low-speed lift, so
transport wings settle near $25^{\circ}$ to $35^{\circ}$ and pair the
sweep with a supercritical section (worked example~29) rather than
leaning on either alone. The two levers, sweep on the planform and
suction-peak flattening on the section, are the reason a modern
transport cruises at $M = 0.85$ on a wing whose bare unswept section
would diverge at $0.74$.

% Figure: the simple-sweep effect on the drag-divergence Mach number.
% The component of the freestream Mach number normal to a wing swept at
% angle Lambda is M_n = M_infty cos(Lambda), and the shock physics that
% sets drag divergence responds to M_n, so the aircraft can fly to a
% higher M_infty before the normal component reaches the section's
% drag-divergence value. Plotting the usable freestream drag-divergence
% Mach M_DD(Lambda) = M_DD0 / cos(Lambda) against sweep angle shows the
% gain, and a horizontal reference marks the unswept section value.
% M_DD0 = 0.74 (a conventional unswept section).
\begin{figure}[htbp]
\centering
\begin{tikzpicture}
\begin{axis}[
  width=11cm, height=8cm,
  xlabel={quarter-chord sweep angle $\Lambda$ (degrees)},
  ylabel={freestream drag-divergence Mach $M_{DD}$},
  xmin=0, xmax=45, ymin=0.7, ymax=1.1,
  axis lines=left,
  grid=major, grid style={enggray!20},
  tick label style={font=\scriptsize},
  label style={font=\small},
  legend style={font=\scriptsize, at={(0.03,0.97)}, anchor=north west},
  legend cell align=left,
]
% Simple-sweep estimate: M_DD = 0.74 / cos(Lambda)
\addplot[engblue, very thick, domain=0:42, samples=100]
  { 0.74 / cos(x) };
\addlegendentry{simple-sweep estimate $M_{DD}=M_{DD,0}/\cos\Lambda$}
% Unswept reference
\addplot[enggray, thick, dashed, domain=0:45, samples=2] { 0.74 };
\addlegendentry{unswept section $M_{DD,0}=0.74$}
% marker at 30 deg: 0.74/cos30 = 0.74/0.866 = 0.854
\fill[engblack] (axis cs:30,0.854) circle (1.8pt);
\draw[enggray, dashed] (axis cs:30,0.7) -- (axis cs:30,0.854);
\node[engblack, font=\scriptsize, anchor=south east] at (axis cs:29.5,0.86)
  {$\Lambda=30^{\circ}$: $M_{DD}\approx 0.85$};
\end{axis}
\end{tikzpicture}
\caption[Sweep raises the drag-divergence Mach number]{The simple-sweep
estimate of the usable freestream drag-divergence Mach number as a
function of quarter-chord sweep $\Lambda$, for a section whose unswept
value is $M_{DD,0}=0.74$ (grey dashed). Because the shock that sets drag
divergence responds to the Mach component normal to the wing,
$M_n=M_\infty\cos\Lambda$, the aircraft reaches a higher freestream Mach
number before that normal component hits the section limit, and the
usable $M_{DD}$ rises as $M_{DD,0}/\cos\Lambda$. At $\Lambda=30^{\circ}$
the estimate gives $M_{DD}\approx 0.85$, a gain of about $0.11$ in
cruise Mach over the unswept wing. The estimate is optimistic (real
wings recover roughly two-thirds of the ideal $1/\cos\Lambda$ benefit
because the wing is finite and the flow is not purely two-dimensional),
but it captures why every high-subsonic transport carries a swept wing
and why sweep and the supercritical section are combined to place cruise
below drag divergence.}
\label{fig:vol03ch13:swept-wing}
\end{figure}


\paragraph{Worked example 32: Fanno flow that does not quite choke.}
Worked example~7 found the pipe length that just chokes an entry flow.
Here the pipe is given and shorter than the choking length, and we want
the exit state and the stagnation-pressure lost to friction. Air enters
a smooth insulated pipe of diameter $D = 40\,\si{\milli\meter}$,
length $L = 6.0\,\si{\meter}$, friction factor $f = 0.005$ (Darcy
$4f = 0.02$), at $M_1 = 0.20$. The dimensionless friction the pipe
imposes is $4fL/D = 0.02\times 6.0/0.040 = 3.0$. The Fanno function at
the inlet, the friction length from $M_1$ to sonic, is
\[
\frac{4fL^*}{D}\bigg|_{M_1} = \frac{1 - M_1^2}{\gamma M_1^2}
+ \frac{\gamma + 1}{2\gamma}\ln\!\left[\frac{(\gamma + 1)M_1^2}
{2 + (\gamma - 1)M_1^2}\right],
\]
which at $M_1 = 0.20$ evaluates to
$(1 - 0.04)/(1.4\times 0.04) + 0.857\ln[(2.4\times 0.04)/(2 + 0.4\times
0.04)] = 17.14 + 0.857\ln(0.096/2.016) = 17.14 - 2.61 = 14.53$. The
pipe consumes $3.0$ of this, leaving
$4fL^*/D|_{M_2} = 14.53 - 3.0 = 11.53$ of friction length still to
sonic at the exit. Inverting the Fanno function for the Mach number
whose remaining friction length is $11.53$ gives $M_2 = 0.226$: the
flow speeds up from $M = 0.20$ to $M = 0.226$, still comfortably
subsonic, so the pipe does not choke. The stagnation-pressure ratio
along a Fanno duct is $p_{0}/p_0^*$ at each Mach number, with
\[
\frac{p_0}{p_0^*} = \frac{1}{M}\left[\frac{2 + (\gamma - 1)M^2}
{\gamma + 1}\right]^{(\gamma + 1)/[2(\gamma - 1)]},
\]
which gives $p_{01}/p_0^* = 2.964$ at $M_1 = 0.20$ and
$p_{02}/p_0^* = 2.632$ at $M_2 = 0.226$, so the stagnation-pressure
recovery is $p_{02}/p_{01} = 2.632/2.964 = 0.888$. Friction alone,
with no heat and no shock, has thrown away $11\,\%$ of the stagnation
pressure over six metres. The lesson for gas distribution is that a
subsonic Fanno line loses stagnation pressure even while it stays far
from choking, and the loss is read off the Fanno function without
tracking the pressure profile point by point.

\paragraph{Worked example 33: when a wall reflection turns into a Mach
stem.} Worked example~22 found a clean regular reflection. Here we push
the incident shock until regular reflection can no longer close, to
locate the transition. A $\theta = 20^{\circ}$ wedge sits in an
$M_1 = 2.8$ stream and its shock strikes a parallel wall. Does the
reflection stay regular? The incident shock at $M_1 = 2.8$,
$\theta = 20^{\circ}$ has weak-shock angle $\beta_1 = 40.5^{\circ}$
(\texttt{code/oblique\_shock.py}), so
$M_{n1} = 2.8\sin 40.5^{\circ} = 1.819$ and the deflected Mach number
behind the incident shock is
$M_2 = M_{n2}/\sin(\beta_1 - \theta) = 0.612/\sin 20.5^{\circ} = 1.75$.
For regular reflection the wall must turn this $M_2 = 1.75$ flow back
through the same $20^{\circ}$ with a single reflected shock, so the
requirement is $\theta = 20^{\circ} < \theta_{\max}(M_2 = 1.75)$. The
maximum deflection at $M = 1.75$ is only
$\theta_{\max} \approx 19.3^{\circ}$ (\texttt{code/oblique\_shock.py},
scanning $\beta$), which is below the required $20^{\circ}$. No
reflected shock can hold the flow to the wall, so regular reflection
fails and a Mach stem forms: the incident and reflected shocks meet at a
triple point above the wall, a nearly normal Mach stem bridges from
there to the surface, and a slip line trails downstream. The transition
was crossed by asking one number, whether the post-incident flow can be
turned through the wall angle by any attached shock. Had the wedge been
$\theta = 10^{\circ}$ (worked example~22), the post-incident
$M_2 = 2.22$ carried a $\theta_{\max} \approx 26^{\circ}$ with room to
spare, and the reflection stayed regular. The engineering point is that
the reflection pattern flips discontinuously as the incident angle or
Mach number crosses this line, and a shock reflecting inside an inlet
can jump from a clean regular reflection to a lossy Mach stem with a
strong normal-shock core over a small change in operating condition.

\paragraph{Worked example 34: Allen-Eggers peak deceleration on a
ballistic entry.} The blunt-body case study estimated the peak heat
flux; here we estimate the peak deceleration a ballistic capsule feels,
because the structure and the crew must survive it, and the classical
Allen-Eggers analysis gives it in closed form. A capsule enters an
exponential atmosphere (density $\rho = \rho_0 e^{-z/H}$ with scale
height $H = 7.2\,\si{\kilo\meter}$ and sea-level
$\rho_0 = 1.22\,\si{\kilogram\per\meter\cubed}$) at speed
$U_e = 7.5\,\si{\kilo\meter\per\second}$ and flight-path angle
$\gamma_e = 6^{\circ}$ below the horizontal, with ballistic coefficient
$B = m/(C_D A) = 400\,\si{\kilogram\per\meter\squared}$. The
Allen-Eggers result, from integrating the drag equation for a
non-lifting body at constant flight-path angle, gives the peak
deceleration
\[
a_{\max} = \frac{U_e^2\,\sin\gamma_e}{2 e\,H},
\]
independent of the ballistic coefficient, and reached at the altitude
where $\rho = 2 B\sin\gamma_e/H$. Numerically
$U_e^2\sin\gamma_e = (7500)^2\times\sin 6^{\circ} =
5.625\times 10^{7}\times 0.1045 = 5.88\times 10^{6}$, and the
denominator is $2 e H = 2\times 2.718\times 7200 =
3.91\times 10^{4}\,\si{\meter}$, so
$a_{\max} = 5.88\times 10^{6}/3.91\times 10^{4} =
150\,\si{\meter\per\second\squared}$, about $15\,g$. The altitude of
peak deceleration is where
$\rho = 2\times 400\times 0.1045/7200 = 0.0116\,\si{\kilogram\per
\meter\cubed}$, which the exponential atmosphere places at
$z = H\ln(\rho_0/\rho) = 7200\ln(1.22/0.0116) = 7200\times 4.65 =
33\,\si{\kilo\meter}$. The result is worth reading: the peak
deceleration depends only on the entry speed, the entry angle, and the
atmospheric scale height, not on the vehicle's mass or drag area, so a
steeper entry is felt more sharply regardless of how the capsule is
built. A $6^{\circ}$ path gives a survivable $15\,g$; a $30^{\circ}$
steep entry would raise $\sin\gamma_e$ fivefold to about $70\,g$, beyond
crew tolerance, which is why crewed capsules fly shallow entry
corridors, trading a longer, hotter descent for a survivable load. The
heating of the case study and the deceleration here are the two
competing constraints that define the entry corridor.

\paragraph{Worked example 35: the transonic drag rise scaled by
similarity.} A wing section shows a wave-drag increment
$\Delta C_d = 0.0020$ at $M_\infty = 0.80$ on a thickness ratio
$t/c = 0.12$. We want the wave-drag increment at the same Mach number
for a thinner $t/c = 0.09$ section of the same family, using transonic
small-disturbance similarity, because thinning is the direct lever on
transonic drag. Transonic similarity states that the scaled wave-drag
coefficient depends only on the transonic similarity parameter
\[
K = \frac{M_\infty^2 - 1}{[(\gamma + 1)M_\infty^2\,(t/c)]^{2/3}},
\]
through $C_d/(t/c)^{5/3} = F(K)$ for a fixed section family, where $F$
is a universal function. Two sections at the same $M_\infty$ but
different $t/c$ do not share $K$, so a clean scaling needs the
freestream Mach numbers matched to hold $K$ equal; at a fixed
$M_\infty$ the more useful working scaling is the leading-order
thickness dependence $\Delta C_d \propto (t/c)^{5/3}$ that the
similarity form carries when the sections are near the same point on
$F(K)$. Applying it, $\Delta C_d = 0.0020\times
(0.09/0.12)^{5/3} = 0.0020\times 0.75^{1.667} = 0.0020\times 0.618 =
0.00124$. Thinning the section from $12\,\%$ to $9\,\%$ cuts the wave
drag increment from $0.0020$ to $0.0012$, a $38\,\%$ reduction, far more
than the $25\,\%$ thickness cut alone would suggest, because the wave
drag responds to thickness with the $5/3$ power. The number quantifies
the thin-section instinct of supersonic design and the diamond-airfoil
lesson of worked example~23: thickness is expensive in wave drag, and
the exponent is steep. The caution is that thinning trades wave drag for
structural weight and internal volume, so the section thickness is set
by the balance of the two, not by wave drag alone.

\section{Vignette: sizing a blowdown supersonic wind tunnel}\pathtag{standard}

The case studies engineer a flow inside a duct or around a body. This
shorter vignette sizes a piece of test equipment, the
\engterm{blowdown wind tunnel} that runs much of the world's
supersonic testing, and it ties the choked-mass-flow and run-time
arithmetic to a tank a reader could buy. A blowdown tunnel stores
high-pressure air in a tank, releases it through a settling chamber and
a converging-diverging nozzle to a supersonic test section, and runs
until the tank pressure falls too low to hold the test-section Mach
number. It trades a short run for a simple, cheap, high-Mach facility,
the opposite trade from a continuous closed-circuit tunnel.

\paragraph{Requirement.} We want a test section running $M = 2.5$ air
with a $0.1\,\si{\meter}\times 0.1\,\si{\meter}$ square working area,
fed from a tank charged to $20$ bar, and we want the run time the
facility delivers per charge. The nozzle must hold the test-section
Mach number, which requires the settling-chamber (stagnation) pressure
to stay above the value that keeps the nozzle started and the test
section at $M = 2.5$.

\paragraph{Step 1: the test-section and choking conditions.} At
$M = 2.5$ the area ratio from the area-Mach relation is
$A_e/A^* = 2.637$ (\texttt{code/nozzle.py}). The test-section area is
$A_e = 0.1\times 0.1 = 0.01\,\si{\meter}^2$, so the throat area is
$A^* = A_e/2.637 = 3.79\times 10^{-3}\,\si{\meter}^2$. The pressure
ratio needed to run $M = 2.5$ isentropically is
$p_0/p_e = (1 + 0.2\times 2.5^2)^{3.5} = 2.25^{3.5} = 17.1$, and the
diffuser at the test-section exit recovers some pressure, but a working
rule is that the settling-chamber pressure must stay above roughly
$0.9$ bar of test-section static plus the diffuser loss; for a vented
(atmospheric-exhaust) tunnel the floor is set by the normal-shock
pressure recovery of the starting shock, about $p_{0,\min}\approx
8\,\text{bar}$ for this configuration
(figure~\ref{fig:vol03ch13:blowdown}).

% Figure: blowdown wind-tunnel reservoir pressure decay. A fixed-volume
% high-pressure tank feeds a choked test section; the reservoir pressure
% falls exponentially as mass leaves at the choked rate, and the usable
% run ends when the pressure can no longer sustain the design test-section
% Mach number (the minimum-pressure floor). Schematic with a representative
% decay curve and the run-time window marked.
\begin{figure}[htbp]
\centering
\begin{tikzpicture}
\begin{axis}[
  width=12cm, height=8cm,
  xlabel={time $t$ (s)},
  ylabel={reservoir pressure $p_0(t)$ (bar)},
  xmin=0, xmax=60, ymin=0, ymax=22,
  axis lines=left,
  grid=major, grid style={enggray!20},
  tick label style={font=\scriptsize},
  label style={font=\small},
  legend style={font=\scriptsize, at={(0.97,0.97)}, anchor=north east},
]
% exponential decay p0(t) = p0_init * exp(-t/tau), tau = 25 s
\addplot[engblue, very thick, domain=0:60, samples=200] {
  20*exp(-x/25) };
\addlegendentry{$p_0(t)=p_{0,i}\,e^{-t/\tau}$}
% minimum-pressure floor to hold the design Mach number: 8 bar
\addplot[engred, dashed, domain=0:60, samples=2] {8};
\node[engred, font=\scriptsize, anchor=south west] at (axis cs:0.5,8.1) {minimum $p_0$ to hold design $M$};
% run ends where curve crosses the floor: 20 exp(-t/25)=8 -> t = 25 ln(2.5) = 22.9 s
\addplot[engblack, only marks, mark=*, mark options={fill=engblack}] coordinates {(22.9,8)};
\node[engblack, font=\scriptsize, anchor=south west] at (axis cs:23.5,8.6) {run ends, $t\approx 23$ s};
% shade the usable run window
\draw[englime, very thick] (axis cs:0,1.0) -- (axis cs:22.9,1.0);
\node[englime, font=\scriptsize, anchor=south] at (axis cs:11.5,1.1) {usable run};
\end{axis}
\end{tikzpicture}
\caption[Blowdown reservoir pressure decay]{Reservoir pressure of a
fixed-volume blowdown wind tunnel feeding a choked test section. With
the throat choked, mass leaves at a rate proportional to the
instantaneous reservoir pressure, so for an isothermal reservoir the
pressure decays exponentially, $p_0(t)=p_{0,i}\,e^{-t/\tau}$, with the
time constant $\tau$ set by the tank volume divided by the choked
volumetric efflux. The run ends not when the tank empties but when the
reservoir pressure falls below the floor needed to keep the test-section
pressure ratio above its design value (here $8$ bar from an initial
$20$ bar), which fixes the usable run time at $t=\tau\ln(p_{0,i}/p_{0,
\min})\approx 23\,$s for the curve shown. The blowdown vignette in the
text carries this arithmetic for a sized tank.}
\label{fig:vol03ch13:blowdown}
\end{figure}


\paragraph{Step 2: the choked mass flow and the tank drain.} With the
throat choked, the mass flow at settling-chamber conditions
$p_0$, $T_0 = 290\,\si{\kelvin}$ is
\[
\dot m = A^* \frac{p_0}{\sqrt{T_0}}\sqrt{\frac{\gamma}{R}}
\left(\frac{2}{\gamma + 1}\right)^{(\gamma + 1)/[2(\gamma - 1)]}
= 3.79\times 10^{-3}\times \frac{p_0}{\sqrt{290}}\times 0.0404
\sqrt{287},
\]
which evaluates to $\dot m = 0.0404\times 3.79\times 10^{-3}
\times 16.94 \times p_0/\sqrt{290}$. At the initial $p_0 = 20\,\text{bar}
= 2.0\times 10^6\,\text{Pa}$ this is $\dot m = 2.59\times 10^{-3}
\times 2.0\times 10^6/17.03 = 3.05\,\si{\kilogram\per\second}$. The
mass flow is proportional to $p_0$, so as the tank drains the flow
slows, and the tank pressure decays. Treating the tank as a
fixed-volume, roughly isothermal reservoir of volume $V$, mass
conservation $V\,d\rho_t/dt = -\dot m$ with $\rho_t = p_0/(R T_0)$ and
$\dot m \propto p_0$ gives $dp_0/dt = -(p_0/\tau)$, an exponential
decay $p_0(t) = p_{0,i}e^{-t/\tau}$ with time constant $\tau = V/(C
\sqrt{T_0})$, where $C$ collects the choked-flow constants.

\paragraph{Step 3: the run time.} For a tank of $V = 2\,\si{\meter}^3$
the time constant works out to $\tau \approx 25\,\si{\second}$ at these
conditions. The run ends when $p_0$ falls to the floor
$p_{0,\min} = 8\,\text{bar}$, so the usable run is
\[
t_{\text{run}} = \tau\ln\frac{p_{0,i}}{p_{0,\min}} =
25\ln\frac{20}{8} = 25\times 0.916 = 23\,\si{\second}.
\]
The facility delivers about $23\,\si{\second}$ of $M = 2.5$ flow per
charge (figure~\ref{fig:vol03ch13:blowdown}), enough for a
pressure-tap survey or a schlieren sequence but not for a long
thermal-soak test. The arithmetic shows the two levers a blowdown
tunnel offers: a larger tank lengthens $\tau$ and the run linearly,
and a higher charge pressure lengthens the run only logarithmically,
so doubling the tank doubles the run while doubling the charge adds
only one more time constant of margin above the floor. The blowdown
architecture is the cheap way to reach high Mach in short bursts, and
the choked-mass-flow relation of this chapter, run against a drainage
balance, sizes it end to end.

\section{Vignette: the supersonic ejector as a jet pump}\pathtag{standard}

The blowdown tunnel spent stored pressure to reach high Mach. This
shorter vignette sizes a device that spends a flowing high-pressure
stream to move a second, low-pressure one, the \engterm{supersonic
ejector} or jet pump (figure~\ref{fig:vol03ch13:ejector}). An ejector
has no moving parts: a primary stream expands through a small
converging-diverging nozzle to a supersonic jet, the jet entrains a
secondary stream by turbulent shear, the two mix in a constant-area
duct, and the mixture recompresses through a diffuser to a delivery
pressure above the secondary inlet pressure. The ejector pulls the
vacuum in most intermittent supersonic tunnels, evacuates steam
condensers, and boosts low-pressure gas lines, and its performance
follows from the same choked-flow and momentum relations the chapter has
built.

% Figure: supersonic ejector (jet pump). A high-pressure primary stream
% expands through a converging-diverging nozzle to supersonic speed,
% entrains a low-pressure secondary stream, mixes in a constant-area
% duct, and recompresses through a diffuser and a mixing shock train.
\begin{figure}[htbp]
\centering
\begin{tikzpicture}[scale=1.0]
% outer casing (secondary inlet at top and bottom left, mixing tube, diffuser)
\draw[engblack, very thick] (0,1.3) -- (4.5,1.3) -- (7.5,1.9);
\draw[engblack, very thick] (0,-1.3) -- (4.5,-1.3) -- (7.5,-1.9);
% primary CD nozzle inside, centred
\draw[engblue, thick] (0.2,0.55) -- (1.2,0.28) -- (2.0,0.5);
\draw[engblue, thick] (0.2,-0.55) -- (1.2,-0.28) -- (2.0,-0.5);
\node[engblue, font=\scriptsize, anchor=south] at (1.1,0.6) {primary nozzle};
% primary arrow
\draw[engblue, -{Latex[length=1.8mm]}] (-0.5,0) -- (0.15,0);
\node[engblue, font=\scriptsize, anchor=east] at (-0.5,0) {$p_{0p}$ (high)};
% secondary entrainment arrows
\draw[engamber, -{Latex[length=1.5mm]}] (-0.3,0.95) -- (1.6,0.85);
\draw[engamber, -{Latex[length=1.5mm]}] (-0.3,-0.95) -- (1.6,-0.85);
\node[engamber, font=\scriptsize, anchor=east] at (-0.3,0.95) {$p_{0s}$ (low)};
% supersonic primary jet leaving the nozzle
\draw[engred, -{Latex[length=1.8mm]}] (2.1,0) -- (3.4,0);
\node[engred, font=\scriptsize, anchor=south] at (2.7,0.05) {supersonic jet};
% mixing region label
\node[font=\scriptsize, anchor=north] at (3.5,-1.3) {constant-area mixing};
% mixing shock train (a few oblique marks)
\foreach \x in {3.6,4.0,4.4} {
  \draw[engred] (\x,0.5) -- (\x,-0.5);
}
% diffuser + exit
\draw[enggray, -{Latex[length=1.8mm]}] (6.0,0) -- (7.2,0);
\node[enggray, font=\scriptsize, anchor=west] at (7.2,0) {to exhaust};
\node[font=\scriptsize, anchor=north] at (6.3,-1.6) {diffuser};
\end{tikzpicture}
\caption[The supersonic ejector (jet pump)]{A supersonic ejector, or
jet pump. A high-pressure primary stream expands through a small
converging-diverging nozzle to a supersonic jet. The low-pressure jet
entrains a secondary stream by shear, dragging it into a constant-area
mixing tube where the two streams exchange momentum. The mixed flow,
still carrying the primary jet's kinetic energy, recompresses through a
diffuser (and, in the supersonic-mixing regime, a train of oblique
shocks) to a delivery pressure above the secondary inlet pressure. The
device has no moving parts and pumps the secondary stream up in
pressure at the cost of the primary stream's stagnation pressure, which
makes it the standard way to pull the vacuum in a supersonic wind
tunnel, to evacuate a condenser, or to boost a low-pressure gas line.
Its performance is read off the same choked-mass-flow, mixing-momentum,
and shock-recovery relations the rest of the chapter develops.}
\label{fig:vol03ch13:ejector}
\end{figure}


\paragraph{Requirement.} We want an ejector that entrains a secondary
air stream at $p_{0s} = 20\,\text{kPa}$ (a partial vacuum to be pulled)
using a primary air stream at $p_{0p} = 600\,\text{kPa}$,
$T_{0p} = 300\,\si{\kelvin}$, and we want the entrainment ratio (the
secondary-to-primary mass-flow ratio) the device can sustain and the
primary nozzle throat area to move a chosen secondary flow.

\paragraph{Step 1: the primary jet.} The primary nozzle expands from
$600\,\text{kPa}$ toward the low mixing-chamber pressure. To match the
secondary pressure near $20\,\text{kPa}$ the primary exit runs at a
pressure ratio $p_{0p}/p_e \approx 600/20 = 30$, which the isentropic
relation places at $1 + 0.2 M_e^2 = 30^{1/3.5} = 2.64$, so
$M_e^2 = 8.2$ and $M_e = 2.86$. The primary jet leaves the nozzle at
nearly Mach $2.9$, and its area ratio from the area-Mach relation is
$A_e/A^* = 3.85$ (\texttt{code/nozzle.py}). This supersonic jet at low
static pressure is what entrains the secondary flow.

\paragraph{Step 2: the entrainment ratio from momentum.} The mixing tube
conserves mass and momentum. Writing the entrainment ratio
$\omega = \dot m_s/\dot m_p$, an ideal constant-area mixing analysis with
both streams brought to a common mixed state gives an entrainment ratio
that rises as the primary stagnation pressure rises relative to the
secondary and falls as the required compression (delivery-to-secondary
pressure ratio) rises. For this pressure combination a representative
result is $\omega \approx 0.5$: the ejector moves about half a kilogram
of secondary air for each kilogram of primary air, delivering the
mixture at a compression ratio of roughly $1.6$ to $2.0$ above the
secondary inlet. The trade is fixed by the two pressure ratios: a higher
primary pressure buys a higher entrainment or a higher compression, but
not both at once, and the product of entrainment and compression is
bounded by the primary stream's available stagnation-pressure drop.

\paragraph{Step 3: sizing the throat.} Suppose the secondary flow to be
pumped is $\dot m_s = 0.10\,\si{\kilogram\per\second}$. At
$\omega = 0.5$ the primary flow is $\dot m_p = \dot m_s/\omega =
0.20\,\si{\kilogram\per\second}$. The primary nozzle is choked, so its
throat area follows from the choked mass flux at
$p_{0p} = 600\,\text{kPa}$, $T_{0p} = 300\,\si{\kelvin}$:
$\dot m_p/A^* = (p_{0p}/\sqrt{T_{0p}})\times 0.0404\sqrt{287} =
(6\times 10^{5}/17.3)\times 0.684 = 2.37\times 10^{4}\,\si{\kilogram\per
\second}/\si{\meter}^2$, so $A^* = 0.20/2.37\times 10^{4} =
8.4\times 10^{-6}\,\si{\meter}^2$, a throat diameter of
$3.3\,\si{\milli\meter}$. The exit area is $3.85 A^* =
3.2\times 10^{-5}\,\si{\meter}^2$ ($6.4\,\si{\milli\meter}$ diameter),
and the constant-area mixing tube is sized to pass the combined
$0.30\,\si{\kilogram\per\second}$ at the mixed Mach number, typically
two to three primary-exit diameters across. The whole device fits in a
hand, has no seals or bearings, and runs as long as the primary supply
holds. The ejector converts a chapter's worth of relations, choked
primary flow, supersonic expansion, and constant-area mixing with a
recompression shock, into a pump with no moving parts, at the cost of
the primary stream's stagnation pressure spent to do the entrainment.

\section{Deepening: the Prandtl-Glauert transformation and its
breakdown}\pathtag{standard}

The subsonic-compressible section quoted the Prandtl-Glauert result
$C_p = C_{p,0}/\sqrt{1 - M_\infty^2}$ as a correction factor. The
transformation behind it repays a closer look, because it shows exactly
what compressibility does to a subsonic flow and exactly why the factor
diverges at $M_\infty = 1$. Small-disturbance subsonic flow obeys the
linearised potential equation
\[
(1 - M_\infty^2)\phi_{xx} + \phi_{yy} = 0,
\]
for the perturbation potential $\phi$. Below $M_\infty = 1$ the
coefficient $1 - M_\infty^2$ is positive and the equation is elliptic,
the same type as the incompressible Laplace equation $\phi_{xx} +
\phi_{yy} = 0$ but with the streamwise coordinate stretched. Introduce
the scaled coordinate $\xi = x$, $\eta = y\sqrt{1 - M_\infty^2}$: the
governing equation becomes Laplace's equation in $(\xi, \eta)$, so a
compressible subsonic flow past a thin body is an incompressible flow
past a body stretched in the flow-normal direction by
$1/\sqrt{1 - M_\infty^2}$. Carrying the boundary condition and the
pressure definition through the stretch delivers the correction factor:
the compressible pressure coefficient on the real body equals the
incompressible pressure coefficient on the stretched body, divided by
$\sqrt{1 - M_\infty^2}$. The physical content is that compressibility
makes a subsonic flow behave as if the body were effectively thicker in
its influence, so the suction peaks deepen and the lift grows, both by
the same $1/\sqrt{1 - M_\infty^2}$ factor. The divergence at
$M_\infty \to 1$ is not a failure of arithmetic but the signal that the
elliptic assumption has expired: as $1 - M_\infty^2 \to 0$ the stretch
becomes infinite, the linearisation that dropped the nonlinear terms is
no longer justified because the local perturbations are no longer small
compared with the vanishing $1 - M_\infty^2$, and a local pocket of the
flow reaches $M = 1$ while the freestream is still subsonic. The
Karman-Tsien and Laitone rules push the usable range higher by keeping a
nonlinear term the Prandtl-Glauert linearisation discards, but they too
diverge at $M_\infty = 1$, because no purely subsonic theory can survive
the appearance of an embedded supersonic pocket. The transformation
also explains the transonic-flow difficulty in one line: at
$M_\infty > 1$ the coefficient $1 - M_\infty^2$ turns negative and the
equation changes from elliptic to hyperbolic, so the mathematics itself
switches character at $M = 1$, and the transonic regime is the narrow
band where both types coexist on the same body.

\section{Deepening: why a supersonic diffuser is hard}\pathtag{standard}

The nozzle section named the supersonic diffuser as the dual of the
supersonic nozzle and left its difficulty as an assertion. The
difficulty is worth making concrete, because it is the reason every
supersonic inlet is a compromise and the reason the starting problem of
the inlet case study exists at all. A supersonic nozzle and a supersonic
diffuser run the same area-Mach relation in opposite directions: the
nozzle accelerates a subsonic flow to sonic at the throat and on to
supersonic in the diverging section, and the diffuser should decelerate
a supersonic flow to sonic at the throat and on to subsonic in the
diverging section. The nozzle works cleanly because acceleration is
stable: a supersonic flow that speeds up in a diverging duct stays
supersonic, and a small perturbation does not run away. Deceleration is
the trouble. An ideal isentropic diffuser would slow the supersonic flow
to exactly $M = 1$ at the throat and then diffuse it subsonically, with
no shock and no loss, but this ideal is unreachable in practice for two
reasons that compound.

First, the isentropic-deceleration state is unstable to the throat Mach
number. If the throat area is set for exactly $M = 1$ and a small
disturbance nudges the flow, the throat cannot hold sonic conditions:
the shock that forms to reconcile the mismatch is not infinitesimal but
jumps to a finite strength, and it parks in the diverging section rather
than at the throat, spending stagnation pressure the ideal design
assumed it would keep. Second, the started state is hard to reach from
rest. On start-up the flow is subsonic everywhere, a normal shock must
be swallowed through the throat to establish supersonic flow ahead of
it, and a throat sized for efficient running (near the isentropic limit)
is too small to swallow that starting shock. This is the Kantrowitz
limit made physical: the most aggressive contraction that will
self-start is far milder than the most aggressive contraction that runs
well once started, and the gap between the two is the band a
fixed-geometry inlet cannot occupy. The design responses are the ones
the inlet case study enumerated, oversizing the throat and accepting a
weaker recovery, adding variable geometry to open the throat for
starting and close it for running, or bleeding the boundary layer to
keep the shock positioned. The deeper point is that compressible-flow
deceleration is intrinsically lossier and less stable than acceleration,
a one-sidedness that traces to the same entropy law that makes shocks
compress rather than expand: a diffuser must fight the flow's preference
to make its pressure rise through an irreversible shock rather than a
reversible diffusion, and every real supersonic inlet is a managed
truce with that preference.

\section{Deepening: real-gas effects and the frozen-versus-equilibrium
question}\pathtag{mastery}

The hypersonic section listed real-gas effects as bullet points:
dissociation, ionisation, radiation. The list understates a shift in the
governing physics that changes how every relation in the chapter must be
read above roughly $M = 8$, and the shift is worth making concrete
because it is why the perfect-gas stagnation temperature of the re-entry
case study is a regime flag rather than a wall temperature. Behind a
strong shock the translational temperature spikes to the perfect-gas
value in the first few mean free paths, and then the energy begins to
redistribute. It flows first into molecular vibration, then into
dissociation as $\mathrm{O}_2$ and $\mathrm{N}_2$ break into atoms, and
at higher temperatures into ionisation. Each of these is a reservoir that
absorbs energy at nearly constant temperature, the way a boiling liquid
holds its temperature while it takes in the heat of vaporisation, so the
true temperature climbs far less than the perfect-gas relation predicts;
the missing enthalpy is stored in the chemical and internal modes. The
effective ratio of specific heats drops from $1.4$ toward $1.1$ or lower
as these modes open, which is why the density ratio across a real
hypersonic shock exceeds the perfect-gas ceiling of six: a lower
effective $\gamma$ raises $(\gamma + 1)/(\gamma - 1)$, and a dissociating
shock layer can reach density ratios of ten or more, thinning the shock
layer and shrinking the standoff distance the case study estimated.

The sharper question is the rate at which these reservoirs fill and
empty relative to the time the gas spends in the flow. If the chemistry
is fast compared with the flow time, the gas is in \engterm{equilibrium}
at each point, and the composition is a function of the local pressure
and temperature alone, so the flow is solved with equilibrium-air
property tables in place of the ideal-gas law. If the chemistry is slow
compared with the flow time, the composition is \engterm{frozen} at
whatever it was upstream, and the gas behaves like a perfect gas of
fixed (dissociated) composition. Real flows sit between: the gas
dissociates behind the shock at a finite rate, convects downstream while
still reacting, and may not have recombined by the time it reaches the
wall, so the composition carries a memory of its history. This
\engterm{non-equilibrium} regime is the hardest case, because the
properties at a point depend on the path the gas took to get there and
the governing equations carry finite-rate chemical source terms. The
wall condition compounds it: a \engterm{catalytic} wall lets the atoms
recombine on the surface, dumping the recombination enthalpy into the
wall as an extra heat load, while a \engterm{non-catalytic} wall lets
the atoms leave still dissociated, carrying that enthalpy away. The
difference between the two wall assumptions can double the stagnation
heat flux, which is why the Fay-Riddell correlation of the case study
carries the wall catalycity explicitly and why the surface material of a
re-entry vehicle is chosen partly for its low catalycity. The
perfect-gas relations of this chapter remain the right first pass and
the right way to identify the regime, but above the dissociation
threshold the working analysis replaces the ideal-gas law with an
equilibrium or finite-rate chemistry model, and the change is not a
correction factor but a change in the equations.

\section{Deepening: shock--boundary-layer interaction}\pathtag{standard}

The failure section named shock-induced separation as a transonic and
supersonic hazard and left the mechanism as an assertion. The
interaction between a shock and the boundary layer it strikes is worth
making concrete, because it is the single most common way the clean
inviscid relations of this chapter fail against a real surface, and it
governs the performance of inlets, transonic wings, and nozzles alike.
The inviscid relations treat a shock as a discontinuity that the flow
crosses in one step. A real surface carries a boundary layer, a thin
region of retarded flow in which the velocity falls from the freestream
value to zero at the wall, and the lower part of that layer is subsonic
even when the outer flow is supersonic. A shock cannot exist in the
subsonic part, so the sharp pressure rise the inviscid shock imposes on
the outer flow cannot be delivered as a discontinuity to the wall:
instead the pressure rise feeds forward through the subsonic sublayer,
thickening the boundary layer ahead of the shock foot and, if the rise
is strong enough, separating it.

The consequences depend on the strength of the interaction. A weak shock
on a turbulent boundary layer thickens the layer and smears the shock
foot into a small lambda-shaped structure without separating the flow,
and the inviscid pressure rise is delivered with only a modest loss. A
strong shock, or any shock on the more fragile laminar boundary layer,
separates the flow: the boundary layer lifts off the surface ahead of
the shock, forms a recirculating bubble, and either reattaches
downstream or, if the shock is strong enough, does not reattach at all.
A separated interaction replaces the single shock with a system, a
separation shock where the flow first turns away from the wall, a
compression fan or reattachment shock where it turns back, and the
lambda-foot pattern that schlieren images of transonic wings show
plainly. The separated shear layer is unsteady, and its oscillation is
the source of the transonic buffet the failure section named: the shock
and the separation bubble pulse together at structural frequencies. The
interaction also sets the drag-rise penalty, because a separated
boundary layer carries a large pressure drag the attached flow does not,
and it caps the pressure recovery of a supersonic inlet, because the
terminal shock sitting on a thick or separated boundary layer loses more
stagnation pressure than the inviscid relations of the inlet case study
credited. The design responses are the ones that keep the boundary layer
thin and energetic where the shock lands: boundary-layer bleed through
porous or slotted walls at the shock foot, vortex generators that
re-energise the layer by mixing in high-momentum outer flow, and shock
shaping that spreads the pressure rise over a longer distance so no
single shock is strong enough to separate the flow. The supercritical
section of section~13.6 is one such shaping, trading a single strong
shock for a weaker one over a flatter upper surface, and the reason it
works is that a weaker shock is less likely to separate the boundary
layer it strikes.

\section{Case study: blunt-body re-entry aerothermodynamics}\pathtag{standard}

The two earlier case studies ran the chapter's relations through a
nozzle and an inlet, devices that engineer the flow inside a duct. The
re-entry problem engineers a body in an open hypersonic stream, and it
inverts the streamlining instinct of the first twelve chapters of this
volume: the right re-entry nose is blunt, not sharp, and the reason is
the stagnation heat flux. We carry a single re-entry condition from the
freestream through the bow shock to the wall, estimate the stagnation
temperature and pressure behind a normal shock, compute the stagnation
heat flux from the Sutton-Graves and Fay-Riddell correlations, and show
why enlarging the nose radius (which a drag-minimising designer would
never do) is the lever that makes the vehicle survive. The arithmetic
reuses \texttt{code/normal\_shock.py} and \texttt{code/isentropic.py};
the heat-flux correlations are evaluated directly.

\paragraph{Requirement and entry condition.} A capsule returns from
low Earth orbit and passes through its peak-heating point at altitude
$h = 70\,\si{\kilo\meter}$, where the standard atmosphere gives
$\rho_\infty = 8.3\times 10^{-5}\,\si{\kilogram\per\meter\cubed}$,
$p_\infty = 5.2\,\text{Pa}$, and $T_\infty = 219\,\si{\kelvin}$. The
vehicle speed there is $U_\infty = 7.0\,\si{\kilo\meter\per\second}$,
so with $a_\infty = \sqrt{1.4\times 287\times 219} = 297\,\si{\meter
\per\second}$ the freestream Mach number is
$M_\infty = 7000/297 = 23.6$, deep in the hypersonic regime. The
designer must choose a nose radius and a thermal-protection approach
that survive the stagnation heat flux at this point, which is the
worst point of the trajectory.

\paragraph{Step 1: the shock detaches, and the standoff distance.} The
forward face of a blunt capsule asks the flow to turn through close to
$90^{\circ}$, far above the maximum deflection any attached oblique
shock can provide at any Mach number (worked example~9 showed
detachment already at $25^{\circ}$ for $M_1 = 2$). The shock therefore
detaches and stands off the nose as a curved bow shock
(figure~\ref{fig:vol03ch13:blunt-body}). Along the stagnation
streamline the bow shock is locally normal to the flow, so the
stagnation-line gas is processed by a normal shock at $M_1 = 23.6$. The
standoff distance $\Delta$ scales with the nose radius and the
density ratio across the shock; a serviceable estimate for a sphere is
$\Delta/R_n \approx 0.78\,(\rho_1/\rho_2)$. The perfect-gas density
ratio saturates at $(\gamma + 1)/(\gamma - 1) = 6$, giving
$\rho_1/\rho_2 = 1/6$ and $\Delta/R_n \approx 0.13$; a real
dissociating shock layer compresses further, so the real standoff is
smaller, around $0.05\,R_n$. For a nose radius $R_n = 0.5\,\si{\meter}$
the bow shock stands roughly $25$ to $65\,\si{\milli\meter}$ ahead of
the nose, the thin hot shock layer marked in
figure~\ref{fig:vol03ch13:blunt-body}.

\paragraph{Step 2: stagnation temperature and pressure, perfect-gas
estimate.} The perfect-gas normal-shock and isentropic relations give a
first estimate that the next step will correct. The stagnation
temperature behind the shock is conserved across it (the shock is
adiabatic), so
\[
T_0 = T_\infty\!\left(1 + \frac{\gamma - 1}{2}M_\infty^2\right)
= 219\,(1 + 0.2\times 23.6^2) = 219\times 112.4 = 2.46\times 10^4\,
\si{\kelvin}.
\]
This number, near $25{,}000\,\si{\kelvin}$, is physically impossible:
long before it is reached the air dissociates and ionises, and the
energy that the perfect-gas model puts into temperature goes instead
into breaking and stripping molecular bonds. The true stagnation
temperature in the shock layer is far lower, around $6000$ to
$8000\,\si{\kelvin}$, with the missing enthalpy stored chemically. The
perfect-gas $T_0$ is a flag for the regime, not a wall temperature
(exercise~10 made the same point). The stagnation pressure behind the
normal shock, by contrast, is well estimated by the perfect-gas
Rayleigh pitot relation of worked example~10, and it is close to the
Newtonian value $p_{0,2} \approx p_\infty + \rho_\infty U_\infty^2 =
5.2 + 8.3\times 10^{-5}\times 7000^2 = 5.2 + 4067 = 4.07\,\text{kPa}$.
The capsule's forward face sees about $4\,\text{kPa}$ of pressure, modest
next to the thermal problem.

\paragraph{Step 3: stagnation heat flux from Sutton-Graves.} The
engineering workhorse for stagnation-point convective heat flux is the
Sutton-Graves correlation
\[
q_w = K\sqrt{\frac{\rho_\infty}{R_n}}\,U_\infty^3,
\qquad K \approx 1.83\times 10^{-4}\ \text{(SI, air)}.
\]
With $\rho_\infty = 8.3\times 10^{-5}\,\si{\kilogram\per\meter\cubed}$,
$R_n = 0.5\,\si{\meter}$, and $U_\infty = 7000\,\si{\meter\per
\second}$,
\[
q_w = 1.83\times 10^{-4}\sqrt{\frac{8.3\times 10^{-5}}{0.5}}\,
(7000)^3
= 1.83\times 10^{-4}\times 0.01288\times 3.43\times 10^{11}
= 8.1\times 10^{5}\,\si{\watt\per\meter\squared},
\]
about $0.81\,\si{\mega\watt\per\meter\squared}$. Reading the same
condition off figure~\ref{fig:vol03ch13:heat-flux-radius}: the
low-Earth-orbit curve ($U_\infty = 7\,\si{\kilo\meter\per\second}$) at
$R_n = 0.5\,\si{\meter}$ sits at about $0.4\,\si{\mega\watt\per
\meter\squared}$ for the figure's lower density
$\rho_\infty = 2\times 10^{-5}$; scaling by
$\sqrt{8.3/2} = 2.04$ for the higher density here recovers the
$0.81\,\si{\mega\watt\per\meter\squared}$ of the arithmetic. The
Fay-Riddell correlation, which carries the boundary-layer chemistry and
the wall-to-stagnation enthalpy ratio explicitly, lands within
$10$ to $20\,\%$ of the Sutton-Graves value for an equilibrium-catalytic
wall, so the simpler correlation is adequate for sizing.

\paragraph{Step 4: why blunt beats sharp.} Suppose a designer, trained
on the streamlined bodies of chapters 11 and 12, proposes a sharp nose
with an effective radius $R_n = 0.01\,\si{\meter}$ to cut the pressure
drag. The Sutton-Graves scaling $q_w \propto 1/\sqrt{R_n}$ multiplies
the stagnation heat flux by $\sqrt{0.5/0.01} = \sqrt{50} = 7.1$,
driving it from $0.81$ to $5.7\,\si{\mega\watt\per\meter\squared}$
(figure~\ref{fig:vol03ch13:heat-flux-radius}, far left). No reusable
material survives a sustained flux that high; the sharp nose burns
through. Enlarging the nose instead, to $R_n = 2.0\,\si{\meter}$, halves
the flux again to $0.41\,\si{\mega\watt\per\meter\squared}$. The lever
is the nose radius, and it points the opposite way from the drag
instinct: bluntness raises the pressure drag (worked example~14 found
$C_D = 1$ for a hemisphere) but lowers the heat flux as
$1/\sqrt{R_n}$, and the pressure drag is the friend here because it
decelerates the vehicle high in the thin atmosphere, before the dense
lower atmosphere can deliver its worst heating. Harvey Allen's
blunt-body insight (NACA, 1951) inverted the wartime assumption that a
re-entry body should be a sharp dart, and it is the reason every crewed
re-entry vehicle from Mercury to Crew Dragon presents a broad, blunt
heat shield to the flow.

% Figure: detached bow shock ahead of a blunt re-entry body versus an
% attached shock on a slender sharp body. Shows the standoff distance,
% the hot shock layer, and why a blunt nose spreads the heat load.
\begin{figure}[htbp]
\centering
\begin{tikzpicture}[scale=1.0]
% === Left panel: blunt body with detached bow shock ===
\begin{scope}[xshift=0cm]
  \node[font=\scriptsize, anchor=south] at (0.2,2.6) {blunt body};
  % freestream arrows
  \foreach \y in {-1.6,-0.8,0,0.8,1.6} {
    \draw[engblue, -{Latex[length=1.4mm]}] (-3.2,\y) -- (-2.5,\y);
  }
  \node[engblue, font=\scriptsize, anchor=east] at (-3.2,2.0) {$M_\infty\!\gg\!1$};
  % blunt nose: arc of large radius centred to the right
  \draw[engblack, very thick] (1.4,1.5) arc (55:-55:1.85);
  \draw[engblack, very thick] (1.4,1.5) -- (2.6,1.5);
  \draw[engblack, very thick] (1.4,-1.5) -- (2.6,-1.5);
  % bow shock: curved, standing ahead of the nose
  \draw[engred, very thick] (-0.1,2.3) .. controls (-1.4,0) .. (-0.1,-2.3);
  \node[engred, font=\scriptsize, anchor=south east] at (-0.7,1.7) {bow shock};
  % shock layer shading hint (hot region between shock and body)
  \node[engamber, font=\scriptsize] at (0.05,0) {hot};
  \node[engamber, font=\scriptsize] at (0.05,-0.45) {layer};
  % standoff distance marker
  \draw[enggray, {Latex[length=1.2mm]}-{Latex[length=1.2mm]}] (-1.05,-1.05) -- (-0.45,-1.05);
  \node[enggray, font=\scriptsize, anchor=north] at (-0.75,-1.1) {$\Delta$};
  % nose radius marker
  \draw[enggray, dashed] (-0.45,0) -- (1.25,0);
  \node[enggray, font=\scriptsize, anchor=south] at (0.5,0.02) {$R_n$};
\end{scope}
% === Right panel: slender sharp body with attached shock ===
\begin{scope}[xshift=7.6cm]
  \node[font=\scriptsize, anchor=south] at (0.2,2.6) {sharp body};
  \foreach \y in {-1.6,-0.8,0,0.8,1.6} {
    \draw[engblue, -{Latex[length=1.4mm]}] (-3.2,\y) -- (-2.5,\y);
  }
  % slender wedge nose
  \draw[engblack, very thick] (-1.0,0) -- (2.6,0.9);
  \draw[engblack, very thick] (-1.0,0) -- (2.6,-0.9);
  % attached oblique shocks springing from the tip
  \draw[engred, very thick] (-1.0,0) -- (1.9,2.0);
  \draw[engred, very thick] (-1.0,0) -- (1.9,-2.0);
  \node[engred, font=\scriptsize, anchor=south west] at (1.0,1.4) {attached shock};
  % small standoff (zero) note at tip
  \node[engamber, font=\scriptsize, anchor=west] at (-0.9,0.55) {$\Delta\!\to\!0$};
  \node[enggray, font=\scriptsize, anchor=west] at (-0.9,-0.6) {$R_n\!\to\!0$};
\end{scope}
\end{tikzpicture}
\caption[Bow shock on a blunt body versus an attached shock on a sharp
body]{Hypersonic flow over a blunt re-entry body (left) and a slender
sharp body (right). The blunt body cannot turn the flow through the
large angle its forward face demands, so the shock detaches and stands
off the nose by a distance $\Delta$ that scales with the nose radius
$R_n$ (for air at high Mach, $\Delta/R_n$ is roughly $0.1$ to $0.2$).
The gas in the shock layer between the bow shock and the body is
processed by a near-normal shock at the stagnation streamline, so it is
hot and slow. The blunt nose pays a large pressure drag but spreads the
stagnation heat load over a large radius, and the stagnation heat flux
falls as $1/\sqrt{R_n}$, so a large $R_n$ lowers the peak flux. The
sharp body keeps an attached shock and low drag, but its tip radius
tends to zero, which drives the local heat flux toward values no
material survives. The design lever for re-entry is bluntness: trade
pressure drag, which decelerates the vehicle high in the atmosphere
where the air is thin, for a survivable stagnation heat flux.}
\label{fig:vol03ch13:blunt-body}
\end{figure}


% Figure: Sutton-Graves stagnation heat flux versus nose radius for two
% entry speeds, showing the 1/sqrt(R_n) falloff that makes bluntness the
% design lever. q = K sqrt(rho/Rn) U^3, K=1.83e-4 SI, rho=2e-5 kg/m3.
% U=7 km/s and U=11 km/s. q in MW/m^2 (divide W/m2 by 1e6).
% prefactor A = K*sqrt(rho)*U^3 / 1e6:
%  U=7e3:  1.83e-4*sqrt(2e-5)*(7e3)^3/1e6 = 1.83e-4*4.4721e-3*3.43e11/1e6
%          = 1.83e-4*4.4721e-3=8.184e-7; *3.43e11=280.7; /1e6=2.807e-4? recompute below in comment
% Careful: 8.184e-7*3.43e11 = 2.807e5 W/m2 base (at Rn=1). /1e6 = 0.2807 MW/m2 at Rn=1.
%  so plot A/sqrt(Rn): U=7 -> 0.281/sqrt(Rn); U=11 -> scale by (11/7)^3=3.88 -> 1.090/sqrt(Rn)
\begin{figure}[htbp]
\centering
\begin{tikzpicture}
\begin{axis}[
  width=12cm, height=7.5cm,
  xlabel={nose radius $R_n$ (m)},
  ylabel={stagnation heat flux $q_w$ (MW/m$^2$)},
  xmin=0.05, xmax=2.0, ymin=0, ymax=5,
  axis lines=left,
  grid=major, grid style={enggray!20},
  tick label style={font=\scriptsize},
  label style={font=\small},
  legend style={font=\scriptsize, at={(0.97,0.97)}, anchor=north east},
  legend cell align=left,
]
% U = 11 km/s (lunar-return class)
\addplot[engred, very thick, domain=0.05:2.0, samples=120] { 1.090/sqrt(x) };
\addlegendentry{$U_\infty=11\,$km/s (lunar return)}
% U = 7 km/s (low-Earth-orbit return)
\addplot[engblue, very thick, domain=0.05:2.0, samples=120] { 0.281/sqrt(x) };
\addlegendentry{$U_\infty=7\,$km/s (LEO return)}
% mark a representative shuttle leading-edge radius
\draw[enggray, dashed] (axis cs:0.3,0) -- (axis cs:0.3,5);
\node[enggray, font=\scriptsize, anchor=south west, rotate=90] at (axis cs:0.3,0.2)
  {$R_n\approx0.3$ m (leading edge)};
\end{axis}
\end{tikzpicture}
\caption[Stagnation heat flux versus nose radius]{Stagnation-point heat
flux from the Sutton-Graves correlation $q_w=K\sqrt{\rho_\infty/R_n}\,
U_\infty^3$ at a peak-heating density $\rho_\infty\approx2\times10^{-5}\,
\si{\kilogram\per\meter\cubed}$, for a low-Earth-orbit return at
$7\,\si{\kilo\meter\per\second}$ (blue) and a lunar return at
$11\,\si{\kilo\meter\per\second}$ (red). The flux falls as
$1/\sqrt{R_n}$, so doubling the nose radius cuts the peak flux by about
thirty percent; this inverse-square-root dependence is the quantitative
reason re-entry vehicles are blunt. The cube on velocity is the steeper
scaling: the lunar-return curve sits a factor $(11/7)^3\approx3.9$ above
the orbital-return curve at every radius, which is why a return from the
Moon or a planet is a far harder thermal problem than a return from low
orbit. The dashed line marks a reinforced-carbon-carbon leading-edge
radius near $0.3\,\si{\meter}$.}
\label{fig:vol03ch13:heat-flux-radius}
\end{figure}


\paragraph{Step 5: the integrated heat load and the material choice.}
Peak flux sizes the surface temperature; the integrated heat load over
the trajectory sizes the thermal-protection thickness. A capsule spends
roughly $100\,\si{\second}$ near peak heating, so the order-of-magnitude
integrated load at the stagnation point is
$Q \approx q_w\,\Delta t \approx 0.8\times 10^{6}\times 100 =
8\times 10^{7}\,\si{\joule\per\meter\squared}$, about
$80\,\si{\mega\joule\per\meter\squared}$. An ablative shield carries
this away by pyrolysis and surface recession: at a heat of ablation of
order $10\,\si{\mega\joule\per\kilogram}$ for a carbon-phenolic or PICA
material, the mass ablated per unit area is roughly
$Q/h_{\text{abl}} \approx 8\,\si{\kilogram\per\meter\squared}$, a few
millimetres of recession plus an insulating char layer that holds the
bondline below its structural limit near $450\,\si{\kelvin}$. A
low-flux, many-cycle profile (the Space Shuttle's shallow entry) favours
a reusable radiative ceramic instead, which radiates the incoming flux
back to the flow at a high surface emissivity rather than carrying it
away with departing mass. The design split between ablative and
radiative protection is set at this step, by the peak flux and the
integrated load together, not by either alone.

\paragraph{Verdict.} The blunt nose is the correct re-entry shape, and
the numbers say why. At the peak-heating point ($M_\infty = 23.6$,
$U_\infty = 7\,\si{\kilo\meter\per\second}$) a $0.5\,\si{\meter}$ nose
sees a stagnation heat flux near $0.8\,\si{\mega\watt\per\meter\squared}$
and an integrated load near $80\,\si{\mega\joule\per\meter\squared}$,
survivable with a few millimetres of ablator or a reusable ceramic
tile, while a $0.01\,\si{\meter}$ sharp nose sees seven times the flux
and burns through. The perfect-gas stagnation temperature of
$25{,}000\,\si{\kelvin}$ is a regime flag, not a wall temperature: the
real shock layer dissociates and ionises, and the thermal-protection
problem is governed by the convective (and, at lunar-return speeds,
radiative) heat flux to the wall, not by the nominal stagnation
temperature. The Columbia accident \cite{acc:caib-columbia} is the
failure-side reading of this same physics: a breach in the
leading-edge protection collapses the local effective radius, the
$1/\sqrt{R_n}$ factor spikes the local flux, and the hot shock-layer gas
enters a structure with no thermal margin. The chapter's relations, run
from freestream to wall, turn the blunt-versus-sharp choice into a
heat-flux number and a material.

\section{Case study: designing a supersonic exhaust nozzle from
chamber to plume}\pathtag{standard}

The relations of this chapter close a complete engineering design when
they are run in sequence rather than singly. We size a
converging-diverging nozzle for a small liquid-propellant rocket
engine, carry the numbers from the chamber to the exit plane, predict
the thrust, and then follow the nozzle off its design point as the
vehicle climbs. The arithmetic is reproducible with
\texttt{code/nozzle.py}, \texttt{code/isentropic.py}, and
\texttt{code/normal\_shock.py}.

\paragraph{Requirement and chamber conditions.} The engine is to
produce a vacuum thrust of about $25\,\text{kN}$ at a chamber pressure
$p_0 = 5\,\text{MPa}$ and chamber temperature $T_0 = 3000\,\si{
\kelvin}$. The combustion products are modelled as a calorically
perfect gas with $\gamma = 1.20$ and specific gas constant
$R = 350\,\si{\joule\per\kilogram\per\kelvin}$, values representative
of a kerosene-oxygen mixture (the lower $\gamma$ and higher $R$ of hot
combustion products, not the $\gamma = 1.4$, $R = 287$ of cold air).
The design exit Mach number is set by the altitude at which the nozzle
should run matched; we choose $M_e = 3.5$, which for these gas
properties expands the flow to a pressure suited to high-altitude
operation.

\paragraph{Step 1: throat area from the choked mass flow.} The chamber
feeds a converging section that chokes at the throat. The choked mass
flux is
\[
\frac{\dot m}{A^*} = \frac{p_0}{\sqrt{T_0}}\sqrt{\frac{\gamma}{R}}
\left(\frac{2}{\gamma + 1}\right)^{(\gamma + 1)/[2(\gamma - 1)]}.
\]
With $\gamma = 1.20$ the exponent is $(\gamma + 1)/[2(\gamma - 1)] =
2.2/0.4 = 5.5$ and $(2/2.2)^{5.5} = 0.909^{5.5} = 0.586$. The radical
is $\sqrt{1.20/350} = \sqrt{3.43\times 10^{-3}} = 0.0586$. So
\[
\frac{\dot m}{A^*} = \frac{5\times 10^{6}}{\sqrt{3000}}\times 0.0586
\times 0.586 = \frac{5\times 10^{6}}{54.8}\times 0.0343
= 3130\,\si{\kilogram\per\second}/\si{\meter}^2.
\]
We first fix the mass flow from the thrust target. The exit velocity
(step 3 below) is $u_e \approx 2400\,\si{\meter\per\second}$, and for a
near-matched nozzle the thrust is close to $\dot m\,u_e$, so
$\dot m \approx 25{,}000/2400 \approx 10.4\,\si{\kilogram\per\second}$.
The throat area follows: $A^* = \dot m/(\dot m/A^*) = 10.4/3130 =
3.32\times 10^{-3}\,\si{\meter}^2$, a throat diameter of
$65\,\si{\milli\meter}$.

\paragraph{Step 2: exit area from the design Mach number.} The
area-Mach relation fixes the exit-to-throat ratio for $M_e = 3.5$ at
$\gamma = 1.20$:
\[
\frac{A_e}{A^*} = \frac{1}{M_e}\left[\frac{2}{\gamma + 1}
\left(1 + \frac{\gamma - 1}{2}M_e^2\right)\right]^{(\gamma + 1)/
[2(\gamma - 1)]}.
\]
The bracket is $\frac{2}{2.2}(1 + 0.10\times 12.25) =
0.909\times 2.225 = 2.023$, raised to the power $5.5$ gives
$2.023^{5.5} = 47.9$, and dividing by $M_e = 3.5$ gives
$A_e/A^* = 13.7$ (\texttt{code/nozzle.py}). The exit area is
$A_e = 13.7\times 3.32\times 10^{-3} = 4.55\times 10^{-2}\,\si{
\meter}^2$, an exit diameter of $241\,\si{\milli\meter}$. The large
area ratio is the visible signature of a high-expansion upper-stage or
vacuum nozzle; a sea-level booster nozzle with a smaller design $M_e$
has a much smaller bell.

\paragraph{Step 3: exit thermodynamic state and velocity.} The
isentropic relations at $M_e = 3.5$, $\gamma = 1.20$ give the static
ratios. The temperature ratio is
$T_0/T_e = 1 + 0.10\times 12.25 = 2.225$, so
$T_e = 3000/2.225 = 1348\,\si{\kelvin}$. The pressure ratio is
$p_0/p_e = 2.225^{\gamma/(\gamma - 1)} = 2.225^{6} = 121$, so
$p_e = 5\times 10^{6}/121 = 41.3\,\text{kPa}$. The exit speed of sound
is $a_e = \sqrt{\gamma R T_e} = \sqrt{1.20\times 350\times 1348} =
753\,\si{\meter\per\second}$, and the exit velocity is
$u_e = M_e a_e = 3.5\times 753 = 2636\,\si{\meter\per\second}$. (The
estimate $u_e \approx 2400$ used in step~1 was close enough to size
the throat; one pass of refinement with $u_e = 2636$ trims the mass
flow to $\dot m = 25{,}000/2636 = 9.5\,\si{\kilogram\per\second}$ and
the throat diameter to $62\,\si{\milli\meter}$, a $5\,\%$ correction
that does not change the design.)

\paragraph{Step 4: thrust at the design point and in vacuum.} The
thrust equation is
\[
F = \dot m\,u_e + (p_e - p_a)A_e.
\]
At the matched design altitude $p_a = p_e$ and the pressure term
vanishes, so $F_{\text{design}} = \dot m\,u_e = 9.5\times 2636 =
25.0\,\text{kN}$, the requirement. In vacuum ($p_a = 0$) the pressure
term adds $p_e A_e = 41{,}300\times 4.55\times 10^{-2} = 1.88\,\text{
kN}$, giving $F_{\text{vac}} = 26.9\,\text{kN}$. The specific impulse
in vacuum is $I_{sp} = F_{\text{vac}}/(\dot m\,g_0) =
26{,}900/(9.5\times 9.81) = 289\,\si{\second}$, a credible figure for a
kerosene-oxygen upper stage.
Figure~\ref{fig:vol03ch13:nozzle-thrust-altitude} traces how the
delivered thrust climbs from its sea-level
value to the vacuum ceiling as the back-pressure penalty $p_a A_e$
falls away.

% Figure: rocket-nozzle thrust versus altitude for a nozzle sized for a
% fixed design altitude. Vacuum thrust is the ceiling; sea-level thrust
% is reduced by the ambient back-pressure term p_a A_e. The momentum
% thrust is constant; the pressure-thrust term grows with altitude as
% the ambient pressure falls.
\begin{figure}[htbp]
\centering
\begin{tikzpicture}
\begin{axis}[
  width=12cm, height=8cm,
  xlabel={altitude (km)}, ylabel={thrust (fraction of vacuum thrust)},
  xmin=0, xmax=60, ymin=0.78, ymax=1.02,
  axis lines=left,
  grid=major, grid style={enggray!20},
  tick label style={font=\scriptsize},
  label style={font=\small},
  legend style={font=\scriptsize, at={(0.97,0.03)}, anchor=south east},
  legend cell align=left,
]
% Thrust(h) = F_vac - p_a(h) * A_e. Model the momentum term + p_e*A_e as a
% constant C and subtract p_a(h)*A_e. Normalise so vacuum (p_a=0) gives 1.
% Use exponential atmosphere p_a = p0*exp(-h/H), H=7 km, and a coefficient
% k = p0*A_e / F_vac chosen so sea-level loss is about 0.18.
% F/F_vac = 1 - k*exp(-x/7), k=0.18.
\addplot[engblue, very thick, domain=0:60, samples=120] {1 - 0.18*exp(-x/7)};
\addlegendentry{delivered thrust $F(h)/F_{\mathrm{vac}}$}
% horizontal vacuum ceiling
\addplot[enggray, dashed, domain=0:60, samples=2] {1.0};
\addlegendentry{vacuum ceiling}
% design-altitude marker: where p_a = p_e, here taken as 12 km
\draw[engred, dashed] (axis cs:12,0.78) -- (axis cs:12,1.02);
\node[engred, font=\scriptsize, anchor=south, rotate=90] at (axis cs:12,0.86)
  {design altitude $p_a=p_e$};
% sea-level loss annotation
\draw[engamber, {Latex[length=2mm]}-{Latex[length=2mm]}]
  (axis cs:1.5,0.82) -- (axis cs:1.5,1.0);
\node[engamber, font=\scriptsize, anchor=west] at (axis cs:2.5,0.91)
  {pressure loss $p_a A_e$};
\end{axis}
\end{tikzpicture}
\caption[Rocket-nozzle thrust against altitude]{Delivered thrust of a
fixed-geometry rocket nozzle as a function of altitude, normalised to
the vacuum thrust. The thrust equation $F = \dot m\,u_e + (p_e - p_a)
A_e$ splits into a constant momentum-plus-exit-pressure part and a
back-pressure penalty $p_a A_e$ that shrinks as the ambient pressure
$p_a$ falls with altitude. At sea level the penalty is largest and the
nozzle delivers the least thrust; in vacuum the penalty vanishes and
the thrust reaches its ceiling. The nozzle is sized so that the exit
pressure matches the ambient pressure ($p_e = p_a$) at one design
altitude, marked by the dashed line; below it the nozzle is
overexpanded, above it underexpanded. The model uses an exponential
atmosphere with scale height $7\,\si{\kilo\meter}$ and a sea-level
penalty of about $18\,\%$, representative of a first-stage nozzle sized
for moderate altitude.}
\label{fig:vol03ch13:nozzle-thrust-altitude}
\end{figure}


\paragraph{Step 5: off-design behaviour at sea level.} Run at sea
level, where $p_a = 101\,\text{kPa}$, this nozzle is badly
overexpanded: the exit pressure $p_e = 41.3\,\text{kPa}$ is far below
ambient. The thrust drops by the back-pressure penalty,
$F_{\text{SL}} = \dot m\,u_e + (p_e - p_a)A_e = 25.0\,\text{kN} +
(41{,}300 - 101{,}000)\times 4.55\times 10^{-2} = 25.0 - 2.72 =
22.3\,\text{kN}$, a $17\,\%$ loss against the design point. The more
serious sea-level problem is flow separation. A rough engineering
criterion (the Summerfield criterion) is that the wall boundary layer
separates from the diverging contour when the wall pressure falls below
roughly $0.35$ to $0.40$ times ambient,
$p_{\text{sep}} \approx 0.4\,p_a = 40\,\text{kPa}$. Since the design
exit pressure $41.3\,\text{kPa}$ sits just at this threshold, a
sea-level firing of this high-altitude nozzle would separate near the
exit plane, with an oblique shock springing from the separation line
and an unsteady, asymmetric plume that imposes side loads on the nozzle
structure. Figure~\ref{fig:vol03ch13:plume-modes} shows the three plume
states; the sea-level firing sits in panel~(a), with the shock-diamond
pattern of an overexpanded plume.

% Figure: three exhaust-plume shapes for a fixed converging-diverging
% nozzle at three back pressures. Overexpanded (p_e < p_a): plume
% contracts through oblique shocks at the lip, shows shock diamonds.
% Matched (p_e = p_a): straight parallel plume. Underexpanded
% (p_e > p_a): plume expands through Prandtl-Meyer fans at the lip.
\begin{figure}[htbp]
\centering
\begin{tikzpicture}[scale=0.95]
% panel (a): overexpanded
\begin{scope}[xshift=0cm]
  % nozzle wall (converging-diverging, only diverging exit shown)
  \draw[engblue, very thick] (-0.8,0.7) -- (0,0.55);
  \draw[engblue, very thick] (-0.8,-0.7) -- (0,-0.55);
  % contracting plume with diamond pattern
  \draw[engred, thick] (0,0.55) -- (0.7,0.18) -- (1.4,0.5) -- (2.1,0.18) -- (2.8,0.45);
  \draw[engred, thick] (0,-0.55) -- (0.7,-0.18) -- (1.4,-0.5) -- (2.1,-0.18) -- (2.8,-0.45);
  % oblique shock crossings (the diamonds)
  \draw[engamber, thin] (0,0.55) -- (0.7,-0.18);
  \draw[engamber, thin] (0,-0.55) -- (0.7,0.18);
  \draw[engamber, thin] (1.4,0.5) -- (2.1,-0.18);
  \draw[engamber, thin] (1.4,-0.5) -- (2.1,0.18);
  \node[engblack, font=\scriptsize, anchor=north] at (1.2,-0.85) {(a) overexpanded};
  \node[engblack, font=\scriptsize, anchor=north] at (1.2,-1.15) {$p_e<p_a$};
\end{scope}
% panel (b): matched
\begin{scope}[xshift=4.3cm]
  \draw[engblue, very thick] (-0.8,0.7) -- (0,0.55);
  \draw[engblue, very thick] (-0.8,-0.7) -- (0,-0.55);
  % straight parallel plume
  \draw[engred, thick] (0,0.55) -- (2.8,0.55);
  \draw[engred, thick] (0,-0.55) -- (2.8,-0.55);
  \node[engblack, font=\scriptsize, anchor=north] at (1.2,-0.85) {(b) matched (design)};
  \node[engblack, font=\scriptsize, anchor=north] at (1.2,-1.15) {$p_e=p_a$};
\end{scope}
% panel (c): underexpanded
\begin{scope}[xshift=8.6cm]
  \draw[engblue, very thick] (-0.8,0.7) -- (0,0.55);
  \draw[engblue, very thick] (-0.8,-0.7) -- (0,-0.55);
  % expanding plume
  \draw[engred, thick] (0,0.55) -- (2.8,1.05);
  \draw[engred, thick] (0,-0.55) -- (2.8,-1.05);
  % expansion fan lines at the lip
  \draw[englime, thin] (0,0.55) -- (1.0,0.95);
  \draw[englime, thin] (0,0.55) -- (1.6,0.78);
  \draw[englime, thin] (0,-0.55) -- (1.0,-0.95);
  \draw[englime, thin] (0,-0.55) -- (1.6,-0.78);
  \node[engblack, font=\scriptsize, anchor=north] at (1.2,-1.25) {(c) underexpanded};
  \node[engblack, font=\scriptsize, anchor=north] at (1.2,-1.55) {$p_e>p_a$};
\end{scope}
\end{tikzpicture}
\caption[Exhaust-plume shapes by expansion state]{The three plume
states of a fixed converging-diverging nozzle as the ambient pressure
$p_a$ varies. (a) Overexpanded ($p_e < p_a$, the nozzle at lower
altitude than its design point): the higher ambient pressure squeezes
the plume through oblique shocks that originate at the lip, reflect off
the plume boundary, and form the periodic shock-diamond pattern visible
in sea-level rocket exhaust. (b) Matched ($p_e = p_a$, the design
altitude): the plume leaves parallel with no lip shock or fan. (c)
Underexpanded ($p_e > p_a$, above the design altitude): the
still-higher exit pressure drives the flow through Prandtl-Meyer
expansion fans at the lip and the plume widens. A fixed nozzle passes
through all three as a launch vehicle climbs, which is why the visible
plume of a first stage broadens with altitude.}
\label{fig:vol03ch13:plume-modes}
\end{figure}


\paragraph{Step 6: a normal shock inside the bell.} If the
overexpansion is severe enough, the readjustment to ambient pressure
happens not at the lip but through a normal shock standing inside the
diverging section. We can locate it for a milder case. Suppose the same
nozzle is fed at a reduced chamber pressure
$p_0 = 1.2\,\text{MPa}$ while still firing at sea level. The fully
supersonic exit pressure would be
$p_e = 1.2\times 10^{6}/121 = 9.9\,\text{kPa}$, far below ambient, so
no fully supersonic solution can match the back pressure; a normal
shock must stand somewhere in the bell. The shock sits at the area
station where the post-shock subsonic diffusion just recovers to
$101\,\text{kPa}$ at the exit. Marching the supersonic branch outward,
trying a shock at the local Mach number $M_x$, computing the
post-shock subsonic Mach number and stagnation pressure from
\texttt{code/normal\_shock.py}, and diffusing isentropically to the
exit area, the back pressure $101\,\text{kPa}$ is matched when the
pre-shock Mach number is about $M_x = 2.05$, which the area-Mach
relation places at $A_x/A^* \approx 3.0$, roughly a fifth of the way
out the bell. Below that chamber pressure the shock moves further in;
above it, toward the exit and then out to the lip. This is the
machinery behind the operating modes of
figure~\ref{fig:vol03ch13:nozzle-modes}, now carried on a concrete
nozzle.

\paragraph{Verdict.} The nozzle meets its $25\,\text{kN}$ vacuum-class
thrust requirement with a $62\,\si{\milli\meter}$ throat, a
$241\,\si{\milli\meter}$ exit, an area ratio of $13.7$, and a vacuum
specific impulse near $289\,\si{\second}$. It must not be fired at sea
level: at the design expansion it sits at the separation threshold, and
any margin lost (a low chamber-pressure transient at start-up, a
throttled condition) drives a normal shock into the bell with the side
loads that have damaged real engines during start and shutdown
transients. The design choices that follow are the ones every launch
program makes: either accept the sea-level penalty and separation risk
of a single high-altitude nozzle, size a shorter sea-level nozzle and
take the vacuum-performance penalty, or add an altitude-compensating
device (a dual-bell contour or an aerospike) that keeps the flow
attached across the climb. The chapter's relations, run end to end,
turn each of those choices into a number.

\section{Vignette: the altitude-compensating aerospike
nozzle}\pathtag{standard}

The exhaust-nozzle case study closed on a dilemma: a fixed bell nozzle
is matched at exactly one altitude and pays a penalty everywhere else,
overexpanded and separation-prone near the pad, underexpanded and
throwing away its exit pressure high up. This shorter vignette follows
the escape the case study named in passing, the \engterm{aerospike}
nozzle that lets the ambient pressure set the expansion at every
altitude, and it sizes the compensation the aerospike delivers against
the fixed bell's loss. The device is worth a page because it turns the
back-pressure operating modes of section~13.5 from a fault to be endured
into a boundary condition to be exploited.

\paragraph{The geometry and its idea.} A bell nozzle expands the flow
against a closed diverging wall, so the exit area, and with it the exit
pressure, is fixed by the metal. An aerospike inverts the surface: the
flow expands along a central spike that tapers to a point, with the
outer edge of the jet open to the atmosphere
(figure~\ref{fig:vol03ch13:aerospike}). The free boundary of the plume
is then the ambient pressure itself. At sea level the high ambient
pressure presses the jet in against the spike, holding the effective
expansion small; at altitude the low ambient pressure lets the jet
stand off and expand further. The nozzle compensates for altitude with
no moving parts, because the atmosphere does the adjusting.

\paragraph{The compensation, in numbers.} Take the case-study engine,
which was sized for a matched exit pressure $p_e = 41.3\,\text{kPa}$ at
its design altitude. Fired at sea level ($p_a = 101\,\text{kPa}$) the
bell suffers the overexpansion penalty the case study computed, a
thrust drop of about $2.7\,\text{kN}$ from the back-pressure term
$(p_e - p_a)A_e$ plus the separation risk of a wall pressure at the
Summerfield threshold. An aerospike carrying the same chamber flow does
not hold $p_e = 41.3\,\text{kPa}$ at sea level: the ambient pressure
closes the jet down to an effective exit near $p_a$, so the
$(p_e - p_a)A_e$ penalty nearly vanishes and the flow stays attached to
the spike. The gain is exactly the penalty the bell paid, recovered
across the low-altitude part of the climb where the fixed nozzle is
worst. At altitude the two converge: once the ambient pressure falls to
the bell's design value the bell is matched and the aerospike offers
little further advantage, and in vacuum both approach the same ceiling
set by the full expansion of the chamber flow. The delivered
thrust-coefficient curves of figure~\ref{fig:vol03ch13:aerospike} carry
this: the bell peaks sharply at its design altitude and falls away on
both sides, while the aerospike tracks near its optimum from the pad to
vacuum.

% Figure: the altitude-compensating aerospike nozzle. A conventional
% bell nozzle expands the flow against a fixed wall, so its exit pressure
% is fixed by geometry and matches the ambient only at one altitude. An
% aerospike expands the flow along a central spike with the outer
% boundary open to the atmosphere, so the ambient pressure itself forms
% the outer boundary of the jet and the effective expansion adjusts with
% altitude. The schematic (left) shows the spike and the plume boundary
% at two altitudes; the plot (right) shows delivered thrust coefficient
% versus altitude for a bell (fixed) and an aerospike (compensating).
\begin{figure}[htbp]
\centering
\begin{tikzpicture}[scale=1.0]
% left panel: schematic of the aerospike
% throat plane at x=0, spike tapering to a point at right
\draw[engblack, very thick] (0,0.9) -- (1.6,0.0);   % upper spike contour
\draw[engblack, very thick] (0,-0.9) -- (1.6,0.0);  % lower spike contour
\draw[enggray, thick] (0,1.5) -- (0,0.9);           % upper cowl lip
\draw[enggray, thick] (0,-1.5) -- (0,-0.9);         % lower cowl lip
% sea-level plume boundary (tight, pinched by high ambient pressure)
\draw[engblue, thick] (0,1.5) .. controls (1.2,1.0) and (2.0,0.4) .. (2.4,0.0);
\draw[engblue, thick] (0,-1.5) .. controls (1.2,-1.0) and (2.0,-0.4) .. (2.4,0.0);
\node[engblue, font=\scriptsize, anchor=west] at (2.5,0.5) {sea level};
% high-altitude plume boundary (wide, low ambient pressure)
\draw[engred, thick, dashed] (0,1.5) .. controls (1.8,1.7) and (3.0,1.4) .. (3.6,0.6);
\draw[engred, thick, dashed] (0,-1.5) .. controls (1.8,-1.7) and (3.0,-1.4) .. (3.6,-0.6);
\node[engred, font=\scriptsize, anchor=west] at (2.9,-1.2) {altitude};
% flow arrows
\draw[englime, -stealth, thick] (-0.5,0.4) -- (0.1,0.35);
\draw[englime, -stealth, thick] (-0.5,-0.4) -- (0.1,-0.35);
\node[engblack, font=\scriptsize] at (-0.55,0.9) {flow};
\node[engblack, font=\scriptsize, anchor=north] at (1.0,-1.9) {spike + open jet boundary};
% right panel: thrust coefficient vs altitude
\begin{axis}[
  at={(6.2cm,-2.0cm)},
  width=7.2cm, height=6.0cm,
  xlabel={altitude (km)},
  ylabel={delivered thrust coefficient $C_F$},
  xmin=0, xmax=30, ymin=1.0, ymax=1.9,
  axis lines=left,
  grid=major, grid style={enggray!20},
  tick label style={font=\scriptsize},
  label style={font=\scriptsize},
  legend style={font=\scriptsize, at={(0.97,0.03)}, anchor=south east},
  legend cell align=left,
]
% Aerospike: near-optimum at every altitude (rises smoothly toward vacuum ceiling)
\addplot[engred, very thick] coordinates
  {(0,1.35)(5,1.50)(10,1.62)(15,1.71)(20,1.77)(25,1.80)(30,1.82)};
\addlegendentry{aerospike (compensating)}
% Bell nozzle designed for ~12 km: peaks there, worse above and below
\addplot[engblue, very thick] coordinates
  {(0,1.18)(5,1.42)(10,1.60)(12,1.66)(15,1.63)(20,1.58)(25,1.55)(30,1.53)};
\addlegendentry{bell (fixed, design $12$ km)}
\end{axis}
\end{tikzpicture}
\caption[The altitude-compensating aerospike nozzle]{Left: an aerospike
expands the exhaust along a central spike with the outer edge of the jet
open to the atmosphere, so the ambient pressure itself forms the free
boundary of the plume. At sea level (blue) the high ambient pressure
pinches the jet against the spike; at altitude (red dashed) the lower
ambient pressure lets the jet stand off and expand further. The
expansion adjusts with altitude at no cost in moving parts. Right: the
delivered thrust coefficient against altitude. A fixed bell nozzle
(blue) is matched at one design altitude (here $12$ km) and loses
performance above and below it, overexpanded and separation-prone near
sea level and underexpanded high up. The aerospike (red) stays near its
optimum across the whole climb, which is the appeal for a single-stage
vehicle that must fire efficiently from the pad to vacuum. The cost is a
spike that must survive the stagnation heating of the recirculating base
flow, which is why the aerospike has flown rarely despite the clean
altitude-compensation argument.}
\label{fig:vol03ch13:aerospike}
\end{figure}


\paragraph{Why it has flown rarely.} The altitude-compensation argument
is clean, and yet almost every flown engine is a bell. The reasons are
in the parts the idealised analysis omits. The spike sits in the
hottest, densest part of the flow and must survive stagnation-level
heating along its whole length, where a bell wall sees the flow only
after it has expanded and cooled, so the aerospike carries a severe
cooling problem that grows with engine size. Truncating the spike to a
short plug, the practical form, leaves a recirculating base flow whose
pressure must be predicted to get the thrust right, and that base flow
is hard to model and sensitive to the ambient conditions. The wetted
area and the manufacturing complexity are larger than a bell's for the
same thrust. The compensation is real, and the linear aerospike of the
X-33 program and the annular aerospike of several test engines
demonstrated it, but the thermal and base-flow penalties have kept the
bell the default. The vignette's lesson is that a device can be right on
the one-dimensional relations of this chapter and still lose to a
simpler competitor on the two- and three-dimensional effects the
one-dimensional analysis sets aside, which is the recurring verdict of
compressible-flow hardware.

\section{Case study: a Mach-3 supersonic inlet from ramp angles to
engine face}\pathtag{standard}

The nozzle case study ran the relations forward, expanding a
high-pressure reservoir to a supersonic plume. A supersonic inlet runs
them backward: it captures a supersonic stream and decelerates it to
the subsonic, high-pressure state a turbojet compressor or a ramjet
combustor requires, while keeping as much of the freestream stagnation
pressure as the shock physics allows. We design the shock system of a
two-ramp external-compression inlet for a vehicle cruising at
$M_\infty = 3.0$, carry the stagnation-pressure recovery through each
shock, compare the result against the single-normal-shock inlet it
replaces, and check that the chosen internal contraction will start.
The arithmetic reuses \texttt{code/oblique\_shock.py} for the ramp
shocks and \texttt{code/normal\_shock.py} for the terminal shock.

\paragraph{Requirement and strategy.} The vehicle flies at
$M_\infty = 3.0$ at $16\,\si{\kilo\meter}$, where
$T_\infty = 217\,\si{\kelvin}$ and $p_\infty = 9.7\,\text{kPa}$
(from \texttt{data/standard\_atmosphere.csv}). The engine face wants
the flow subsonic, at $M \approx 0.4$, with the highest stagnation
pressure the inlet can deliver. A single normal shock at $M_\infty = 3$
would do the deceleration in one step but throw away two thirds of the
stagnation pressure: from \texttt{code/normal\_shock.py},
$p_{02}/p_{01} = 0.328$ at $M_1 = 3$. The strategy that beats it is to
decelerate the flow through a staircase of weak oblique shocks, each
losing little because its entropy rise scales as the cube of
$(M_n^2 - 1)$, and to finish with a weak terminal normal shock at the
lowest Mach number the ramps can deliver. We use two compression ramps
(figure~\ref{fig:vol03ch13:inlet-shocks}); a single ramp would already
help, and three or four would help more, with diminishing returns and
growing length and weight.

\paragraph{Step 1: first ramp.} Choose a first ramp angle
$\theta_1 = 8^{\circ}$. From the $\theta$-$\beta$-$M$ relation at
$M_1 = 3.0$, $\theta = 8^{\circ}$, the weak shock angle is
$\beta_1 = 25.6^{\circ}$ (\texttt{code/oblique\_shock.py}). The
upstream normal Mach component is
$M_{n1} = M_1 \sin\beta_1 = 3.0 \times \sin 25.6^{\circ} = 1.296$.
The normal-shock relations at $M_{n1} = 1.296$ give the post-shock
normal component $M_{n2} = 0.787$ and the stagnation-pressure ratio
across this shock $p_{02}/p_{01} = 0.983$. The Mach number along the
deflected flow behind ramp~1 is
$M_2 = M_{n2}/\sin(\beta_1 - \theta_1) =
0.787/\sin 17.6^{\circ} = 2.60$. The flow is turned $8^{\circ}$,
compressed, and slowed from $M = 3.0$ to $M = 2.60$, and the first
shock kept $98.3\,\%$ of the stagnation pressure.

\paragraph{Step 2: second ramp.} Add a second ramp that turns the flow
a further $\theta_2 = 10^{\circ}$ relative to the ramp-1 flow
direction, so the local deflection is $10^{\circ}$ against the
oncoming $M_2 = 2.60$ stream. At $M = 2.60$, $\theta = 10^{\circ}$, the
weak shock angle is $\beta_2 = 31.9^{\circ}$. The normal component is
$M_{n2,\text{in}} = 2.60 \times \sin 31.9^{\circ} = 1.374$, which gives
across the second shock $p_{0}\text{-ratio} = 0.965$ and a deflected
Mach number $M_3 = M_{n2,\text{out}}/\sin(\beta_2 - \theta_2) =
0.752/\sin 21.9^{\circ} = 2.02$. After two ramps the flow runs at
$M_3 = 2.02$, having kept $0.983 \times 0.965 = 0.949$ of the
freestream stagnation pressure across both oblique shocks.

% Figure: external-compression supersonic inlet shock system. Two
% compression ramps throw two oblique shocks that meet at the cowl lip;
% a terminal normal shock just inside the lip drops the flow to
% subsonic. Schematic, not to scale; angles chosen for legibility.
\begin{figure}[htbp]
\centering
\begin{tikzpicture}[scale=1.0]
% ramps (the compression surface)
% first ramp from (0,0) to (3,0.9), second ramp to (5,2.1)
\draw[engblue, very thick] (0,0) -- (3,0.9) -- (5,2.1);
\node[engblue, font=\scriptsize, anchor=north] at (1.5,0.25) {ramp 1};
\node[engblue, font=\scriptsize, anchor=north] at (4.1,1.35) {ramp 2};
% cowl lip (upper body)
\draw[engblue, very thick] (5.2,3.0) -- (6.6,3.0);
\node[engblue, font=\scriptsize, anchor=south] at (5.9,3.05) {cowl};
% engine-face duct (subsonic)
\draw[enggray, thick] (5.0,2.1) -- (6.6,2.4);
\draw[enggray, thick] (6.6,3.0) -- (6.6,2.4);
\node[enggray, font=\scriptsize] at (6.95,2.7) {engine};
\node[enggray, font=\scriptsize] at (6.95,2.45) {face};
% freestream arrow
\draw[engblack, -{Latex[length=2mm]}, thick] (-1.4,1.6) -- (-0.3,1.6);
\node[engblack, font=\scriptsize, anchor=south] at (-0.85,1.62) {$M_\infty=3$};
% first oblique shock: from ramp-1 corner (0,0) up to lip
\draw[engred, thick] (0,0) -- (5.2,3.0);
\node[engred, font=\scriptsize, anchor=west, rotate=30] at (2.4,1.5) {shock 1};
% second oblique shock: from ramp-2 corner (3,0.9) up to lip
\draw[engamber, thick] (3,0.9) -- (5.2,3.0);
\node[engamber, font=\scriptsize, anchor=north west] at (3.7,1.6) {shock 2};
% terminal normal shock just inside the lip
\draw[englime, very thick] (5.2,2.1) -- (5.2,3.0);
\node[englime, font=\scriptsize, anchor=west] at (5.25,2.45) {normal};
\node[englime, font=\scriptsize, anchor=west] at (5.25,2.2) {shock};
% Mach-number labels in each zone
\node[engblack, font=\scriptsize] at (1.3,1.05) {$M_1=3$};
\node[engblack, font=\scriptsize] at (3.7,1.05) {$M_2$};
\node[engblack, font=\scriptsize] at (4.7,1.45) {$M_3$};
\node[engblack, font=\scriptsize] at (5.9,2.6) {$M_4<1$};
\end{tikzpicture}
\caption[External-compression supersonic inlet]{The shock system of a
two-ramp external-compression inlet at $M_\infty=3$, drawn
schematically. The freestream meets the first ramp and turns through an
oblique shock (shock~1), then the second ramp and a second oblique
shock (shock~2); the two ramps are angled so both shocks focus on the
cowl lip, the shock-on-lip condition that captures the full stream tube
without spillage. Each oblique shock is weak, so each loses little
stagnation pressure. A terminal normal shock just inside the lip drops
the now-decelerated supersonic flow to the subsonic state the engine
face requires. Stacking two weak oblique shocks ahead of a weak
terminal normal shock recovers far more stagnation pressure than a
single strong normal shock would at $M_\infty=3$, the quantitative
content of the inlet case study.}
\label{fig:vol03ch13:inlet-shocks}
\end{figure}


\paragraph{Step 3: terminal normal shock.} The flow at the cowl lip is
still supersonic at $M_3 = 2.02$, so a terminal normal shock drops it
to subsonic. From \texttt{code/normal\_shock.py} at $M_1 = 2.02$:
$M_4 = 0.573$ and $p_{04}/p_{03} = 0.722$. The terminal shock is the
single largest loss in the system, because it does the supersonic-to-
subsonic step in one jump, and it is far weaker than it would have been
at $M_\infty = 3$ because the ramps have already brought the flow down
to $M = 2.02$. The overall stagnation-pressure recovery of the
three-shock system is
\[
\frac{p_{0,\text{exit}}}{p_{0,\infty}} = 0.983 \times 0.965 \times
0.722 = 0.685.
\]
A short subsonic diffuser then slows the $M_4 = 0.573$ flow to the
$M \approx 0.4$ the engine face wants, an almost loss-free step because
subsonic diffusion at this Mach number is nearly isentropic.

\paragraph{Step 4: the comparison.} The single-normal-shock inlet at
$M_\infty = 3$ recovers $p_{02}/p_{01} = 0.328$. The two-ramp inlet
recovers $0.685$, slightly more than twice as much
(figure~\ref{fig:vol03ch13:inlet-recovery}).
% Figure: total-pressure recovery of a supersonic inlet versus
% freestream Mach number. Lower curve: single normal shock (pitot
% inlet). Upper curve: a two-shock external-compression inlet plus a
% terminal normal shock. Values are computed at sampled Mach numbers
% and entered as coordinate lists (no in-axis trig needed).
\begin{figure}[htbp]
\centering
\begin{tikzpicture}
\begin{axis}[
  width=12cm, height=8cm,
  xlabel={freestream Mach number $M_\infty$},
  ylabel={stagnation-pressure recovery $p_{0,\text{exit}}/p_{0,\infty}$},
  xmin=1, xmax=4, ymin=0, ymax=1.05,
  axis lines=left,
  grid=major, grid style={enggray!20},
  tick label style={font=\scriptsize},
  label style={font=\small},
  legend style={font=\scriptsize, at={(0.03,0.03)}, anchor=south west},
  legend cell align=left,
]
% single normal shock recovery p02/p01 at the freestream Mach number
\addplot[engred, very thick, mark=*, mark size=1.3pt] coordinates {
  (1.0,1.000) (1.25,0.987) (1.5,0.930) (1.75,0.835)
  (2.0,0.721) (2.5,0.499) (3.0,0.328) (3.5,0.213) (4.0,0.139) };
\addlegendentry{single normal shock (pitot inlet)}
% two-ramp external compression + terminal normal shock (illustrative
% optimised-ramp recovery, the kind a two-shock inlet attains)
\addplot[englime, very thick, mark=square*, mark size=1.3pt] coordinates {
  (1.0,1.000) (1.25,0.995) (1.5,0.980) (1.75,0.960)
  (2.0,0.935) (2.5,0.875) (3.0,0.800) (3.5,0.720) (4.0,0.640) };
\addlegendentry{two oblique shocks + normal shock}
% horizontal reference at 0.5
\draw[enggray, dashed] (axis cs:1,0.5) -- (axis cs:4,0.5);
\node[enggray, font=\scriptsize, anchor=south east] at (axis cs:4,0.5) {half recovery};
\end{axis}
\end{tikzpicture}
\caption[Inlet stagnation-pressure recovery]{Stagnation-pressure
recovery against freestream Mach number for two inlet strategies. The
lower curve is the recovery of a single normal shock, the pitot inlet,
which falls steeply because the normal-shock loss grows with upstream
Mach number; by $M_\infty=3$ it keeps only a third of the freestream
stagnation pressure and by $M_\infty=4$ barely a seventh. The upper
curve is a two-ramp external-compression inlet, in which two weak
oblique shocks decelerate the flow before a weak terminal normal shock
finishes the job; because each oblique shock loses little (its entropy
rise scales as the cube of $M_n^2-1$) the stacked system keeps much
more of the stagnation pressure. The gap between the curves is the
margin a multi-shock inlet buys, and it widens with Mach number, the
reason every practical supersonic inlet above about $M_\infty=1.6$
uses ramps or a centre-body spike rather than a plain pitot mouth.}
\label{fig:vol03ch13:inlet-recovery}
\end{figure}

Stagnation pressure is the
currency the engine spends to make thrust, so a recovery that doubles
roughly doubles the pressure available at the compressor face, with a
direct effect on thrust and specific fuel consumption. The penalty is
geometric: the ramps add length, wetted area, and weight, and the
shock-on-lip condition that captures the full stream tube without
spillage holds only at the design Mach number, so an inlet optimised
for $M_\infty = 3$ spills flow and loses recovery at other speeds. This
is why high-Mach inlets carry variable geometry.

\paragraph{Step 5: the starting check.} Recovery is not the only
constraint. An inlet with internal contraction (area falling from the
lip to a downstream throat) is conditionally stable: the converging
internal passage that gives good pressure recovery also makes the
started state hard to reach from rest. The
\engterm{Kantrowitz limit} is the most aggressive internal contraction
that will self-start at a given flight Mach number; a more aggressive
contraction will not swallow its starting normal shock and the inlet
never starts (figure~\ref{fig:vol03ch13:kantrowitz}).
% Figure: Kantrowitz starting limit and the isentropic contraction
% limit for a supersonic inlet, as throat-to-capture area ratio versus
% flight Mach number. Curves are sampled and entered as coordinate
% lists; no in-axis trigonometry.
\begin{figure}[htbp]
\centering
\begin{tikzpicture}
\begin{axis}[
  width=12cm, height=8cm,
  xlabel={flight Mach number $M_\infty$},
  ylabel={throat-to-capture area ratio $A_t/A_c$},
  xmin=1, xmax=5, ymin=0, ymax=1.05,
  axis lines=left,
  grid=major, grid style={enggray!20},
  tick label style={font=\scriptsize},
  label style={font=\small},
  legend style={font=\scriptsize, at={(0.97,0.97)}, anchor=north east},
  legend cell align=left,
]
\addplot[engblue, very thick, mark=*, mark size=1.3pt] coordinates {
  (1.0,1.000) (1.5,0.768) (2.0,0.661) (2.5,0.620)
  (3.0,0.600) (3.5,0.589) (4.0,0.582) (5.0,0.574) };
\addlegendentry{Kantrowitz (self-start) limit}
\addplot[engred, very thick, mark=square*, mark size=1.3pt] coordinates {
  (1.0,1.000) (1.5,0.578) (2.0,0.373) (2.5,0.254)
  (3.0,0.181) (3.5,0.133) (4.0,0.100) (5.0,0.060) };
\addlegendentry{isentropic (max steady) contraction}
\node[engblack, font=\scriptsize, align=center, text width=2.6cm]
  at (axis cs:2.9,0.86) {self-starts at rest geometry};
\node[engblack, font=\scriptsize, align=center, text width=3.0cm]
  at (axis cs:3.4,0.42) {starts only with variable geometry or overspeed};
\node[engblack, font=\scriptsize, align=center, text width=3.0cm]
  at (axis cs:3.7,0.12) {cannot run steadily (unstartable)};
\end{axis}
\end{tikzpicture}
\caption[Kantrowitz and isentropic inlet limits]{The two contraction
limits that bound supersonic-inlet design, plotted as the
throat-to-capture area ratio against flight Mach number. A smaller area
ratio is a more aggressive internal contraction. The upper (Kantrowitz)
curve is the most aggressive fixed contraction that will self-start: an
inlet whose geometry sits above this curve swallows its starting normal
shock on its own. The lower (isentropic) curve is the most aggressive
contraction the flow will sustain once started, where the throat just
reaches $M=1$. Between the two curves the inlet runs well once started
but will not start from its running geometry, so it needs variable
geometry, a starting bleed, or an overspeed-then-throttle schedule.
Below the isentropic curve no steady supersonic solution exists and the
inlet is unstartable as built. The widening gap between the curves at
high Mach number is why high-Mach inlets carry the most elaborate
starting hardware.}
\label{fig:vol03ch13:kantrowitz}
\end{figure}

At $M_\infty = 3$
the Kantrowitz throat-to-capture area ratio is about $A_t/A_c = 0.60$,
while the isentropic limit (the most aggressive contraction the flow
sustains once started, throat at $M = 1$) is about $0.18$. Our
external-compression design does its turning ahead of the lip, so its
internal contraction is mild and sits well above the Kantrowitz curve;
it self-starts. A mixed-compression inlet, which moves some of the
shock system inside the lip to recover even more stagnation pressure,
would push the internal contraction toward the isentropic limit, into
the band between the two curves where the inlet runs well once started
but cannot start from its running geometry. Such an inlet needs
variable geometry (a translating spike or a hinged ramp that opens the
throat for starting and closes it for running), a starting bleed, or an
overspeed-then-throttle schedule. The unstart hazard in section~13.8 is
the same physics seen from the failure side: a disturbance that pushes
the terminal shock forward of the throat expels the whole shock system
out the front of the inlet.

\paragraph{Verdict.} The two-ramp external-compression inlet meets the
deceleration requirement with a stagnation-pressure recovery of
$0.685$ at $M_\infty = 3$, against $0.328$ for the single-normal-shock
inlet it replaces, and its mild internal contraction self-starts
without moving parts. The design choices that follow are the ones every
supersonic-inlet program makes: accept the simplicity and wide
off-design Mach range of external compression at the cost of some
recovery, or
chase the higher recovery of mixed compression and pay for it with
variable geometry and an unstart-protection control law. Adding a third
ramp would lift the recovery toward $0.75$ at the cost of length and a
narrower good-recovery Mach band. The chapter's two relations, the
$\theta$-$\beta$-$M$ relation for the ramps and the Rankine-Hugoniot
relations for the terminal shock, turn each of those choices into a
recovery number and a starting margin.

\section{Failure: shock-induced separation, transonic instability, and aeroelastic flutter}\pathtag{core}

Every chapter in this work has a failure section. The
compressible-flow failure modes cluster into three families:
shock-induced boundary-layer separation and the associated
non-linear lift behaviour, transonic flutter and buffet, and
aeroelastic instability of structures in high-speed flow.

\subsection{Shock-induced separation and transonic buffet}

A shock impinging on a boundary layer raises the pressure
abruptly across a small streamwise distance. If the boundary
layer is laminar, it nearly always separates under the shock; if
turbulent, it may stay attached for weak shocks but separates
under stronger ones (typically when the pre-shock Mach number
exceeds about $1.3$). The shock-induced separation is unsteady at
typical transonic flight conditions: the shock oscillates
streamwise as the separated region pulses, producing
\engterm{transonic buffet} on the wing surface.

Transonic buffet is a flight-envelope boundary for civil
transport aircraft. The aircraft's operating Mach number must
stay below the buffet onset (typically $M = 0.85$ to $0.90$ in
cruise) to keep the structural-fatigue loads acceptable and the
pilot's workload manageable. Above buffet onset, the wing
oscillates at frequencies in the structural range, the cockpit
shakes, and the cabin experiences vibration. Modern transport
aircraft fly with a small margin to buffet onset under typical
cruise conditions; the margin shrinks at high cruise altitude and
heavy weight.

\subsection{Wing flutter at transonic and supersonic speeds}

Aeroelastic flutter is a self-excited instability in which the
wing's structural oscillation extracts energy from the airflow.
The classical case is the binary bending-torsion flutter familiar
from low-speed flight: the wing's bending mode couples with its
torsional mode through aerodynamic forces, and when the coupling
exceeds the wing's structural damping, the oscillation grows
exponentially until structural failure.

At transonic and supersonic speeds, new flutter modes appear:
\engterm{shock-induced flutter} (oscillation of a shock on the
wing surface couples with the wing's structural modes),
\engterm{panel flutter} (a flexible skin panel between supporting
ribs oscillates under the supersonic external pressure), and
\engterm{single-degree-of-freedom torsional flutter} in
transonic flow. The design discipline requires flutter analysis
across the full flight envelope, with computational
aerolasticity methods that couple unsteady CFD with structural
dynamics.

The original Tacoma Narrows suspension bridge of $1940$
\cite{acc:tacoma-narrows-1941} provides the canonical low-speed
analogue of aeroelastic flutter at the structural scale. The
deck's torsional self-excitation under steady wind at about
$19\,\si{\meter\per\second}$ was a single-degree-of-freedom
torsional flutter mode, in which the angular displacement of the
deck modified the aerodynamic moment in a way that fed energy
back into the deck's motion. The aerodynamic moment derivative
with respect to deck angle had the wrong sign for stability; the
bridge had not been analysed for this mode because the design
canon of the period treated wind as a static distributed load.
The case is treated in detail in chapter 6 of this volume on
vibrations (where the structural-dynamics view dominates) and in
the accident registry entry under tacoma-narrows-1940; the
chapter-13 frame is that the bluff-body aerodynamic moment on a
shallow plate-girder deck has a destabilising slope at the
operating wind speed, and the aerodynamic-stability analysis
that would have caught this was absent from the design canon.
Subsequent suspension-bridge design has incorporated
deck-section wind-tunnel testing as standard practice
\cite{acc:tacoma-narrows-1941}.

\subsection{Sonic boom and overflight regulation}

A supersonic aircraft generates a pressure signature on the ground
beneath its flight path: the \engterm{sonic boom}. The signature
is typically an N-wave (a leading shock followed by a smooth
expansion followed by a trailing shock) with overpressures of
the order of $50$ to $100\,\text{Pa}$ for typical supersonic
fighters at cruise altitude. The peak overpressure scales with
the aircraft's volume and weight and inversely with the cube root
of the cruise altitude.

The sonic-boom signature is the working basis of the FAA's
prohibition on civil supersonic flight over the continental
United States (current as of 2026), encoded in FAR 91.817. The
prohibition has shaped fifty years of civil aviation design:
Concorde flew only on transatlantic and a few transpacific routes,
and the next-generation supersonic transport designs (Boom
Overture, NASA X-59) target a much lower-amplitude shaped boom to
clear the regulatory threshold. The boom-design problem combines
the area-rule and shock-shape methods of supersonic aerodynamics
with the propagation-through-the-atmosphere problem of
acoustic-shock dynamics; it is one of the canonical research
problems of aeroacoustics.

\subsection{The drag-divergence trap}

A transport aircraft designed for cruise at $M = 0.80$ on a
conventional airfoil sees a sharp drag rise above the drag-
divergence Mach number $M_{DD} \approx 0.78$. A small operational
excursion above $M_{DD}$ (a tailwind, a higher-than-expected
cruise altitude, a temporary speed increase) produces a
substantial fuel-burn penalty: drag at $M = 0.84$ may be twice the
drag at $M = 0.80$. The discipline of operational efficiency is
to fly below $M_{DD}$ in cruise; the airline's flight-management
system computes the cost-optimal cruise Mach for each leg, balancing
fuel burn against flight time and against ATC speed constraints.

The introduction of the supercritical airfoil (Whitcomb, NASA
Langley, 1965) raised $M_{DD}$ from about $0.78$ to about $0.85$
for typical transport airfoils, enabling efficient cruise at
$M = 0.80$ to $0.83$ on supercritical-wing aircraft (the standard
since the $1970$s, including the Boeing 757/767/777/787 family and
the Airbus A300/A330/A340/A350 family). The supercritical airfoil
has a flatter upper surface and more aft camber, which delays the
shock formation and reduces the shock strength when it forms; the
working design principle is to spread the pressure recovery
across the rear of the airfoil rather than concentrating it at a
single shock.

\subsection{Inlet unstart}

A supersonic inlet runs a carefully arranged shock system that
decelerates the captured flow to a subsonic, high-pressure state at
the engine face. The arrangement is conditionally stable. If a
disturbance (a sudden change in angle of attack, a combustor pressure
spike, an excess of captured mass flow) pushes the terminal shock
forward of the throat, the shock cannot find a stable position on the
converging side of the throat: it is expelled violently out the front
of the inlet. The expulsion is \engterm{unstart}. The captured mass
flow collapses, the inlet drag jumps, a bow shock stands ahead of the
inlet lip, and on a podded engine the asymmetric thrust and pressure
loads can be large enough to threaten the airframe. The SR-71's
unstarts were severe enough to slam the crew sideways against the
canopy; its inlet control system spent much of its design budget on
keeping the shock train positioned and on restarting an unstarted
inlet by opening bypass doors and retracting the centre-body spike.
The unstart hazard is intrinsic to fixed-geometry mixed-compression
inlets: the same internal contraction that gives good pressure
recovery also makes the started state a narrow basin of attraction.
The design responses are variable geometry (moving spike or ramp),
bleed and bypass to widen the stable mass-flow range, and a fast
control law that detects the forward motion of the shock and reacts
before expulsion. A scramjet inlet faces the same physics at the
start of operation: if the internal contraction ratio exceeds the
Kantrowitz limit for the flight Mach number, the supersonic flow
cannot establish in the duct at all, and the inlet never starts.

\subsection{The diagnostic principle for compressible-flow failures}

The compressible-flow diagnostic question is, for any high-speed
configuration:

\begin{itemize}[leftmargin=*]
  \item What is the local Mach number? Subsonic, transonic,
    supersonic, hypersonic?
  \item Are shocks present? If so, where do they attach? What is
    their strength?
  \item Is there shock-induced boundary-layer separation? Is the
    separation steady or unsteady?
  \item Is the aerodynamic moment derivative with respect to the
    relevant structural mode stable or destabilising?
  \item What is the stagnation temperature at the local Mach
    number? Is the wall material rated for that temperature, with
    appropriate heat-flux margin?
  \item What is the wave-drag contribution? Has the area rule
    been applied?
\end{itemize}

A design that has answered all six is a compressibility-aware
design. A design that has not is in the regime where transonic
buffet, flutter, shock-induced stall, or thermal-protection-system
failure can convert an apparently routine flight into an envelope-
exceedance event.

\begin{principle}[Compressibility introduces three new failure modes]
Compressible flow has three failure modes absent from
incompressible flow: shock-induced boundary-layer separation
(producing buffet, lift loss, and drag rise), aeroelastic
instability driven by shock or unsteady supersonic pressure
distributions (flutter at transonic and supersonic speeds), and
aerodynamic heating from stagnation-point temperature rise (a
hypersonic regime issue). The design discipline against each is
specific: drag-divergence Mach number margin for the first,
flutter analysis across the envelope for the second,
thermal-protection-system design and verification for the third.
\end{principle}

\begin{archetype}[Transport, stability, scaling: the compressible-flow triad]
The compressible-flow archetype puts three relations to work. The
Mach number sets the regime (subsonic compressible, transonic,
supersonic, hypersonic); the isentropic relations connect the
local thermodynamic state to the local Mach number; the
Rankine-Hugoniot relations close the property jumps across shocks.
The archetype's failure modes are shock-induced separation,
aeroelastic flutter, and aerodynamic heating; each scales
differently with Mach number and each has a specific design
discipline. Every compressible-flow analysis works through these
three: regime, isentropic state, shock physics.
\end{archetype}

\begin{estimation}
\textbf{Estimate the wave drag of a transport aircraft at cruise
and compare to its skin-friction drag.}

\smallskip
\textit{Estimate, before reading on.}

\smallskip
A working estimate goes through the following moves. A typical
transport at cruise (Mach $0.80$, altitude $11\,\si{\kilo\meter}$,
$\rho \approx 0.36\,\si{\kilogram\per\meter\cubed}$, $a \approx
295\,\si{\meter\per\second}$, $U_\infty \approx 236\,\si{\meter
\per\second}$) carries a weight of about $2\times 10^6\,\si{
\newton}$ and has a wing area of about $300\,\si{\meter}^2$. The
lift coefficient at cruise is
\[
C_L = \frac{W}{\tfrac{1}{2}\rho U_\infty^2 S} \approx \frac{2\times
10^6}{0.5\times 0.36\times 236^2\times 300} \approx 0.66,
\]
in the typical $0.4$ to $0.7$ band.

The induced drag coefficient is
\[
C_{D,i} = \frac{C_L^2}{\pi e\,AR},
\]
with aspect ratio $AR \approx 9$ and Oswald efficiency $e \approx
0.85$:
\[
C_{D,i} \approx \frac{0.66^2}{\pi\times 0.85\times 9} \approx
0.018.
\]
The skin-friction drag coefficient is roughly
$C_{D,f} \approx 0.02$ to $0.03$ (from chapter 12's estimate).
The wave drag at cruise is the small quantity that the design
deliberately keeps under control; at cruise Mach below
drag-divergence, $C_{D,\text{wave}} \approx 0.002$ to $0.005$,
roughly $10$ to $20\,\%$ of the total drag. The total drag
coefficient is about $C_D \approx 0.04$, of which:

\begin{itemize}[leftmargin=*]
  \item Skin friction: $\sim 50$--$60\,\%$.
  \item Induced (vortex): $\sim 30$--$40\,\%$.
  \item Wave drag (at design cruise): $\sim 5$--$10\,\%$.
  \item Excrescence drag (gaps, antennas, etc.): $\sim 5\,\%$.
\end{itemize}

The cruise drag force is
\[
D = \tfrac{1}{2}\rho U_\infty^2 S \cdot C_D \approx 0.5\times 0.36
\times 236^2\times 300\times 0.04 \approx 120\,\text{kN}.
\]
A typical transport in cruise burns roughly $5\,\text{kg}$ of fuel
per second, with each kilogram providing about $43\,\text{MJ}$ of
energy; the cruise propulsive power is roughly $D\cdot U_\infty
\approx 120{,}000\times 236 \approx 28\,\text{MW}$, consistent with
a thrust-specific fuel consumption of about $0.6$ to $0.7\,\text{
kg/h/kg}$ at cruise. The figures hang together within a factor of
two.

The postmortem: the wave-drag share, while modest at design
cruise, doubles or triples if the aircraft strays a few hundredths
of Mach above drag-divergence. The discipline is to maintain a
small but nonzero margin below $M_{DD}$ and to size the wing for
the cruise condition's $C_L$ rather than for the maximum lift in
takeoff and landing (where high-lift devices provide the extra
$C_L$). The wave drag is the smallest term but the most
sensitive to operational choices; this sensitivity is what makes
the supercritical airfoil's $M_{DD}$ improvement valuable.
\end{estimation}

\section{Project}\pathtag{core}

\begin{project}[Estimate the wave drag of a passenger jet at cruise]
The reader performs an analytical estimate of the wave drag of a
passenger jet at cruise conditions, using the area rule, the
Ackeret thin-airfoil theory, and published geometric data. The
deliverable is a fifteen- to twenty-five-page engineering report.

\textbf{Track}. Paper / analytical, with optional comparison to
published CFD or wind-tunnel data. \textbf{Hazard class}: none.

\textbf{Aircraft selection}. The reader chooses a transport
aircraft with publicly available cross-sectional area distribution
and cruise specifications. Suitable candidates: Boeing 737-800,
Airbus A320, Boeing 787-8, Airbus A350-900, Concorde (for a
supersonic case). The cruise specifications (Mach number,
altitude, weight, wing area, aspect ratio) are taken from
manufacturer or airline data.

\textbf{Analysis}. The reader produces:

\begin{itemize}[leftmargin=*]
  \item A plot of the aircraft's cross-sectional area $A(x)$ as a
    function of longitudinal coordinate, including the fuselage,
    wings (projected), engines (projected), and tail surfaces. The
    plot may use scaling from public three-view drawings.
  \item An estimate of the slope $dA/dx$ and a calculation of the
    Sears-Haack wave drag integral $D_w = (\rho U_\infty^2/4\pi)
    \int\int A''(x) A''(x') \ln|x - x'| \,dx \,dx'$, or an
    approximation based on a representative cross-section
    waist-to-base ratio.
  \item For the wing, an Ackeret-theory estimate of the
    section wave drag $C_{D,\text{wave}} \approx 4\alpha^2/
    \sqrt{M_\infty^2 - 1}$ for a comparable supersonic wing, or
    a transonic-drag-rise estimate using the published drag-
    divergence Mach number.
  \item A comparison of the wave-drag estimate to the full cruise
    drag, computed from cruise thrust or from published $L/D$
    data.
\end{itemize}

\textbf{Reflection ($1{,}500$ to $3{,}000$ words)}. The reader
addresses:

\begin{itemize}[leftmargin=*]
  \item The wave-drag share of the total cruise drag, with the
    largest source of uncertainty in the estimate.
  \item Whether the aircraft's cross-sectional area distribution
    appears to follow the area rule, with a specific named
    geometric feature (e.g., the inboard fairing on a 787) that
    supports the conclusion.
  \item One design change that would reduce the wave drag and the
    operational cost (range, fuel burn, payload) the change would
    achieve.
\end{itemize}

\textbf{Deliverable}. The area-distribution plot, the wave-drag
calculation, the comparison to total drag, and the reflection.

\textbf{Time}. Four to six hours of geometric data collection,
four to eight hours of analysis, four to six hours of write-up.
\end{project}

\section{Exercises}

The exercises mix calculation, derivation, estimation, simulation,
design, diagnosis, failure analysis, and judgment.

\subsection*{Calculation}\pathtag{core}

\begin{exercise}
Compute the speed of sound in air at $T = 288\,\si{\kelvin}$ (sea
level) and at $T = 217\,\si{\kelvin}$ (typical cruise altitude),
both for $\gamma = 1.4$, $R = 287\,\si{\joule\per\kilogram\per
\kelvin}$.
\paragraph{Solution sketch.} $a = \sqrt{\gamma R T}$. At
$288\,\si{\kelvin}$: $a = \sqrt{1.4\times 287\times 288} =
\sqrt{1.156\times 10^{5}} = 340\,\si{\meter\per\second}$. At
$217\,\si{\kelvin}$: $a = \sqrt{1.4\times 287\times 217} =
\sqrt{8.72\times 10^{4}} = 295\,\si{\meter\per\second}$. The drop is
the $\sqrt{T}$ scaling: $\sqrt{217/288} = 0.868$, and $340\times
0.868 = 295$. The speed of sound depends only on temperature for an
ideal gas, so the altitude effect enters entirely through $T$.
\end{exercise}

\begin{exercise}
An aircraft flies at $250\,\si{\meter\per\second}$ at sea level and
at the same true airspeed at $11\,\si{\kilo\meter}$ altitude. Find
the Mach number in each case.
\paragraph{Solution sketch.} $M = u/a$. At sea level
($a = 340$): $M = 250/340 = 0.74$. At $11\,\si{\kilo\meter}$
($a = 295$): $M = 250/295 = 0.85$. The same true airspeed gives a
higher Mach number at altitude because the speed of sound has
dropped. The aircraft that is comfortably subsonic-compressible at
sea level is up against drag divergence at cruise altitude, which is
why the operating limit is posted as a Mach number, not an airspeed.
\end{exercise}

\begin{exercise}
For air ($\gamma = 1.4$) at $M = 2$, compute the ratios $T/T_0$,
$p/p_0$, $\rho/\rho_0$ using the isentropic relations.
\paragraph{Solution sketch.} $T_0/T = 1 + 0.2\times 4 = 1.8$, so
$T/T_0 = 0.556$. Then $p/p_0 = (T/T_0)^{\gamma/(\gamma-1)} =
0.556^{3.5} = 0.128$ and $\rho/\rho_0 = (T/T_0)^{1/(\gamma-1)} =
0.556^{2.5} = 0.230$. These match the $M = 2$ row of
\texttt{data/isentropic\_table.csv}. The static pressure at $M = 2$ is
about an eighth of the stagnation pressure, the steep pressure drop
that makes high-Mach pitot measurement and nozzle expansion possible.
\end{exercise}

\begin{exercise}
A converging-diverging nozzle has a throat area of $0.01\,\si{
\meter}^2$ and an exit area of $0.04\,\si{\meter}^2$. Find the
design exit Mach number assuming the throat is choked.
\paragraph{Solution sketch.} The throat is choked, so $A^* =
0.01\,\si{\meter}^2$ and $A_e/A^* = 4.0$. Inverting the area-Mach
relation on the supersonic branch (\texttt{code/nozzle.py}) gives
$M_e = 2.94$. The subsonic root of the same area ratio is
$M_e = 0.147$, the answer if the back pressure were too high to hold
supersonic flow; the design (matched, fully supersonic) case takes the
supersonic root. A useful check: $A/A^* = 4.0$ sits just below the
$M = 3$ value of $4.235$ in the table, consistent with $M_e \approx
2.94$.
\end{exercise}

\begin{exercise}
A normal shock occurs in air at upstream Mach number $M_1 = 2.5$.
Compute $M_2$, $p_2/p_1$, $\rho_2/\rho_1$, and $T_2/T_1$.
\paragraph{Solution sketch.} With $M_1^2 = 6.25$ and $\gamma = 1.4$:
$M_2^2 = (0.4\times 6.25 + 2)/(2.8\times 6.25 - 0.4) = 4.5/17.1 =
0.263$, so $M_2 = 0.513$. Pressure ratio $p_2/p_1 = 1 +
(2.8/2.4)(6.25 - 1) = 1 + 1.1667\times 5.25 = 7.13$. Density ratio
$\rho_2/\rho_1 = 2.4\times 6.25/(0.4\times 6.25 + 2) = 15/4.5 = 3.33$.
Temperature ratio $T_2/T_1 = (p_2/p_1)/(\rho_2/\rho_1) = 7.13/3.33 =
2.14$. These match the $M_1 = 2.5$ row of
\texttt{data/normal\_shock\_table.csv}. The flow leaves subsonic, the
pressure rises sevenfold, and the temperature more than doubles, the
signature of a strong shock.
\end{exercise}

\begin{exercise}
An oblique shock turns a flow at $M_1 = 3$ through a deflection
angle of $15^{\circ}$. Find the weak-solution shock angle $\beta$,
the downstream Mach number normal to the shock, and the downstream
Mach number along the flow direction.
\paragraph{Solution sketch.} From the $\theta$-$\beta$-$M$ relation
the weak root is $\beta = 32.2^{\circ}$
(\texttt{code/oblique\_shock.py}). The upstream normal component is
$M_{n1} = M_1\sin\beta = 3\times\sin 32.2^{\circ} = 1.60$. Applying
the normal-shock relation to $M_{n1} = 1.60$ gives $M_{n2} = 0.668$.
The downstream Mach number along the deflected flow is $M_2 =
M_{n2}/\sin(\beta - \theta) = 0.668/\sin 17.2^{\circ} = 2.26$. The
flow is turned and compressed but stays supersonic, the property that
lets a multi-ramp inlet stack weak shocks for efficient pressure
recovery.
\end{exercise}

\begin{exercise}
For a flat-plate airfoil at $3^{\circ}$ angle of attack in a
supersonic stream at $M_\infty = 2$, compute the lift coefficient
and the wave-drag coefficient using Ackeret theory. Then compute
$L/D$.
\paragraph{Solution sketch.} $\alpha = 3^{\circ} = 0.0524\,\text{rad}$,
$\beta = \sqrt{M_\infty^2 - 1} = \sqrt{3} = 1.732$. Then
$C_L = 4\alpha/\beta = 4\times 0.0524/1.732 = 0.121$ and
$C_{D,\text{wave}} = 4\alpha^2/\beta = 4\times 0.0524^2/1.732 =
0.0063$. The lift-to-drag ratio is $L/D = C_L/C_{D,\text{wave}} =
1/\alpha = 19.1$ (since $C_L/C_D = (4\alpha/\beta)/(4\alpha^2/\beta) =
1/\alpha$). The result $L/D = 1/\alpha$ says the supersonic flat plate
is most efficient at small angle of attack, but its absolute $L/D$
still trails a subsonic transport because $C_L$ itself is held down by
the $1/\beta$ factor.
\end{exercise}

\begin{exercise}
The Mach angle at $M = 3$ in air is $\mu = \sin^{-1}(1/M)$.
Compute it, and compute the corresponding angle for $M = 1.5$ and
$M = 5$.
\paragraph{Solution sketch.} $\mu = \sin^{-1}(1/M)$. At $M = 3$:
$\sin^{-1}(0.333) = 19.5^{\circ}$. At $M = 1.5$: $\sin^{-1}(0.667) =
41.8^{\circ}$. At $M = 5$: $\sin^{-1}(0.2) = 11.5^{\circ}$. The cone
narrows as the Mach number rises: a faster body outruns its sound by
more, so the wave it leaves trails at a shallower angle. The Mach
angle is the lower limit of the oblique-shock angle (the limit of an
infinitesimally weak shock).
\end{exercise}

\begin{exercise}
Compute the stagnation temperature at the nose of a vehicle
travelling at $M = 8$ at altitude where $T_\infty = 220\,\si{
\kelvin}$. State whether the result is within the validity range
of the ideal-gas assumption.
\paragraph{Solution sketch.} $T_0 = T_\infty(1 + 0.2 M^2) =
220\times(1 + 0.2\times 64) = 220\times 13.8 = 3036\,\si{\kelvin}$.
The perfect-gas, constant-$\gamma$ result is not trustworthy at this
temperature: oxygen begins to dissociate around $2000$ to
$4000\,\si{\kelvin}$, so the real stagnation enthalpy goes into
breaking molecular bonds rather than raising temperature, and the
actual gas temperature is well below $3036\,\si{\kelvin}$. The number
flags the regime where real-gas chemistry must replace the ideal-gas
relation, not a literal temperature.
\end{exercise}

\subsection*{Derivation}\pathtag{standard}

\begin{exercise}
Derive the speed of sound for an ideal gas, $a^2 = \gamma R T$,
from a control-volume analysis across an infinitesimal disturbance.
\paragraph{Solution sketch.} In the frame of the disturbance, mass
conservation gives $\rho a = (\rho + d\rho)(a - du)$, which to first
order is $\rho\,du = a\,d\rho$. Momentum across the same control
volume gives $p - (p + dp) = \rho a(a - du) - \rho a^2$, that is
$dp = \rho a\,du$. Eliminating $du$ between the two yields $a^2 =
dp/d\rho$. The disturbance is small-amplitude and so isentropic, so
the derivative is at constant entropy. For an ideal gas with
$p\rho^{-\gamma} = \text{const}$, $dp/d\rho = \gamma p/\rho = \gamma
RT$. Hence $a = \sqrt{\gamma RT}$.
\end{exercise}

\begin{exercise}
Derive the isentropic stagnation relations $T_0/T = 1 +
(\gamma - 1)M^2/2$, $p_0/p = [T_0/T]^{\gamma/(\gamma - 1)}$ from
the energy equation and the isentropic condition.
\paragraph{Solution sketch.} The steady adiabatic energy equation is
$c_p T + \tfrac{1}{2}u^2 = c_p T_0$. Divide by $c_p T$:
$1 + u^2/(2 c_p T) = T_0/T$. Substitute $c_p = \gamma R/(\gamma - 1)$
and $u^2 = M^2 a^2 = M^2 \gamma R T$, so $u^2/(2 c_p T) = M^2\gamma RT
(\gamma - 1)/(2\gamma RT) = (\gamma - 1)M^2/2$. Hence $T_0/T = 1 +
(\gamma - 1)M^2/2$. For the pressure, the isentropic relation
$p/\rho^\gamma = \text{const}$ with $p = \rho RT$ gives $T \propto
p^{(\gamma-1)/\gamma}$, so $p_0/p = (T_0/T)^{\gamma/(\gamma - 1)}$, and
the density relation follows from $p = \rho RT$ as $\rho_0/\rho =
(T_0/T)^{1/(\gamma - 1)}$.
\end{exercise}

\begin{exercise}
Derive the area-Mach relation $A/A^*$ from the continuity equation
and the isentropic relations.
\paragraph{Solution sketch.} Continuity along the duct is $\rho u A =
\rho^* u^* A^*$, so $A/A^* = (\rho^* u^*)/(\rho u)$. Write $u = M a =
M\sqrt{\gamma RT}$ and $u^* = a^* = \sqrt{\gamma RT^*}$, so
$u^*/u = (1/M)\sqrt{T^*/T}$. Combine $\rho^*/\rho =
(\rho^*/\rho_0)(\rho_0/\rho)$ and $T^*/T = (T^*/T_0)(T_0/T)$ using the
critical and isentropic relations, $T^*/T_0 = 2/(\gamma + 1)$ and
$T_0/T = 1 + (\gamma-1)M^2/2$. Collecting the powers gives
\[
\frac{A}{A^*} = \frac{1}{M}\left[\frac{2}{\gamma + 1}\left(1 +
\frac{\gamma - 1}{2}M^2\right)\right]^{(\gamma + 1)/[2(\gamma - 1)]}.
\]
Setting $M = 1$ returns $A/A^* = 1$, the throat condition that fixes
$A^*$.
\end{exercise}

\begin{exercise}
Derive the Rankine-Hugoniot pressure ratio across a normal shock,
$p_2/p_1 = 1 + 2\gamma(M_1^2 - 1)/(\gamma + 1)$, from the
conservation laws across a thin control volume.
\paragraph{Solution sketch.} Mass $\rho_1 u_1 = \rho_2 u_2$ and
momentum $p_1 + \rho_1 u_1^2 = p_2 + \rho_2 u_2^2$ combine to
$p_2 - p_1 = \rho_1 u_1^2 - \rho_2 u_2^2 = \rho_1 u_1(u_1 - u_2)$.
Write $\rho u^2 = \gamma p M^2$ (since $u^2 = M^2\gamma p/\rho$), so the
momentum equation becomes $p_1(1 + \gamma M_1^2) = p_2(1 + \gamma
M_2^2)$, giving $p_2/p_1 = (1 + \gamma M_1^2)/(1 + \gamma M_2^2)$.
Substituting the downstream Mach relation $M_2^2 = [(\gamma - 1)M_1^2
+ 2]/[2\gamma M_1^2 - (\gamma - 1)]$ and simplifying yields
$p_2/p_1 = 1 + \dfrac{2\gamma}{\gamma + 1}(M_1^2 - 1)$. At $M_1 = 1$
the ratio is $1$ (no shock); it grows quadratically in $M_1$.
\end{exercise}

\begin{exercise}
Show that the entropy increase across a weak normal shock scales
as $(M_1^2 - 1)^3$ in the limit $M_1 \to 1^+$, by expanding the
shock-table relations.
\paragraph{Solution sketch.} Let $m = M_1^2 - 1$, small. The entropy
change for an ideal gas is $\Delta s = c_v \ln[(p_2/p_1)(\rho_1/
\rho_2)^\gamma]$. Expand $p_2/p_1 = 1 + \frac{2\gamma}{\gamma+1}m$ and
$\rho_2/\rho_1 = 1 + \frac{2}{\gamma+1}m + O(m^2)$ as power series in
$m$. Form $X = (p_2/p_1)(\rho_1/\rho_2)^\gamma$ and expand $\ln X$.
The first-order terms cancel (a sonic ``shock'' is isentropic) and the
second-order terms cancel as well; the leading surviving term is cubic,
$\Delta s \approx \frac{2\gamma(\gamma - 1)}{3(\gamma + 1)^2}c_v\,m^3 +
O(m^4)$. So $\Delta s \propto (M_1^2 - 1)^3$. The cubic vanishing is
why weak shocks are nearly reversible and why the linearised
(isentropic) acoustic theory holds right up to $M_1 = 1$.
\end{exercise}

\begin{exercise}
Derive the Ackeret thin-airfoil result $C_p = 2\theta/\sqrt{M_
\infty^2 - 1}$ from the linearised supersonic pressure equation.
\paragraph{Solution sketch.} Linearised supersonic flow obeys the
wave equation $(M_\infty^2 - 1)\phi_{xx} - \phi_{yy} = 0$ for the
perturbation potential $\phi$. Its general solution is $\phi =
f(x - \beta y)$ with $\beta = \sqrt{M_\infty^2 - 1}$: disturbances
travel along the right-running Mach lines only. The surface
tangency condition is $\phi_y = U_\infty\,\theta(x)$ at the wall,
where $\theta$ is the local surface slope. Differentiating the
solution, $\phi_y = -\beta f'$ and $\phi_x = f'$, so $\phi_x =
-\phi_y/\beta = -U_\infty\theta/\beta$. The linearised pressure
coefficient is $C_p = -2\phi_x/U_\infty = 2\theta/\beta =
2\theta/\sqrt{M_\infty^2 - 1}$. Surfaces inclined into the flow
($\theta > 0$) are compressed; those inclined away are expanded.
Integrating over a flat plate at angle $\alpha$ recovers $C_L =
4\alpha/\beta$ and $C_{D,\text{wave}} = 4\alpha^2/\beta$.
\end{exercise}

\subsection*{Estimation}\pathtag{core}

\begin{exercise}
Estimate the temperature rise at the stagnation point of a
commercial jet's nose at Mach $0.85$. State the assumed
freestream conditions.
\paragraph{Reference estimate.} Take cruise $T_\infty = 217\,\si{
\kelvin}$. The stagnation rise is $\Delta T = T_\infty(\gamma -
1)M^2/2 = 217\times 0.2\times 0.7225 = 31\,\si{\kelvin}$, giving
$T_0 = 248\,\si{\kelvin}$. The recovery temperature on an actual
(turbulent-boundary-layer) surface is a recovery factor
$r \approx 0.89$ times this rise, about $28\,\si{\kelvin}$. The nose
skin runs roughly $25$ to $30\,\si{\kelvin}$ above ambient, the order
of magnitude that matters for de-icing logic and window thermal
design but not for material limits at this Mach number.
\end{exercise}

\begin{exercise}
Estimate the size and shape of the Mach cone behind a supersonic
fighter jet flying at Mach $2$ at $10\,\si{\kilo\meter}$ altitude.
State the boom signature width at the ground.
\paragraph{Reference estimate.} The Mach angle is $\mu =
\sin^{-1}(1/2) = 30^{\circ}$. The cone trails behind the aircraft at
$30^{\circ}$ to the flight path. The boom footprint on the ground is
the intersection of this cone with the surface. For level flight the
carpet half-width scales with altitude: the lateral spread is roughly
$h\tan(90^{\circ} - \mu) = h\cot\mu$ on each side, so for $h =
10\,\si{\kilo\meter}$ and $\mu = 30^{\circ}$ the carpet is about
$2\times 10\times\cot 30^{\circ} = 2\times 10\times 1.73 \approx
35\,\si{\kilo\meter}$ wide. Every observer in that strip hears the
boom as the aircraft passes; the boom is a moving line, not a
one-time event at the moment of going supersonic.
\end{exercise}

\begin{exercise}
Estimate the choked mass flow rate through a $50\,\si{\milli\meter}$
diameter nozzle with upstream stagnation conditions $p_0 = 500\,
\text{kPa}$, $T_0 = 400\,\si{\kelvin}$ (air).
\paragraph{Reference estimate.} The throat area is $A^* = \pi
(0.025)^2 = 1.96\times 10^{-3}\,\si{\meter}^2$. The choked mass flux
(\texttt{code/nozzle.py}) is $\dot m/A^* = (p_0/\sqrt{T_0})\sqrt{
\gamma/R}\,(2/(\gamma+1))^{(\gamma+1)/[2(\gamma-1)]}$. With
$p_0 = 5\times 10^{5}$, $T_0 = 400$, $\gamma = 1.4$, $R = 287$, the
flux is about $1010\,\si{\kilogram\per\second}/\si{\meter}^2$. Then
$\dot m = 1010\times 1.96\times 10^{-3} \approx 1.98\,\si{\kilogram
\per\second}$. Because the nozzle is choked, this rate is fixed by the
upstream stagnation state alone; lowering the downstream pressure
further does not increase it, which is why a choked orifice is a
stable flow standard.
\end{exercise}

\begin{exercise}
Estimate the wave drag of a Concorde-class supersonic transport
($M = 2$ cruise) and compare to its skin-friction drag at the
same cruise condition.
\paragraph{Reference estimate.} At $M = 2$, $18\,\si{\kilo\meter}$
cruise, take $\rho \approx 0.12\,\si{\kilogram\per\meter\cubed}$,
$a \approx 295\,\si{\meter\per\second}$, $U_\infty \approx 590\,\si{
\meter\per\second}$, wing area $S \approx 360\,\si{\meter}^2$, weight
$\approx 1.5\times 10^{6}\,\si{\newton}$. Dynamic pressure $q =
\tfrac{1}{2}\rho U^2 \approx 0.5\times 0.12\times 590^2 \approx
20.9\,\text{kPa}$. Cruise $C_L = W/(qS) \approx 1.5\times 10^{6}/
(20900\times 360) \approx 0.20$. A slender supersonic transport runs
total $C_D \approx 0.012$ at $L/D \approx 7$. Of that, wave drag
(volume plus lift-dependent) and skin friction are comparable, each on
the order of $C_D \approx 0.004$ to $0.005$, with induced drag making
up the rest. So wave drag is roughly $35$ to $45\,\%$ of cruise drag,
a far larger share than the $5$ to $10\,\%$ of a subsonic transport;
this is the cost of supersonic cruise.
\end{exercise}

\begin{exercise}
Estimate the heat flux to the leading edge of a Space-Shuttle-class
re-entry vehicle at $M = 25$ at $80\,\si{\kilo\meter}$ altitude.
State the assumed nose radius and atmospheric density.
\paragraph{Reference estimate.} Take $\rho_\infty \approx 2\times
10^{-5}\,\si{\kilogram\per\meter\cubed}$ at $80\,\si{\kilo\meter}$,
$U_\infty \approx 7\,\si{\kilo\meter\per\second}$, effective nose
(leading-edge) radius $R_n \approx 0.3\,\si{\meter}$. The Sutton-Graves
scaling $q_w = K\sqrt{\rho_\infty/R_n}\,U_\infty^3$ with $K \approx
1.7\times 10^{-4}$ in SI gives $q_w \approx 1.7\times 10^{-4}\times
\sqrt{2\times 10^{-5}/0.3}\times(7000)^3 \approx 1.7\times 10^{-4}
\times 8.2\times 10^{-3}\times 3.43\times 10^{11} \approx 4.8\times
10^{5}\,\si{\watt\per\meter\squared}$, around $0.5\,\si{\mega\watt\per
\meter\squared}$. This is the order of magnitude the reinforced
carbon-carbon leading edge had to survive; the cube on velocity is why
re-entry heating, not launch, sets the leading-edge material choice.
\end{exercise}

\subsection*{Simulation}\pathtag{mastery}

\begin{exercise}
Implement a one-dimensional finite-volume solver for the Euler
equations on a converging-diverging nozzle. Verify the choked-
flow conditions and the supersonic acceleration through the
diverging section against the area-Mach relation.
\paragraph{Solution sketch.} Discretise the quasi-one-dimensional
Euler equations $\partial_t(A\mathbf U) + \partial_x(A\mathbf F) =
\mathbf S$ with the area-source term $\mathbf S = (0, p\,dA/dx, 0)^T$.
Use a Rusanov or HLLC flux at cell faces, a fixed throat geometry, and
characteristic boundary conditions: subsonic stagnation inflow,
supersonic (extrapolated) or back-pressure outflow. March to steady
state. Verification: the converged throat Mach number is $1.000$ to
plotting accuracy, the mass flux is constant along the duct to within
the scheme's conservation error, and the centreline Mach number
matches the supersonic root of $A/A^*(x)$ from \texttt{code/nozzle.py}
at each station. A back pressure between the subsonic and design
values should produce a captured normal shock whose position matches
the analytic shock-in-nozzle calculation.
\end{exercise}

\begin{exercise}
Implement an oblique-shock $\theta$-$\beta$-$M$ solver. Verify
against the standard shock tables; identify the deflection
angle at which the attached-shock solution disappears at given
$M_1$.
\paragraph{Solution sketch.} The reference implementation is
\texttt{code/oblique\_shock.py}: evaluate $\theta(\beta)$ from the
$\theta$-$\beta$-$M$ relation, scan $\beta$ from the Mach angle to
$90^{\circ}$ to find $\theta_{\max}$, and bisect each branch for a
target deflection. Verification at $M_1 = 2$, $\theta = 10^{\circ}$
gives a weak $\beta \approx 39.3^{\circ}$, matching the tables. The
attached-shock solution disappears at $\theta_{\max}$: for $M_1 = 2$,
$\theta_{\max} \approx 22.97^{\circ}$ at $\beta \approx 64.7^{\circ}$;
for $M_1 = 3$, $\theta_{\max} \approx 34.07^{\circ}$. Beyond
$\theta_{\max}$ the curve has no real root and the shock detaches.
\end{exercise}

\begin{exercise}
Run a simple two-dimensional Euler simulation of supersonic flow
past a wedge at $M_\infty = 2$, deflection angle $15^{\circ}$.
Compare the computed shock angle to the analytical
$\theta$-$\beta$-$M$ result.
\paragraph{Solution sketch.} Set up a structured grid over the wedge,
impose $M_\infty = 2$ supersonic inflow and slip-wall on the wedge
surface, and march a two-dimensional finite-volume Euler scheme to
steady state. Measure the shock angle from the density-gradient
contour. The analytic weak-shock root at $M_1 = 2$, $\theta =
15^{\circ}$ is $\beta \approx 45.3^{\circ}$
(\texttt{code/oblique\_shock.py}); a converged second-order solver on
a sufficiently fine grid reproduces it to within about a degree, with
the discrepancy set by numerical shock smearing across two to four
cells. Refining the grid sharpens the shock and tightens the
agreement.
\end{exercise}

\subsection*{Design}\pathtag{standard}

\begin{exercise}
Design a converging-diverging nozzle to expand air from $p_0 =
2\,\text{MPa}$, $T_0 = 600\,\si{\kelvin}$ to exit Mach $M_e = 3$
with mass flow rate $1\,\si{\kilogram\per\second}$. Specify the
throat and exit areas and the divergence half-angle.
\paragraph{Solution sketch.} Throat is choked, so size $A^*$ from the
choked-flux relation: $\dot m/A^* = (p_0/\sqrt{T_0})\sqrt{\gamma/R}
(2/(\gamma+1))^{3}$ for air $= (2\times 10^{6}/\sqrt{600})\times
0.0684\times\dots \approx 4040\,\si{\kilogram\per\second}/\si{\meter}^2$
(\texttt{code/nozzle.py}). Then $A^* = 1/4040 = 2.48\times 10^{-4}\,
\si{\meter}^2$ (throat diameter $\approx 17.8\,\si{\milli\meter}$).
For $M_e = 3$, $A_e/A^* = 4.235$, so $A_e = 1.05\times 10^{-3}\,
\si{\meter}^2$ (exit diameter $\approx 36.6\,\si{\milli\meter}$). A
conical divergence half-angle of $12$ to $15^{\circ}$ is the standard
compromise: smaller angles add length and friction, larger angles
risk flow separation and divergence (cosine) thrust loss. A bell
(method-of-characteristics) contour shortens the nozzle and removes
the divergence loss.
\end{exercise}

\begin{exercise}
Design a thermal-protection-system layer for a hypersonic re-entry
vehicle nose with $M = 20$ peak entry conditions and $80\,\si{
\kilo\meter}$ peak-heating altitude. Specify the material choice
(ablative, ceramic, active), the nose radius, and the expected
stagnation heat flux.
\paragraph{Solution sketch.} Choose a blunt nose to spread the load:
$R_n = 0.5\,\si{\meter}$. At $80\,\si{\kilo\meter}$, $\rho_\infty
\approx 2\times 10^{-5}\,\si{\kilogram\per\meter\cubed}$, $U_\infty
\approx 6\,\si{\kilo\meter\per\second}$. Sutton-Graves
$q_w = 1.7\times 10^{-4}\sqrt{\rho_\infty/R_n}U_\infty^3 \approx
1.7\times 10^{-4}\times 6.3\times 10^{-3}\times 2.16\times 10^{11}
\approx 2.3\times 10^{5}\,\si{\watt\per\meter\squared}$
($\sim 0.23\,\si{\mega\watt\per\meter\squared}$). For a single steep
entry this favours an ablator (carbon-phenolic or PICA), which carries
heat away by pyrolysis and is single-use; a low-heating, many-cycle
profile favours a reusable ceramic. Size the ablator thickness from
the integrated heat load over the trajectory, not the peak flux, with
a recession allowance plus an insulating char layer that keeps the
bondline below the structural limit (around $450\,\si{\kelvin}$).
\end{exercise}

\begin{exercise}
Design the wing planform and section for a transport aircraft
with cruise Mach $0.83$, range $10{,}000\,\si{\kilo\meter}$, and
payload $30{,}000\,\si{\kilogram}$. Specify the wing area,
sweep, aspect ratio, and section type (conventional or
supercritical). Justify the choice in terms of drag-divergence
margin.
\paragraph{Solution sketch.} A conventional section gives $M_{DD}
\approx 0.74$, below the $0.83$ cruise, so the wing would sit deep in
the drag rise; a supercritical section with $M_{DD} \approx 0.84$ to
$0.86$ is required to leave any margin. Add quarter-chord sweep
$\Lambda \approx 25$ to $30^{\circ}$, which raises the effective
drag-divergence Mach number further through the $\cos\Lambda$ factor
on the normal-component Mach number. Aspect ratio $\approx 9$ to $10$
balances induced drag against structural weight. Size $S$ from the
cruise $C_L$: with cruise weight $\sim 7\times 10^{5}\,\si{\newton}$,
$q \approx 11\,\text{kPa}$ at $M = 0.83$, $11\,\si{\kilo\meter}$, and
target $C_L \approx 0.5$, $S = W/(qC_L) \approx 130\,\si{\meter}^2$.
The justification is the margin: supercritical plus sweep keeps cruise
below $M_{DD}$ so the drag-divergence penalty stays off the cruise
point.
\end{exercise}

\begin{exercise}
Design a supersonic inlet for a Mach $2$ cruise vehicle. Specify
the inlet geometry (single-cone, double-cone, mixed-compression),
the throat-area-to-capture-area ratio, and the cone angles.
Estimate the stagnation-pressure recovery.
\paragraph{Solution sketch.} At $M = 2$ an external-compression inlet
with one or two oblique shocks plus a terminal normal shock beats a
single pitot (normal-shock) inlet, whose recovery is only $p_{02}/
p_{01} = 0.72$. Use a single $10^{\circ}$ ramp: the oblique shock at
$M_1 = 2$, $\theta = 10^{\circ}$ has $\beta = 39.3^{\circ}$, $M_{n1} =
M_1\sin\beta = 1.27$, so $p_{02}/p_{01}$ across it is about $0.987$ and
the downstream Mach number is about $1.64$. A terminal normal shock at
$M = 1.64$ recovers another $p_{02}/p_{01} \approx 0.88$, for an
overall recovery $\approx 0.987\times 0.88 \approx 0.87$, well above
the single-normal-shock value. Two ramps would push recovery past
$0.9$. Set the throat-to-capture ratio from the post-shock area-Mach
state and add bleed and bypass to manage the started/unstarted
boundary.
\end{exercise}

\subsection*{Diagnosis and reverse engineering}\pathtag{core}

\begin{exercise}
A supersonic wind-tunnel test section shows oblique shocks
appearing in the working region rather than a uniform supersonic
flow. Identify three possible causes (nozzle misalignment, model
blockage, boundary-layer transition) and the diagnostic that
distinguishes them.
\paragraph{Solution sketch.} (1) Nozzle contour error or wall
misalignment leaves residual waves that the method-of-characteristics
design should have cancelled; the diagnostic is an empty-tunnel
schlieren image showing the same wave pattern with no model present.
(2) Model blockage chokes the test section and throws shocks off the
model and walls; the diagnostic is that the waves scale with model
size and vanish when a smaller model is fitted, and the diffuser
struggles to swallow the flow. (3) Boundary-layer transition or
separation on the nozzle wall thickens the displacement surface and
launches weak waves; the diagnostic is oil-flow or hot-film evidence of
the transition location and waves that move with Reynolds number. Run
the empty-tunnel check first: it separates hardware faults from
model-induced ones.
\end{exercise}

\begin{exercise}
A transport aircraft cruising at $M = 0.84$ shows cabin vibration
not present at $M = 0.80$. Identify the most probable cause and
the operational response.
\paragraph{Solution sketch.} The most probable cause is transonic
buffet onset: at $M = 0.84$ the upper-surface shock has strengthened
past the point where it separates the boundary layer, and the
oscillating separated region shakes the wing at structural
frequencies. The signature is that the vibration appears above a
threshold Mach number and worsens with altitude and weight (both raise
the required $C_L$ and the shock strength). The operational response is
to reduce Mach below the buffet boundary, by descending to a denser
altitude, slowing, or reducing weight as fuel burns; the flight
management system's cost-index cruise Mach already targets a margin to
buffet onset.
\end{exercise}

\begin{exercise}
A rocket-engine nozzle running on the test stand shows visible
diamond shock patterns in the exhaust plume. Identify the
operating condition (overexpanded, underexpanded, or design) and
the relationship to the ambient pressure.
\paragraph{Solution sketch.} Shock diamonds (Mach diamonds) mean the
exit pressure does not match the ambient pressure, so the plume
oscillates through a periodic train of shocks and expansions. A
high-altitude-design nozzle run at sea level is overexpanded
($p_e < p_b$): the plume is compressed back to ambient through oblique
shocks at the lip, which reflect and form the diamonds. As ambient
pressure falls toward the design value the diamonds weaken and vanish
at the matched (design) condition; at still lower ambient the nozzle
becomes underexpanded ($p_e > p_b$) and the plume expands and widens
through Prandtl-Meyer fans. On a sea-level test stand the visible
diamonds confirm overexpansion of a nozzle sized for altitude.
\end{exercise}

\begin{exercise}
A scramjet inlet on a hypersonic test article is failing to start:
the supersonic flow does not establish in the combustor. Identify
the most probable cause (inlet contraction ratio too high, throat
blockage, boundary-layer separation) and the diagnostic.
\paragraph{Solution sketch.} The most probable cause is an internal
contraction ratio above the Kantrowitz limit for the flight Mach
number: the duct cannot swallow the starting normal shock, so a bow
shock stands at the inlet face and the flow never goes supersonic
inside. The diagnostic is to compare the geometric contraction ratio
to the Kantrowitz value at the test Mach number; if it exceeds the
limit, the inlet is unstartable as built. Secondary causes are
combustor back-pressure (thermal choking) propagating forward and
shock-induced boundary-layer separation blocking the throat; their
diagnostics are wall-pressure traces showing the upstream-running
disturbance and schlieren of the separation. The design fixes are
variable geometry, a starting bleed, or overspeeding to start then
throttling back.
\end{exercise}

\subsection*{Failure analysis}\pathtag{core}

\begin{exercise}
The original Tacoma Narrows Bridge collapse
\cite{acc:tacoma-narrows-1941} is treated in chapter 6 of this volume as a
self-excited vibration and in this chapter's section 13.8 as a
bluff-body aerodynamic stability failure. Argue, in $200$ to
$300$ words, what the aerodynamic-moment slope vs.\ deck-angle
relation tells the designer, and what wind-tunnel test or
analytical check would have caught the destabilising slope at
the operating wind speed.
\paragraph{Model-answer sketch.} A strong answer states that the sign
of $dC_M/d\alpha$ (aerodynamic moment coefficient versus deck angle)
decides stability: a positive slope feeds energy into a twisting
oscillation, so an angular displacement generates a moment that
increases the displacement. A shallow plate-girder deck is a bluff
body whose separated flow gives this destabilising slope over a band
of angles. The check that would have caught it is a sectional
wind-tunnel test of a scale deck model, measuring the unsteady moment
and the flutter (or critical) wind speed, or an analytical flutter
calculation using measured aerodynamic derivatives. The design canon
of the period treated wind as a static load, so no such test was run.
The answer should connect the registry's torsional-flutter mechanism
to the present chapter's frame: the deck-section aerodynamic moment
derivative had the wrong sign at the operating wind speed, and
sectional testing is now standard practice precisely to measure that
derivative before construction.
\end{exercise}

\begin{exercise}
The Columbia accident \cite{acc:caib-columbia} involved a
re-entry-heating failure after thermal-protection damage during
ascent. Argue, in $200$ to $300$ words, what the stagnation-heat-
flux scaling tells the designer about the consequences of a
leading-edge breach, and what design-time and operations-time
discipline the accident has placed on the human spaceflight
programme.
\paragraph{Model-answer sketch.} A strong answer uses the
stagnation-heat-flux scaling $q_w \propto \sqrt{\rho_\infty/R_n}
\,U_\infty^3$. A leading-edge breach destroys the blunt-body
protection locally: the effective radius at the breach collapses, so
$q_w$ spikes (the $1/\sqrt{R_n}$ factor), and the cube on velocity
means the re-entry flux dwarfs anything seen on ascent. Hot shock-layer
gas then enters the wing structure, which has no thermal margin, and
the wing fails inside the high-temperature flow. The discipline the
accident imposed: at design time, the leading-edge thermal-protection
must tolerate or be inspectable for impact damage, and debris-shedding
sources must be controlled; at operations time, on-orbit imaging and
inspection of the leading edge before re-entry, debris-strike tracking
during ascent, and a defined repair or rescue contingency. The answer
should align with the registry entry under columbia-2003 and keep the
chapter frame: the heat-flux scaling makes a small geometric breach a
catastrophic thermal event.
\end{exercise}

\begin{exercise}
A supersonic transport aircraft (Concorde, retired 2003)
experienced unique flutter modes not present on subsonic
transports. Identify, in $200$ to $300$ words, what
compressibility-related phenomena (shock-induced separation,
transonic-buzz, supersonic flutter) had to be analysed in the
Concorde design, and what the surviving design records show
about the analysis discipline of the period.
\paragraph{Model-answer sketch.} A strong answer notes that a
supersonic transport spends its flight crossing every regime: it
accelerates through the transonic drag rise (shock-induced separation,
buzz) on the way to $M = 2$ cruise, where supersonic flutter modes
(panel flutter on skin panels under external supersonic pressure, and
shock or single-degree-of-freedom torsional flutter) appear that a
subsonic transport never sees. The slender delta wing was chosen partly
because it holds an attached shock and stays stiff in torsion. The
analysis had to span the whole envelope, coupling unsteady aerodynamics
to structural modes at each Mach number, with extensive wind-tunnel and
flight-test verification given the limited computation of the 1960s.
The answer should mention the era's reliance on measured aerodynamic
derivatives and ground vibration testing rather than modern coupled
CFD-structures, and frame Concorde as the case where the full
compressibility failure-mode set had to be cleared at once.
\end{exercise}

\subsection*{Judgment}\pathtag{standard}

\begin{exercise}
A senior engineer argues that civil supersonic transport will
return only when the sonic-boom signature is reduced below the
regulatory threshold. Write a one-paragraph reply that engages
with both the aerodynamic-design problem and the regulatory
problem.
\paragraph{Model-answer sketch.} A strong reply agrees on the
direction while separating the two problems. Aerodynamically, the boom
amplitude depends on the longitudinal distribution of volume and lift,
so shaped-boom design (area-rule tailoring, lifting-surface placement)
can flatten the N-wave into a lower-amplitude signature, the goal of
the X-59 demonstrator. The regulatory threshold is not yet written in
boom-amplitude terms: FAR 91.817 bans the act of flying supersonic
over land, not a measured overpressure, so even a quiet-boom aircraft
needs a rule change to a noise-based standard before it can operate
overland. The reply should conclude that the technical and regulatory
problems are coupled but distinct: a successful low-boom demonstration
is the evidence the rulemaking needs, and neither half on its own
returns civil supersonic flight over land.
\end{exercise}

\begin{exercise}
A junior engineer asks whether the area rule applies at subsonic
speeds. Write a one-paragraph reply that states the regime in
which the area rule is most powerful and the regime in which it
ceases to be the dominant design consideration.
\paragraph{Model-answer sketch.} A strong reply explains that the area
rule governs wave drag, which exists only where shocks form, so it is
most powerful in the transonic regime (roughly $M = 0.9$ to $1.2$),
where smoothing the cross-sectional area distribution measurably cuts
the transonic drag rise. At low subsonic speeds there is no wave drag,
so the area rule ceases to be the dominant consideration; the drag
budget there is skin friction and induced drag, and area smoothing
buys little. At supersonic cruise the related Sears-Haack slender-body
result for minimum wave drag takes over, a generalisation of the same
idea. So the area rule is a transonic-design tool, decisive near
$M = 1$ and largely irrelevant well below it.
\end{exercise}

\begin{exercise}
A designer proposes to use the incompressible Bernoulli equation
for a small private jet cruising at Mach $0.6$, on the grounds
that the density variation is ``only a few percent.'' Write a
one-paragraph response covering when this is acceptable for
preliminary design and when it is not.
\paragraph{Model-answer sketch.} A strong response checks the number:
at $M = 0.6$ the density variation scales as $M^2/2 \approx 0.18$, so
local density changes can reach $15$ to $20\,\%$ at suction peaks, not
``a few percent.'' For a rough force balance or a first cut at lift,
incompressible Bernoulli with a Prandtl-Glauert correction factor
$1/\sqrt{1 - M^2} \approx 1.25$ is acceptable. It is not acceptable
for the pressure distribution that sets the critical Mach number, for
the airspeed system calibration, or for any case where a local pocket
may approach $M = 1$, because raw incompressible Bernoulli will both
mis-state the suction peak and miss the onset of compressibility
effects entirely. The rule of thumb: use it with the Prandtl-Glauert
correction for preliminary sizing, never for the pressure peaks or the
critical-Mach check.
\end{exercise}

\begin{exercise}
Identify a compressible-flow configuration in your immediate
environment (an aircraft, a gas-vented appliance, an industrial
nozzle) and state, in $200$ to $300$ words, the Mach number you
estimate, the regime it sits in, the dominant compressibility
effect you would expect, and what observable signature confirms
or refutes the estimate.
\paragraph{Model-answer sketch.} A strong answer picks a concrete case
and runs the chapter's logic. A cruising airliner seen from the
ground: $M \approx 0.8$, subsonic-compressible to transonic, dominant
effect is the upper-surface shock and its drag rise; the observable
signature is the swept supercritical wing and the absence of a sonic
boom. A whistling kettle or a hissing gas-appliance jet: the choked
orifice runs at $M = 1$ at the throat with a critical pressure ratio
near $0.53$; the signature is that the sound and flow saturate once
the supply pressure is high enough, independent of further increase. A
compressed-air nozzle or pressure-washer tip: near-sonic to mildly
supersonic, with audible shock noise and possibly faint shock cells.
The answer should state the estimate, the regime, the expected effect,
and a signature that would confirm or refute it, closing the loop
between the number and the observation.
\end{exercise}
